% before: is-indefinite-datum
The lattices of lifted forms live on definite unitary groups, but their metrics are compared through
Hermitian spaces of indefinite signature. The comparison only uses linear algebra over the
coefficient field: lattices, metrics at the infinite places, the matrices of the lifts and the
scalars by which they change metrics. Gluing the images of the lattices of the different tuples of
Hecke characters, place by place, produces the lattice $RR$ whose dual is the auxiliary lattice of
Chapter~\ref{chapter-inert-prime-balance}.

% before: is-numerical-bounds
Each estimate on the indefinite side compares a quantity of the datum with an explicit number. The
numbers depend on the scale and on two constants: $C$ absorbs the implied constants of the error
terms, and $C_H$ fixes the height budget $x^{C_H}b$. The covolume bound is the first such
comparison, and Chapter~\ref{chapter-covolumes} proves it.

% before: is-lower-norm-bound
Comparing the metrics before and after a lift needs a lower bound for the factor by which the lift
changes the Petersson metric. The lower norm bound gives it at the base tuple of characters, lift by
lift, part by part and place by place. It bounds the norms of the lifts of step 9 of the
announcement's timeline from below.

% before: is-surjectivity-bound
At a twisted place the image of the lattice of a lift can differ from the lattice of lifted forms.
The surjectivity bound limits this index by the local losses $F_0$ and an error term, and bounds the
weighted sum of the error terms over the twisted places. Both enter the division inequality through
the covolume of the glued lattice $RR$, which is assembled from these images at the twisted places.

% before: is-functional-bound
The frame slope inequality of Chapter~\ref{chapter-hermitian-slopes} needs an upper bound for
$\mu_{\max}(RR)$, because $\mu_{\min}(CC)=-\mu_{\max}(RR)$ for the dual lattice $CC$. The functional
bound supplies it, measured against the height budget and the mean of the archimedean scalars.

% before: is-twisted-division-bound
The local defects at the twisted places come from Tate-Shafarevich groups of Tate twists, and
Chapter~\ref{chapter-tate-shafarevich} bounds them on average. The surjectivity bound and the
functional bound are proved only for families in which their weighted sum stays within the error
budget $D_0x^{\varepsilon}(B_0+N_c)$.

% before: is-point
The bounds on the lifts concern explicit theta kernels, whose metrics are computed on indefinite
unitary groups. A point is a choice of these indefinite data at one member, with its rows, its
columns and its scale data. The conditions attached to a point say that these data realize the
definite-side datum.

% before: is-complex-data
The rows and columns of a point are read off from Hilbert spaces of automorphic forms through
explicit coordinates. Complex data isolate what this needs: the spaces, the theta lifts between them,
the parts and levels, and the scalars. The indefinite-side data of
Definition~\ref{is-indefinite-datum} are then the matrices of the lifts in these coordinates and the
Gram matrices of the coordinate vectors.

% before: is-realisation
A realization attaches to a definite-side datum the indefinite unitary groups, Weil representations
and theta lifts that produce complex data. Its components are those of the theta tower of
Chapter~\ref{chapter-unitary-theta}, together with the Hecke modules on which the lifts act and the
scalars by which they change metrics.

% before: is-one-map-data
The lower bound for the eigenvalues of the lifts is computed by writing each lift as a single theta
map between spaces of algebraic modular forms. One-map data are such presentations, together with
the weights and vectors that tie them to the realization.

% before: is-bounded-rows
The bounded rows are the first of the four conditions that Chapter~\ref{chapter-assembly} imposes
on the selected realizations. They ask for extremal vectors pairing to the delta polynomial and for
presentations of the lifts as single maps, and they are what the lower norm bound uses from the
realization.

% before: is-primitive-rows
Primitive rows make the realization rational and diagonal: the Hecke operators are rational and
adjoint to one another, the parts are joint eigenspaces, and the lifts act on each part by the
scalars $\lambda$ and $\kappa\prod\lambda$. These are the conditions under which the scalar conditions of
an indefinite-side datum hold and the eigenvalues can be computed.

% before: is-row-side-family
The division inequality is applied to rows and columns that come from one family of definite forms
and are identified with the lattices of the definite-side datum. A row-side family carries these
identifications, and with them the types and isotypic parts that the selected rows of
Chapter~\ref{chapter-assembly} refer to.

% before: is-division-rows
The division rows say that the identifications of a row-side family respect the arithmetic of the
definite-side datum. The glued lattices, the spaces of forms at the used tuples, the theta lifts,
the local indices and the division pairings must agree on the two sides, so that the division
inequality for the rows and columns is the division link for the datum.

% before: is-realisation-characterization
A realization is a list of data; the conditions that tie these data to the definite-side datum and to
each other form its characterization. Later statements use parts of it by name: the core conditions on
the characters, the group, the family, the models and the theta maps, the conventions for the scalars,
and the conditions on the tuples.

% before: is-realisation-join
The near-orthogonality of the columns is proved in a space of automorphic forms, while the division
inequality uses the coordinates of the realization. The join identifies the two: periods, base
vectors, rows, column norms and the pairings with Poincaré series.

% before: is-record-object
The surjectivity bound is proved place by place at the twisted places, through auxiliary lattices
and numerical parameters attached to a point. A record object collects them, and the estimates of
the proof are statements about record objects.

% before: is-canonical-error
The surjectivity bound uses the canonical error terms of a point, which are defined from the indices
themselves. Any error terms that satisfy the three conditions can be replaced by the canonical ones,
at the cost of doubling the constant.

% before: is-hook-lower-bound
The lower norm bound compares $\ell_6(j)$ with an explicit lower bound for the eigenvalues of the
$j$-th lift. That comparison is numerical once the size constant of the family is at least $8$.

% before: is-onto-record-numbers
The estimates at the twisted places compare the lattices of a record object with numbers built from
its sizes. These numbers are explicit expressions in the ranks, the scale and a few further
parameters of the record object, and the same few enter all of the estimates.

% before: is-onto-page-estimates
The first group of estimates involves only the numbers of a record object. They bound the sizes
against $b$ and the terms of the bad primes against the constants $C_{b,1}$ and $C_{b,2}$, and they
fix the sizes of the residue fields at the level places.

% before: is-onto-geometric-estimates
At a twisted place and its conjugate, a record object carries several lattices in the same space.
The geometric estimates say that they are lattices, that they lie in one another and are stable under
the group, and that the index of $M_s(j)$ against a sublattice grows at most linearly with the
dimension, up to a remainder whose weighted sum over the places is small.

% before: is-onto-depth-estimates
The local numbers $e_0(v,j)$ of a record object are bounded by the function $F_0$ of the numerical
bounds, by a different argument for the good primes, for $q$ and $d$, for the prime divisors of $Q_S$
and for the level primes. At a good prime larger than $Y$ both sides vanish.

% before: is-onto-local-estimates
The remaining estimates concern the eigenspace $V_\psi(j)$ and its part $Z_g(j)$ in $V_0(j)$. The
rank and twist bounds compare lattices in $Z_g(j)$ at a twisted place, while the share and bad-place
bounds compare $d_0(j)=\dim V_\psi(j)$ with $n_Z(j)$ and with the rank $D$.

% before: is-onto-recorded-numbers
A record object recorded at a point carries the numbers of the member of the point. Its coordinate
spaces are those of the rational structures of the lifts, and its eigenspace $V_\psi(j)$ is the part
of the $j$-th lift on which the subgroup $P_0$ of the face acts through a fixed character.

% before: is-onto-page-numerics
For a recorded object the first four page estimates are inequalities among the ranks, the size data,
the host and the parameter choices of the member. They hold at the large scales of every standard
family, given the conditions on the constants $c$ and $3C\le x$.

% before: is-onto-twist-transport
The twist condition of a family bounds an index of theta lattices over a discrete valuation ring
finite over the coefficients of the twisting data, while the twist bound of a record object is an
index at a place of the coefficient field. A semilinear map that carries a basis to an independent
family transports the first bound to the second, at the cost of the ramification index of the
place.

% before: is-onto-share-flag-data
The share bound compares the dimension of a lift with that of an eigenspace of a finite group, through
a model of the lift by cusp forms on a flag domain. These cusp forms are holomorphic functions of
tuples of matrices, one for each embedding of a CM type, that transform under an upper triangular
group and an arithmetic group. Their sizes are measured by the leading minors of Gram matrices, and
the majorant at a point of the domain is a positive form that measures arithmetic elements there.

% before: is-onto-share-two-level
A model of a lift uses several copies of these data, indexed by a finite set on which the group of
the realization acts. The group moves the copies into one another by unitary elements, and so acts on
the cusp forms. The level hypothesis is a congruence condition at a prime that splits completely, and
the archimedean condition is the input for counting cusp forms when the signature is indefinite.

% before: is-onto-share-majorant
The arguments below use a few properties of the majorant. It is nonnegative, unchanged by the upper
triangular group, and equal to the standard norm in the coordinates of an orthonormal point; when
the form is definite, it is also unchanged by the unitary group.

% before: is-onto-share-free-summands
When a group permutes the summands of a direct sum freely, every eigenspace of a character of a
subgroup has its expected share of the dimension. The trace of a nontrivial element vanishes, since
it moves every summand, so the averaging projection onto the eigenspace has trace the dimension
divided by the order.

% before: is-onto-share-components
The cusp forms of two-level data are functions on a disjoint union of copies of the flag domain. They
split into the forms supported on single copies, and a pulled-back representation permutes these
pieces as the group permutes the copies.

% before: is-onto-share-borel-weil
In the definite case the flag domain consists of all tuples of invertible matrices, and the finiteness
of the space of cusp forms reduces to a statement about one matrix variable. A holomorphic function
on $GL_r(\mathbf C)$ that transforms under the upper triangular group by a character with nonnegative
exponents is a polynomial.

% before: is-onto-share-definite-finite
By the Borel-Weil theorem every slice of a section in one matrix variable is a polynomial of bounded
degree divided by a power of the determinant. A function all of whose slices lie in
finite-dimensional spaces is determined by its values on a finite grid, so the sections span a
finite-dimensional space.

% before: is-onto-share-unipotent
The freeness of the action in the definite case rests on a property of algebraic matrices that are
close to isometries at every embedding. If the traces of their powers are congruent to $r$ modulo a
large integer $g$, the product formula makes these traces equal to $r$, and the matrix is unipotent
and close to $1$.

% before: is-onto-share-trace
The trace congruence comes from the level hypothesis. An element that is unipotent modulo every prime
over a completely split prime $g$ has trace congruent to $r$ modulo each of these primes, hence
modulo $g$.

% before: is-onto-share-small
At a point of the flag domain where no nonzero lattice endomorphism contracts the majorant at every
embedding, an element of a group unipotent modulo the primes over $g$ that is nearly an isometry is
the identity. Its difference
with the identity is a lattice endomorphism, and the unipotence criterion bounds its Hilbert-Schmidt
norm below $1$.

% before: is-onto-share-degree
Semistability is measured by the degree of a lattice, the logarithm of an index minus half the sum
of the logarithms of Gram determinants for a family of majorants. The degree is additive: that of a
lattice is the degree of its intersection with a subspace plus a quotient degree, computed from the
Gram matrices of the complementary vectors after projecting away from the subspace.

% before: is-onto-semistable-basis
The degree is defined with a chosen basis of the lattice. A change of basis multiplies the index and
the Gram determinants by norms of one determinant, which cancel by the product formula, so the degree
depends only on the lattice.

% before: is-onto-semistable-det
Degrees are compared through Gram determinants. Two Hermitian forms, one positive definite and
dominating the other, have determinants in the same order, by diagonalizing one form relative to the
other.

% before: is-onto-semistable-volume
The index of a lattice over the span of a basis is controlled by the Gram determinant of the basis for
the Hermitian form. For an integral lattice the square of the index is at most the norm of that
determinant, and for a self-dual lattice it is at least that norm, because the volume ideal of a
self-dual lattice is the unit ideal.

% before: is-onto-semistable-witt
The majorant is built in coordinates adapted to the form. In a Witt basis the form is a sum of
hyperbolic planes and an anisotropic diagonal part, and a diagonal majorant can be chosen separately
on each piece.

% before: is-onto-semistable-graded
For a diagonal majorant, degrees can be bounded along the flag of coordinate subspaces. The degree of
a lattice is at most the degree of its part in a coordinate subspace plus the degree of its
projection to the remaining coordinates, so a lattice whose successive projections have only
submodules of nonpositive degree has the same property.

% before: is-onto-semistable-line
The successive projections of the lattice to single coordinates are rank one modules. For them the
degree is an explicit sum, and nonpositive degree passes to submodules.

% before: is-onto-semistable-pair
The two coordinates of a hyperbolic plane receive reciprocal weights in the majorant. The pairing of
the form between them is integral on the lattice, so the integral index bound for the plane makes the
degree of the second coordinate nonpositive once that of the first is zero.

% before: is-onto-semistable-middle
On the anisotropic part the majorant is the absolute value of the diagonal form at each embedding. The
form then dominates its own Gram matrices, and the integral index bound gives nonpositive degree for
every submodule.

% before: is-onto-semistable-majorant
The diagonal weights in Witt coordinates define a majorant of the form. Its degrees are those of the
transported lattice for the diagonal weights, so the bounds proved in Witt coordinates carry over.

% before: is-onto-share-semistable
Points without contraction come from a majorant of the Hermitian form for which the self-dual lattice
is semistable. For such a majorant no nonzero lattice endomorphism shrinks every vector at every
embedding.

% before: is-onto-share-balanced
A majorant of this kind gives a point of the flag domain: at each embedding of the CM type, a basis
orthonormal for the form and for the majorant. At this point the majorant of the flag data is the
given one, so there is no contraction. The semistability of the lattice also bounds the majorants of
its nonzero vectors from below.

% before: is-onto-share-free
In the definite case the unitary group at every embedding is compact, and the transition elements of
two-level data are isometries for the majorant. An element of the group that fixes a copy has a
transition element that is unipotent modulo the primes over $g$ and nearly an isometry, and the
previous lemmas make it the identity.

% before: is-onto-share-definite
The first dimension statement follows. Each copy has a finite-dimensional space of cusp forms by the
Borel-Weil theorem, and every subgroup for which the level hypothesis holds permutes the copies
freely, so the eigenspaces of its characters have their expected shares.

% before: is-onto-share-pullback
In the indefinite case the copies need not be permuted freely, and the eigenspaces are bounded below
by multiplying cusp forms by translates of one auxiliary form. Translating a cusp form of one copy by a
transition element gives a bounded section of another copy, since the transition elements are
unitary.

% before: is-onto-sections-fourier
Cusp forms are defined by the vanishing of constant terms, and products are controlled through their
Fourier-Jacobi coefficients along an isotropic vector. For a bounded section only coefficients of
nonnegative index occur, and the coefficients of a product are the convolution of those of its
factors.

% before: is-onto-share-sections
Products of cusp forms with bounded sections are cusp forms, and the flag domain is connected, so a
holomorphic function on it that vanishes near one point vanishes everywhere. With these two facts,
multiplication by translates carries cusp forms of a smaller weight injectively into an eigenspace.

% before: is-onto-step-free-ball-majorant
In the indefinite case the auxiliary form is built near an orthonormal point that no nontrivial element of
the group moves close to itself. The comparison of nearby orthonormal frames rests on linear algebra for a
form of signature $(p,r-p)$: a majorant $H$ of the form whose positive eigenspace for $JH$ is close to the
positive coordinate space is close to the identity, with explicit constants.

% before: is-onto-step-free-ball-minor
Two orthonormal points of the flag domain differ by a matrix $g$ that preserves the form $J$. The leading
$p_0\times p_0$ minor of $g$ is at least $1$ in absolute value, and it grows with the part of $g$ that mixes
the positive and the negative coordinates.

% before: is-onto-step-free-ball
The leading minors measure how far a point of the flag domain is from an orthonormal point $x_0$, up to
the action of $\mathcal B^\Psi$. A group has a free ball at $x_0$ if its nontrivial elements move every point
near $x_0$ away from $x_0$.

% before: is-onto-step-free-ball-bound
If an orthonormal point is near $x_0$ in this sense, the matrix that relates the two frames, and its
inverse, change lengths by a factor close to $1$. So an element of the group that maps a point near $x_0$
to a point near $x_0$ is nearly an isometry for the majorant at $x_0$.

% before: is-onto-step-free-ball-exists
At a point without contraction, an element of the group that is unipotent modulo the primes over $g$ and
nearly an isometry for the majorant is the identity. With the previous bound, the group has a free ball of
radius $1/(49r^2)$ at the orthonormal point below a given point without contraction.

% before: is-onto-step-orbit-situation
The interpolation at the centre of a free ball is carried out for any group with a free ball at an
orthonormal point, in the indefinite case. The weight is $\delta_s$, with $s$ large compared with
$r^2|\Psi|p_0(r-p_0)$, which provides the curvature the construction needs.

% before: is-onto-step-orbit-interpolation-large-span
The interpolating cusp form is obtained from approximate ones. Only the values at finitely many points
matter, and these values range over a finite-dimensional space, so approximate values $1$ and $0$ give
exact ones.

% before: is-onto-step-orbit-interpolation-large-weight-near
The sums over $\Gamma$ are controlled by a weight that measures, through leading $p_0\times p_0$ minors
relative to $x_0$, how far a point is from $x_0$. Moving a point within a small ball of its chart changes
this weight by a bounded factor, uniformly over the orthonormal points.

% before: is-onto-step-orbit-interpolation-large-weight-power
A power of this weight serves as the cutoff of the interpolation. Its invariance under $\mathcal B^\Psi$,
its values in the chart at $x_0$, the local comparison and the bound for its sums over $\Gamma$ all come from
the corresponding properties of the weight.

% before: is-onto-step-orbit-interpolation-large-weight
For $\nu>4|\Psi|r^2$ the sums over $\Gamma$ of this power are bounded by the counting argument of the
Poincaré series. This gives the cutoff when $p_0(r-p_0)\geq4$.

% before: is-onto-step-orbit-interpolation-large-pole
The $\bar\partial$ equation is solved with a weight that has a logarithmic pole along the orbit of $x_0$.
The pole function is invariant under $G$ and $\mathcal B^\Psi$, at most $\dim X\log(\rho(x_0,x)+\delta)$ up
to a constant near $x_0$, bounded below on a shell around $x_0$, and its Levi form is bounded below by a
multiple of that of $\rho$.

% before: is-onto-step-orbit-interpolation-large-pole-exists-cutoff
The pole function is a function of one variable applied to $\rho(x_0,\cdot)$. That function must behave like
$\log(u+\delta)$ near $0$, vanish before $\rho_0$, and satisfy a lower bound for $g'+ug''$ with the factor
$\theta/\rho_0$, which becomes the lower bound for the Levi form.

% before: is-onto-step-orbit-interpolation-large-pole-exists-block
Moving a point within a small ball of its chart changes the leading $p_0\times p_0$ minor relative to an
orthonormal point by a bounded factor. The factor can be taken as close to $1$ from below as needed, and at most
$2$ from above.

% before: is-onto-step-orbit-interpolation-large-pole-exists-translate
The pole function is defined through the translates of the free ball under $G$. Translating both arguments of
$\rho$ by an element that preserves $h$ does not change it, and the free ball determines the translating element.

% before: is-onto-step-orbit-interpolation-large-pole-exists-lower
The pole function vanishes away from the translates of the free ball. To see that it vanishes on a
neighbourhood of such a point, $\rho(x_2,\cdot)$ is bounded below near the point uniformly in the orthonormal
point $x_2$.

% before: is-onto-step-orbit-interpolation-large-pole-exists-sharp
The Levi form of a composite $g\circ f$ involves the squared gradient of $f$. It is controlled by the Levi form
of $f$ when the squared gradient is at most $f$ times the Levi form.

% before: is-onto-step-orbit-interpolation-large-pole-exists-sharp-rules
The bound holds for $-\log(1-\sum_i|g_i|^2)$ with holomorphic $g_i$, passes to finite sums, and passes from a
function to a larger one that agrees with it at the point. These rules give it for $-\log\det(1-WW^*)$ with $W$
holomorphic.

% before: is-onto-step-orbit-interpolation-large-pole-exists-transfer
In the chart at $x_1$, $\rho(x_2,\cdot)$ differs from $\rho_{x_1}$ by the logarithm of the squared absolute value
of a holomorphic function, so both have the same Levi form. It is also $-\log\det(1-WW^*)$ for a holomorphic
$W$, which gives the bound for its gradient that the chain rule needs.

% before: is-onto-step-orbit-interpolation-large-pole-exists-function
The pole function at $x$ is a function of $\rho(x_0,\gamma x)$ for the element $\gamma$ of $G$ that brings $x$
into the free ball, and $0$ if there is none. This makes it invariant under $G$ and $\mathcal B^\Psi$.

% before: is-onto-step-orbit-interpolation-large-pole-exists-function-local
Near each point of a chart the pole function is either a function of $\rho(x_2,\cdot)$ for one orthonormal
point $x_2$, or $0$. Its smoothness and its Levi form are therefore read off from those of $\rho(x_2,\cdot)$.

% before: is-onto-step-orbit-interpolation-large-pole-exists
Pole functions exist at the centre of a free ball, with a factor $\theta>1$ in the lower bound for the Levi
form. The free ball keeps the poles at the different points of the orbit of $x_0$ apart.

% before: is-onto-step-orbit-interpolation-large-cutoff
The datum of the $\bar\partial$ problem is a cutoff in $\rho_{x_0}$, equal to $1$ near $x_0$ and to $0$
away from it. Its derivative is bounded and supported in the shell where the pole function is bounded
below.

% before: is-onto-step-orbit-interpolation-large-block
For a weight that is constant on the two blocks of indices, the potential $\varphi_\lambda$ in the chart at
$x_0$ is a combination of the logarithms of the minors $d_{p_0}$. It is therefore bounded below where
$\rho_{x_0}$ is bounded.

% before: is-onto-step-orbit-interpolation-large-sublevel
The sets where $\rho_{x_0}$ is bounded have finite volume for the measure $D_{x_0}(z)\,dz$. Up to
$\mathcal B^\Psi$, their points are orthonormal points that differ from $x_0$ by bounded matrices, and the
saturations of compact sets have finite volume by the invariance of the measure.

% before: is-onto-step-orbit-interpolation-large-gradient
Hörmander's estimate needs the derivative of the cutoff to be controlled by the curvature of the weight.
On the shell this follows from a bound for the gradient of $\rho_{x_0}$ by its Levi form.

% before: is-onto-step-orbit-interpolation-large-solve-levi
The weight of the $\bar\partial$ problem has a Levi form bounded below by a positive multiple of that of
$\rho_{x_0}$. The gradient bound shows that this multiple is nonnegative.

% before: is-onto-step-orbit-interpolation-large-solve
With these inputs, Hörmander's estimate solves $\bar\partial$ for the cutoff, with a weight that contains
the pole function. The bound for the solution does not depend on $\delta$, so $\delta$ can later be taken
small.

% before: is-onto-step-orbit-interpolation-large-near-zero
Near $0$ in the chart, $\rho_{x_0}$ is bounded by a multiple of the squared norm of the coordinates that
mix the two blocks. There are $\dim X$ such coordinates, the exponent of the pole.

% before: is-onto-step-orbit-interpolation-large-blowup-product
The weight $\max(|x|^2,t)^{-d}$ in $d$ complex variables has total mass tending to infinity as $t\to0$. By
the mean value property, its integral against $|u|^2$ then dominates any multiple of $|u(0)|^2$.

% before: is-onto-step-orbit-interpolation-large-blowup
The same holds with $|x|^2$ replaced by $|p|^2$, for a holomorphic map $p$ whose derivative at the point
is surjective. A holomorphic change of coordinates reduces it to the product case.

% before: is-onto-step-orbit-interpolation-large-point
In the chart at an orthonormal point, the weight $(\rho+\delta)^{-\dim X}$ concentrates at $0$ as
$\delta\to0$. So the value of a holomorphic function at $0$ is bounded by a small multiple of its weighted
integral, uniformly over the orthonormal points.

% before: is-onto-step-orbit-interpolation-large-ball
The pole and the cutoff weight are compared on a small ball of the chart at each orthonormal point. On it
$\rho$ is small, $N_\lambda$ is bounded below, and the cutoff weight changes by a bounded factor.

% before: is-onto-step-orbit-interpolation-large-approx
For every $\eta>0$ these estimates give a $\mathcal B$-equivariant function within $\eta$ of $1$ at $x_0$
and small, relative to the cutoff weight, at the points $gx_0$ with $g\ne1$. Its restriction to the chart
at $x_0$ is square integrable for the weight of the Poincaré series.

% before: is-onto-step-orbit-interpolation-large-interpolation
The Poincaré series of these functions are cusp forms whose values at $x_0$ and at representatives of the
points $gx_0$ modulo $\Gamma$ approach $1$ and $0$. As $\Gamma$ has finite index, there are finitely many
representatives, and the approximate values give an exact interpolation.

% before: is-onto-step-orbit-interpolation-large-curvature
For the weight $\delta_s$ the curvature condition of the interpolation holds with $m=s-2r+2$. It comes
from the differences between consecutive entries of $\delta_s-k$, which are those of $-k$ except across
$p_0$.

% before: is-onto-step-orbit-interpolation-large
When $p_0(r-p_0)\geq4$, Alpöge's $L^2$ method gives the interpolating cusp form. Holomorphic sections near
$x_0$ come from estimates for $\bar\partial$ with a weight that has a logarithmic pole at $x_0$, and their
Poincaré series are cusp forms.

% before: is-onto-step-orbit-interpolation-small-integral-coordinates
The integral of $e^{-s\rho_{x_0}}$ is computed in coordinates in which $\rho_{x_0}$ depends only on the mixed
blocks. Multiplying each mixed block by the diagonal block above it is such a change of coordinates, and it
preserves Lebesgue measure.

% before: is-onto-step-orbit-interpolation-small-integral-minors
In these coordinates the density is a product of Gram minors of matrices with a mixed column or row. Two
determinant inequalities bound such minors below by $1-|y|^2$ times Gram minors of the remaining columns.

% before: is-onto-step-orbit-interpolation-small-integral-density
When $p_0=1$ or $p_0=r-1$ the mixed block is a single column or row. For $s\geq2(r-1)$ the density times
$e^{-s\rho}$ is then bounded by Gram minors of a unitriangular block, which do not involve the mixed entries.

% before: is-onto-step-orbit-interpolation-small-integral-blocks
As the mixed entries are at most $1$ in absolute value, the bound leaves an integral over the coordinates at a
single $\tau$. It is finite for $r\le4$, the ranks of the signatures with $p_0(r-p_0)\le3$.

% before: is-onto-step-orbit-interpolation-small-integral
When $p_0(r-p_0)\le3$ the bound on $s$ does not give $\nu>4|\Psi|r^2$, which the counting argument for the
sums over $\Gamma$ needs. It is replaced by the finiteness of the integral of $e^{-s\rho}$ over the chart,
which holds for $s\geq2(r-1)$ in these signatures.

% before: is-onto-step-orbit-interpolation-small-orbit-sum-fundamental
Sums over $\Gamma$ are compared with integrals by unfolding over a fundamental set in a chart. Such a set meets
the orbit of almost every point under $\Gamma$ and $\mathcal B^\Psi$ in one point.

% before: is-onto-step-orbit-interpolation-small-orbit-sum-fundamental-exists
By proper discontinuity, every point of the chart has a neighbourhood that meets each orbit at most once.
Countably many of these neighbourhoods, with the saturations of the earlier ones removed, form a fundamental set.

% before: is-onto-step-orbit-interpolation-small-orbit-sum-unfold
An integral over the chart unfolds to the integral over a fundamental set of the sum over $\Gamma$. An integral
over a subset of another chart is then bounded by the integral over the fundamental set, times the number of
elements of $\Gamma$ that move the subset into itself modulo $\mathcal B^\Psi$.

% before: is-onto-step-orbit-interpolation-small-orbit-sum-harnack
The averaging needs $e^{-s\rho(x,\cdot)}$ to change by a bounded factor on a small ball of each chart. This
follows from the bound for the leading minor under small perturbations, and the symmetry of $\rho$ on orthonormal
points exchanges the roles of $x_0$ and $x$.

% before: is-onto-step-orbit-interpolation-small-orbit-sum
A bound for that integral, uniform over the orthonormal points, bounds the sums over $\Gamma$ of
$e^{-s\rho(x_0,\gamma x)}$. The sum is averaged over a small ball and unfolded to the integral.

% before: is-onto-step-orbit-interpolation-small-weight
With $s=\nu/2$ the integral bound and the unfolding give the sum bound in the small signatures for
$\nu\geq4(r-1)$. The cutoff weight is then the same power of the minor ratio as when $p_0(r-p_0)\geq4$.

% before: is-onto-step-orbit-interpolation-small-poincare
The Poincaré series converge and are bounded as before, once the sums over $\Gamma$ of
$e^{-(\nu/2)\rho(x_0,\gamma x)}$ are bounded. That bound is the only step of the earlier argument that used
$\nu>4|\Psi|r^2$.

% before: is-onto-step-orbit-interpolation-small-interpolation
In the small signatures the interpolation runs as for the large ones. The cutoff weight is the power of the
minor ratio whose sums over $\Gamma$ are bounded through the integral over the chart, and the Poincaré series
converge by the corresponding bound.

% before: is-onto-step-orbit-interpolation-small
The small signatures $p_0(r-p_0)\le3$ are treated by the same method. The sums over $\Gamma$ are then
controlled through a uniform bound for integrals over balls of the flag domain.

% before: is-onto-step-orbit-interpolation
Since $p_0(r-p_0)$ is a positive integer when $0<p_0<r$, the two cases cover every indefinite signature.
Together they give the interpolating cusp form at the centre of any free ball of radius $1/(49r^2)$.

% before: is-onto-share-interpolation
The auxiliary form is assembled from cusp forms on single copies. At a point without contraction, a
cusp form of weight $\delta_s$ can be chosen nonzero at the point and zero at its translates by the
elements of the stabilizer of the copy outside its group $\Gamma_c$. This uses the level hypothesis and a
lower bound for $s$ in terms of the dimension of the flag domain.

% before: is-onto-share-translates
The auxiliary form has weight $\delta_s$, and its translates by the group must be independent at some
point of every copy. Summing the forms of the previous lemma over representatives of the orbits makes
the matrix of translates at a suitable point a multiple of a permutation matrix.

% before: is-onto-share-indefinite
Multiplying a cusp form of weight $\kappa-\delta_s$ by the translates of the auxiliary form and
projecting to an eigenspace gives $|H|$ elements of the eigenspace. If all of them vanish, the form
vanishes where the determinant of the translates is nonzero, hence everywhere. So the eigenspace has
at least $1/|H|$ of the dimension of the space of weight $\kappa-\delta_s$, which the ranks compare
with the full space.

% before: is-onto-share-counts-model
The counts of cusp forms on a copy use a few parts of its archimedean model. The flat model gives local
coordinates at finitely many base points, and further sections and an invariant function control norms
near the base points and towards the cusps. A quotient of norms built from these parts measures how fast a
section vanishes at the base points.

% before: is-onto-share-counts-jets
A combination of many holomorphic functions can be made to vanish to high order at every base point. Near
a base point the flat model gives coordinates, and the Taylor coefficients up to a given order at all base
points form a space whose dimension is the number of base points times a binomial coefficient.

% before: is-onto-share-counts-reduction-lll
The quotient of norms has to attain its maximum, and for this the points at which the first minimum of the
lattice is bounded below have to lie, up to the group and the triangular matrices, in a compact set. The
argument reduces a basis of the lattice for a positive definite quadratic form. A basis reduced in the
sense of Lenstra, Lenstra and Lovász has vectors of bounded length and controls the coordinates of every
lattice vector.

% before: is-onto-share-counts-reduction-lattice
The quadratic forms come from the lattice $\Lambda$ placed in the complex space $(\mathbf C^r)^\Psi$ through the
embeddings in $\Psi$. Since $\Psi$ contains one embedding from each pair of complex conjugate embeddings, a
basis of $\Lambda$ over $\mathbf Z$ becomes a real basis of this space.

% before: is-onto-share-counts-reduction-normal
A point of the flag domain can be moved by a triangular matrix to an orthonormal point at which the
majorant becomes the standard norm. The triangular matrix comes from orthonormalizing the columns for the
Hermitian form of the majorant, which is positive definite.

% before: is-onto-share-counts-reduction-forms
At an orthonormal point the majorants define a positive definite quadratic form on $\Lambda$. Its
determinant is bounded by that of a fixed form built from $h$, and it bounds the values of $h$ at every
complex embedding, the effect of multiplying by an element of $E$, and the coordinates of rational
combinations.

% before: is-onto-share-counts-reduction-basis
When the first minimum is bounded below, the reduction of bases applies to the quadratic form of an
orthonormal point. It gives a basis of $\Lambda$ whose vectors have bounded values of the form, and in which
the coordinates of every vector are controlled by its value.

% before: is-onto-share-counts-reduction-types
A basis of $\Lambda$ carries two kinds of arithmetic data: its Gram matrix for $h$ and the integer matrices
of multiplication by a basis of $\mathcal O_E$. Bases with the same data are carried to each other by a
unitary transformation that preserves $\Lambda$.

% before: is-onto-share-counts-reduction-compact
The compact set is a set of orthonormal points with bounded entries. For a matrix with Gram matrix $J$, the
entries are bounded by those of its inverse, which are coordinates of vectors of $E^r$ and so can be bounded
through the quadratic form.

% before: is-onto-share-counts-reduction
The reduced basis of a point with large first minimum has a bounded Gram matrix and bounded multiplication
matrices, so it has one of finitely many types. The unitary transformation that carries it to a fixed basis
of its type, corrected by one of finitely many cosets of $\Gamma$, moves the orthonormal point into a compact
set.

% before: is-onto-share-counts-maximum
For a section that vanishes to high order at the base points, the quotient of norms is small near the base
points, and the invariant function of the model makes it small where the first minimum is small. On the
remaining compact set it attains a maximum, at a point where the flat coordinates do not all vanish.

% before: is-onto-share-counts-no-maximum
In the chart at such a maximum, the logarithm of the quotient is a sum of logarithms of holomorphic
functions and half the potential of the weight $L$. In a direction along which the ratios of the flat
coordinates are constant, only the potential contributes to the Levi form, and its positivity excludes a
maximum.

% before: is-onto-share-counts-upper
The last three lemmas show that no nonzero bounded section of the weights in question vanishes to high
order at all base points. Counting the Taylor coefficients then bounds the number of independent bounded
sections.

% before: is-onto-share-counts-section-levi
The remaining counts need bounded sections of the weights $(n+4)L-2\delta_s$ and $F-\delta_s$ with prescribed
behaviour at the base points. Both are built by one method, which needs lower bounds for Levi forms in the
charts at orthonormal points. The bounds compare the Levi form of the potential $\varphi_\lambda$ of the weight with
those of a function $\rho_{x_0}$ built from the minors of order $p_0$, of the logarithm of a density $D_{x_0}$, and
of the potential of $L$.

% before: is-onto-share-counts-section-diagonal
At the centre of the chart of an orthonormal point the Levi forms of the potentials can be computed. The terms of
degree two of the leading minors of the Gram matrix give a diagonal form, whose coefficients are differences of
the entries of the weight with signs from the signature.

% before: is-onto-share-counts-section-transfer-lu
To compare charts at different points, a point of the flag domain is written in the chart of a fixed orthonormal
point, up to a triangular factor. Gaussian elimination does this wherever the leading minors of the point, read in
the basis of the orthonormal point, do not vanish, and the coordinate and the factor depend holomorphically on the
point.

% before: is-onto-share-counts-section-transfer-jacobian
When one chart is written in another up to a triangular factor, the derivative of the change of coordinates is
determined by that factor. Its determinant is a character of the triangular factor, of the weight $k$ that also
gives the density $D_{x_0}$; in particular it does not vanish.

% before: is-onto-share-counts-section-transfer
A point of a chart is the centre of the chart of an orthonormal point, up to a holomorphic change of coordinates
and a triangular factor. The triangular factor changes the potential by a pluriharmonic function, so Levi forms
anywhere in a chart are Levi forms at the centre of another chart.

% before: is-onto-share-counts-section-curvature
The diagonal form shows when the Levi form of the potential of a weight dominates a multiple of that of
$\rho_{x_0}$. The entries of the weight must not increase within the first $p_0$ indices and within the remaining
ones, and they must increase by at least the multiple across the index $p_0$.

% before: is-onto-share-counts-section-density
The weight $k$ with entries $2i-r-1$ turns the potential into the logarithm of the density $D_{x_0}$, as the last
minor has absolute value one in the chart of an orthonormal point. So shifting a weight by $k$ and by multiples
of $L$ changes Levi forms linearly, and curvature constants reduce to the previous lemma.

% before: is-onto-share-counts-section-gauge
The sections are first built in a chart, as holomorphic functions that are square integrable against a gauge of
the weight, and then summed over $\Gamma$. A gauge transforms under the triangular matrices like the square of the
norm of a section and is invariant under $\Gamma$.

% before: is-onto-share-counts-section-kronecker-points
The construction at a base point needs local coordinates there, and it must recognize the other translates of the
base points. At every point $\gamma y_jb$ the flat coordinates $s_k/w_0$ vanish and have invertible derivative in the
chart, as a change of charts with a triangular factor has invertible derivative.

% before: is-onto-share-counts-section-kronecker-gauge
The simplest gauge of a weight is the square of the norm factor $N_\lambda$ on the flag domain. It transforms as a
gauge must because the Gram matrices change only through the diagonal of a triangular factor, and not at all under
$\Gamma$.

% before: is-onto-share-counts-section-kronecker-solve-simple-zero
A function constructed with a singular weight is controlled at a common zero of the flat coordinates through the
integrability of the weight. Near a simple common zero of $n$ functions in $n$ variables the sum of their squares
is comparable to $|z|^2$, and $|z|^{-2n}$ is not integrable in real dimension $2n$; so a holomorphic function with
such an integral must vanish there.

% before: is-onto-share-counts-section-kronecker-solve-chart-matrices
In the chart of an orthonormal point the function $\rho_{x_0}$ will be a sum of logarithms of determinants of
matrices $1-ZZ^*$. Two inequalities for these determinants are needed: each lies between $0$ and $1$, and the
determinant of $1-ZW^*$ bounds the product of two of them. The second makes $\rho_{x_0}$, corrected by a pluriharmonic
function, minimal at any given point of the chart.

% before: is-onto-share-counts-section-kronecker-solve-chart-domain
The flag domain is defined through the signs of all the leading minors of the Gram matrices. When the columns are
independent, it suffices that the Gram matrix of the first $p_0$ columns is positive definite: vectors orthogonal to
these columns then have negative norm, which gives the signs of the remaining minors. This reduces the chart to one
condition on a $p_0\times p_0$ matrix for each embedding.

% before: is-onto-share-counts-section-kronecker-solve-chart-block
In the chart of an orthonormal point, the Gram matrix of the first $p_0$ columns is congruent, through a unitriangular
matrix, to a matrix $1-ZZ^*$ whose $Z$ depends holomorphically on the coordinates. So the chart is the set where these
matrices are positive definite, and $\rho_{x_0}$ is minus the sum of the logarithms of their determinants.

% before: is-onto-share-counts-section-kronecker-solve-chart
The $\bar\partial$-equation will be solved on the chart $\Omega(x_0)$ itself, which needs it to be pseudoconvex. The
function $\rho_{x_0}$, corrected by $\log(1+|z|^2)$, serves as a plurisubharmonic exhaustion of the chart, and the
chart is connected, so that holomorphic functions on it obey the identity theorem.

% before: is-onto-share-counts-section-kronecker-solve-hormander
The holomorphic function is obtained from a smooth one by an $L^2$ solution of the $\bar\partial$-equation. The
estimate needed is the one for a weight whose Levi form is only semi-positive, with the datum controlled pointwise
by the Levi form, on an open set with a plurisubharmonic exhaustion.

% before: is-onto-share-counts-section-kronecker-solve-riemann
The $L^2$ solution is produced off the zero set of one flat coordinate, where the singular weight is smooth. A
holomorphic function that is square integrable near a hypersurface extends across it, which brings the solution
back to the whole chart.

% before: is-onto-share-counts-section-kronecker-solve-interpolation
These ingredients combine into an interpolation statement on an open set of $\mathbf C^n$. A cutoff equal to $1$
near an isolated common zero of $n$ holomorphic functions is corrected to a holomorphic function, with an integral
weighted by a negative power of the sum of their squares. The weight's strict positivity is needed only where the
cutoff is not locally constant.

% before: is-onto-share-counts-section-kronecker-solve
In the chart of an orthonormal point that is a translate of a base point, a cutoff function equal to $1$ near the
centre is corrected to a holomorphic function by solving the $\bar\partial$-equation with $L^2$ estimates. The weight
includes a negative power of the flat coordinates, which forces the correction to vanish at the centre.

% before: is-onto-share-counts-section-kronecker-extend-measure
Square integrability on the chart has to give square integrability on compact sets of the flag domain. For
functions invariant under the triangular matrices, an integral over a compact set is bounded by a multiple of the
integral against $D_{x_0}(z)\,dz$ over the points of the chart that reach the set modulo triangular matrices.

% before: is-onto-share-counts-section-kronecker-extend
The holomorphic function of the chart has to become an equivariant function on the whole flag domain. The
triangular factorization defines it on a dense open set, and the integrability bounds let it extend.

% before: is-onto-share-counts-section-kronecker-vanish
The same singular weight shows that the extended function vanishes at the other translates of the base points. The
lemma needs only that such a translate is not reached from the small ball of the chart by a triangular factor.

% before: is-onto-share-counts-section-kronecker
With curvature constants, a function can be built at one base point that is nonzero there and vanishes at all
other translates of base points under $\Gamma$. In the chart of an orthonormal point it is square integrable
against a gauge that is bounded below, which is what its Poincaré series needs.

% before: is-onto-share-counts-section-discontinuous
The sums over $\Gamma$ in the Poincaré series are controlled through the way $\Gamma$ acts modulo the triangular
matrices. Near two given points of the flag domain, only finitely many elements of $\Gamma$ carry one point into
the other up to a triangular matrix.

% before: is-onto-share-counts-section-poincare-measure
Integrals over the charts at different orthonormal points have to be compared, and sums over $\Gamma$ of integrals
over a fundamental set have to be unfolded. Both work for functions invariant under the triangular matrices,
against the density $D_{x_0}$.

% before: is-onto-share-counts-section-poincare-mean
A value of a holomorphic function is bounded by an integral of its square over a small ball of a chart. At
orthonormal points the ball and the constant can be chosen once for all points, and near a fixed point the same
holds for the equivariant functions, through the triangular factorization.

% before: is-onto-share-counts-section-poincare-count-separation
The sums over $\Gamma$ are compared with a count of matrices that are uniformly separated. Proper discontinuity
gives this separation at an orthonormal point: the elements of $\Gamma$ other than $1$ move it, read in its own
frame, by a fixed positive amount.

% before: is-onto-share-counts-section-poincare-count-sum
Matrices that preserve a form of signature $(p,q)$ have entries bounded by the determinant of a diagonal block.
Separated families of such matrices therefore have few members with a small block determinant, and a dyadic count
sums negative powers of these determinants.

% before: is-onto-share-counts-section-poincare-count
The sums over $\Gamma$ are controlled by a function $\rho_Y$ on the flag domain that is invariant under the
triangular matrices and restricts to $\rho_{x_0}$ on the chart of $x_0$. It changes boundedly on small balls of
the charts, and its exponentials are summable over $\Gamma$ once the exponent exceeds $2|\Psi|r^2$.

% before: is-onto-share-counts-section-poincare-ball
Square integrability in the chart of $x_0$, read in the chart of a translate $\gamma x$, gives a bound for the
integral over a small ball that decays like $e^{-\nu\rho_Y(\gamma x)}$. Summing the square roots of these bounds over
$\Gamma$ converges, as $\nu/2$ exceeds $2|\Psi|r^2$.

% before: is-onto-share-counts-section-poincare
Summing a square integrable function of a chart over $\Gamma$ gives a bounded section. The sum converges
absolutely, it is square integrable on fundamental sets, it is bounded pointwise through the gauge, and it
commutes with multiplication by sections.

% before: is-onto-share-counts-section-separate
The Poincaré series of the function built at one base point keeps a single nonzero term at that base point and
vanishes at the others. So curvature constants for a weight give bounded sections of that weight that separate
the base points.

% before: is-onto-share-counts-section-cusp-minors
Constant terms are integrals along the unipotent flows $\phi_\xi$ of an isotropic vector. To follow a point under
these flows with complex parameters, one needs the leading minors of the moved point. They change by factors that
are affine in the imaginary part of the parameter, with coefficients given by projections of the isotropic vector
onto the spans of the first columns.

% before: is-onto-share-counts-section-cusp-flow
The factors of the previous lemma decide when a complex flow stays in the flag domain, and how the norm factor
$N_\lambda$ changes along it. For a weight whose entries do not increase within the two blocks and increase from the
first entry to the last at one embedding, an imaginary flow at that embedding makes $N_\lambda$ grow like a power
of the flow parameter.

% before: is-onto-share-counts-section-cusp-constant
The constant term of a bounded section is bounded by the same constant against $N_\lambda$, and it is holomorphic.
It is invariant under the real flows by the periodicity of the section under the transvections of $\Gamma$, and the
identity theorem extends this invariance to complex flows that stay in the flag domain.

% before: is-onto-share-counts-section-cusp
In the weight $(n+4)L-2\delta_s$ the interpolating sections have to be cusp forms. For a weight whose entries do
not increase within the two blocks and increase from the first entry to the last at one embedding, the constant
terms of every bounded section vanish.

% before: is-onto-share-counts-section-values
When the bounded sections of a weight are cusp forms, the sections that separate the base points are cusp forms
too. Their combinations take any prescribed values at the base points.

% before: is-onto-share-counts-section
The bound in the weight $F+TL+2\delta_s$ multiplies cusp forms by a bounded section of weight
$(n+4)L-2\delta_s$ that does not vanish identically. Such a section interpolates the value $1$ at the base
points, and there is at least one base point.

% before: is-onto-share-counts-interpolation
The lower count multiplies functions that separate the base points by monomials in the flat coordinates.
These functions are quotients of Poincaré series of weight $F-\delta_s$ that are nonzero at one base point
and vanish at the others.

% before: is-onto-share-counts-lower
The products of these functions with the monomials of degree at most $T$, homogenized by $w_0$, are cusp
forms of the lowered weight. Their leading terms at the base points show that they are independent.

% before: is-onto-share-counts
The ranks come from counting cusp forms on each copy. For weights $F+TL$ that satisfy the gap
conditions, the archimedean model of a copy bounds the number of independent cusp forms above and
below by multiples of binomial coefficients in $T$.

% before: is-onto-share-binomial
The upper and lower counts differ in their binomial coefficients. For $T$ large compared with $n^2$
these coefficients are close to one another.

% before: is-onto-share-ranks
Summing the counts over the copies compares the dimensions of the spaces of weights $\kappa$,
$\kappa-\delta_s$ and $\kappa+2\delta_s$. The ratio of the binomial coefficients gives the loss
$n(n+6)/T$.

% before: is-onto-step-runs
The estimates at the twisted places are taken at a point together with local data for the core of
its definite-side datum, chosen place by place. These choices come in pairs, one for the places over
$p_{de}$ and one for the others, and they must be pinned to the realization at the twisted places.

% before: is-onto-step-links-scalars
The characterization of a realization fixes its scalars by conventions read off the theta maps and the
rows at the base tuple. Each convention prescribes one scalar under a condition that allows at most one
value, so scalars satisfying all of them exist for every realization and face.

% before: is-onto-step-links
The point conditions state the compatibilities of the realization with the definite-side datum in
the form in which the point is built. Most later statements use them at the scalars of the
characterization of the realization, where they become the links.

% before: is-onto-step-runs-at-prime
A pair of runs is assembled from one run for each of the two special primes of the host. Each run
must satisfy the laws at every lift and be pinned at the twisted places whose own prime is its
special prime.

% before: is-onto-step-level-structure-congruence
The level maps of the one-map data are endomorphisms of the lattice of the core, tied to the level of
the realization through level lifts. This tie makes them multiply like the elements of the group,
modulo the level ideal, so they induce an action on the level module.

% before: is-onto-step-move-laws-situation
The level action on cusp forms changes the level structure of each framed point. On the moduli scheme
this change is made by a level automorphism, which moves the framed points; the statements below share
the moduli scheme of a PEL datum at a level, its universal object and the framed points over algebras.

% before: is-onto-step-move-laws-inverse
A level automorphism comes from twisting the universal object by a family of similitudes, and showing
that it is invertible needs the inverse family. Inverting each map modulo $m$ is not enough: the inverses
have to fit together as $m$ varies, and their multipliers have to be read off a pair of vectors with
pairing $1$.

% before: is-onto-step-move-laws-automorphism
The moduli scheme represents objects at the level, so a twist of the universal object by a family of
similitudes is classified by a map of the moduli scheme to itself. The twist by the inverse family gives
the inverse map, so every family of similitudes whose multiplier is $1$ modulo the full level gives a
level automorphism.

% before: is-onto-step-move-laws-unique
A level automorphism is determined by its matrix modulo the level ideal. Both automorphisms classify
the same twisted object, which fixes the map of the moduli scheme; the rigidity of level structures then
fixes the map of the universal abelian scheme.

% before: is-onto-step-move-laws-moves
For a given level automorphism, the move of a framed point is defined through the identification of the
universal abelian scheme with its pullback. Its existence and uniqueness and its behaviour under
composition, base change and triangular changes of frame are the properties the level action needs.

% before: is-onto-step-move-laws-framed
Combining the existence and uniqueness of level automorphisms with the moves they induce gives the laws
of an action of the automorphisms of the level module on framed points. These laws are what the level
action on sections, its integrality and its multiplicativity rest on.

% before: is-onto-step-level-structure-cusp
The action condition on a run asks for an action of the group on its cusp forms. The level action of an
automorphism of the level module moves framed points, and it has to keep cusp forms of a weight among
cusp forms of that weight.

% before: is-onto-step-level-structure
The two runs of a pair share one level structure, an action of the group of the realization on the
level modules of the core. It has to give the level of the realization and the action on the cusp
forms of every lawful run at once.

% before: is-onto-step-rational
The class sums at the twisted places come from theta series whose values lie in ray class fields of
the subfield of the host. When these ray class fields lie in the coefficient field, the class sums
are rational.

% before: is-onto-step-kernels
At a twisted place, the lattice that the export carries to the local model is spanned by theta
series of the kernels attached to the character of the own prime. Its description through
representatives of the classes of a level uses the individual kernels and their theta series.

% before: is-onto-step-kernel-unique-character
A base theta kernel comes with a character of the ideles that is equal to $\delta_v$. Its values at the
finite places on the principal elements determine it, by continuity and approximation, so two such
characters agree.

% before: is-onto-step-kernel-unique
The base theta lattice is spanned by the theta series of all base theta kernels at once. Under the
characterization all of them have the same theta series, so one kernel suffices.

% before: is-onto-step-kernel-span-move
Moving an element of the finite adelic group by a rational element and an element of the level changes
its lattice by the rational element alone. The integral coefficient vectors move with it, through the
equivariance of the family tie.

% before: is-onto-step-kernel-span-twisted
The action of the level and of the real points on the theta series uses only a few properties of a base
theta kernel. They are those of a character of finite order with its ramified places, residue characters
at those places, a Weil datum and a level vector. A tuple with these properties makes sense for any Hecke
character and any finite set of places, and so does the level of the data at those places.

% before: is-onto-step-kernel-span-twisted-base
A base theta kernel has these properties for the character of its twisted place. Most of them come from the
conditions on the characters of the level data and on the ramified sets. The ramified places need the
hypothesis on the twisted places, because the character of the tuple at the place is the character of the
place times a character of order a power of the residue characteristic.

% before: is-onto-step-kernel-span-level-vector
The test vector of a twisted kernel is a level vector, so the action of a stabilizer pair on it is the
product of characters computed in the realization's frames. At the level of the data these characters
take values in the roots of unity of the coefficient field, and on the level lifts of $1$ they cancel.

% before: is-onto-step-kernel-span-level
An element of the level changes the theta series of a twisted kernel by a scalar. The scalar is a root of
unity in the coefficient field, so it does not change the $\mathcal O_L$-span of the theta series.

% before: is-onto-step-kernel-span-real-compact
An element with finite components $1$ and compact archimedean components acts on a vacuum vector by a
character, and it moves the bases that form the pure tensors. A change of orthonormal bases is a linear
substitution of the creation variables. So the element acts on the pure tensors of every vacuum vector by
the same character and substitution.

% before: is-onto-step-kernel-span-real-fock-weights
The compact elements act on pure tensors with exponents that must be integers. They are read from the
weight of the lowest $K$-type of the model of $W_j$ at the base embedding, and from the dominant weight of
the lift. Both are the positive and negative parts of the weights of the lift, shifted by half the rank of
$V_j$ and by the numbers $e$ of the Harish-Chandra parameter.

% before: is-onto-step-kernel-span-real-fock
The theta lifts of the realization intertwine the action on the pure tensors of its test vector with the
action on the coefficient space. The real points act on pure tensors in a way that does not depend on the
vacuum vector, and the pure tensors of the test vector determine their polynomials. So the intertwining
holds for every level vector.

% before: is-onto-step-kernel-span-real-weil
The characters of the Weil datum of a twisted kernel differ from those of the realization at the base
tuple by the character of the kernel, which has finite order. On an element with finite components $1$ the
two Weil data differ by a value of this character at a determinant root, and a root that is an $n$-th power
makes this value $1$.

% before: is-onto-step-kernel-span-real
After the rational part and the level, an element of the finite adelic group still differs from a
representative by an element with finite components $1$. Such an element acts on the theta series of a
twisted kernel through the coefficient vector.

% before: is-onto-step-kernel-span
For one kernel, the theta series at all elements of the finite adelic group come from those at
representatives of the classes of the level, through the rational points and the level. So the
representatives span the same module.

% before: is-onto-step-span
The two previous lemmas together describe the base theta lattice through finitely many elements of
the adelic group. This description is the span over the classes of lattices of the own place.

% before: is-onto-step-export
The twist bound compares lattices at a twisted place inside the local model of its twist choice.
The export carries the base theta lattice and the theta lattice of the tuple into that model, with
the published inputs on cusp forms, class field theory and the discrete series.

% before: is-onto-step-numerics
The geometric estimates hold beyond a scale fixed by numerical conditions on the member. These
conditions compare the parameters of the member with its scale and its host.

% before: is-onto-step-index-conditions
The geometric estimates rest on index bounds for the cusp forms of the runs. Those bounds hold under
the index-parameter conditions, which the core ties and the laws of the runs give at large scales.

% before: is-onto-step-leaves-kappa
The theta module at a level is built in the weight $\kappa'_l$, the weight of the lift at that level.
Its entries are affine functions of the Harish-Chandra parameters of the datum, truncated at $0$, so
the bounds on the parameters become bounds on the weight.

% before: is-onto-step-leaves-theta-module
The lattices of a record object at a place are spans of cusp forms of the run used there, read
through the tie. The theta module of the run is built from its theta series and short theta series
by scaling with explicit integers and by multiplication in the ring of forms, and the estimates are
statements about these pieces.

% before: is-onto-step-leaves-weights
The theta module involves three weights of the parameter choices and the integer $T_O$. They are
related linearly, which is what the products in the theta module need to land in the right weight.

% before: is-onto-step-leaves-integral
Integrality of a cusp form is tested by valuations that are additive on products. So products of
integral forms are integral, and the theta module consists of integral forms as soon as its building
blocks do.

% before: is-onto-step-leaves-scalar
The scalars that enter the theta module are products of powers of the same integers. Each scaled
summand is therefore a multiple of one record scalar times the corresponding plain summand.

% before: is-onto-step-leaves-hypotheses-discrete
A module of cusp forms is finitely generated once its nonzero elements cannot be small at every
archimedean place at once. A lower bound for the sum of the logarithms of the norms of a nonzero element
gives this, and a discrete subgroup of a real vector space is finitely generated.

% before: is-onto-step-leaves-hypotheses-sup
The lower bound comes from integrality. A nonzero integral cusp form is evaluated at a CM point near a
point where it is nearly maximal, and the product formula bounds the mean logarithm of its supremum
norms from below.

% before: is-onto-step-leaves-hypotheses
The index-parameter conditions contain, besides their conclusion, conditions on the scalars and on
the scaled modules. Three consequences are used here: the integrality condition (H2), the spanning
of the cusp forms by short theta series, and the finiteness of the integral forms.

% before: is-onto-step-leaves-short
The span estimate needs the short theta series to form a finitely generated module. A nonzero
multiple of them lies in the integral forms, which are finitely generated.

% before: is-onto-step-leaves-products
Containment in a multiple of the short theta lattice requires the plain theta module to consist of
short theta series of the weight $\kappa$. Its summands are products of theta series with short
theta series, so this is a statement about such products.

% before: is-onto-step-leaves-move-laws
The level action of an automorphism of the level module evaluates a section at the moved framed
points. Its properties follow from the laws of the moves: they exist, commute with base change and
compose.

% before: is-onto-step-leaves-level-integral
Integral cusp forms are those with integral values at the framed points that extend over complete
discrete valuation rings. A move changes only the level structure of a framed point, so the extension
of the point serves for the moved point.

% before: is-onto-step-leaves-level-theta
The theta module is built from theta series of the weight $w_L$ and from short series. Its stability
under the level action needs that the level action keeps the theta series of a weight among theta
series.

% before: is-onto-step-leaves-level-short
On short series the level action multiplies by a scalar and moves the vector of the series. The
scalar is integral, so the $\mathcal O_F$-span of the short series of a weight is stable.

% before: is-onto-step-leaves-level
The group of the record object acts on the cusp forms through the level action of the pair of runs.
Stability of the lattices under this group comes from the stability of the integral forms and of the
theta module under level actions.

% before: is-onto-step-leaves-form-tie
At a single place the supremum, span and containment estimates depend only on the tie between the
record object and the cusp forms of a run. With the lemmas above they follow from the lattice
property of the integral forms, the spanning by short theta series and (H2).

% before: is-onto-step-leaves-ties
The record object is tied to the run used at each place. Besides the integral forms and the theta
module, the tie identifies the short theta lattice, the record scalar, the action of the group and
the lattice $\Theta_0$ of traces of orbit families.

% before: is-onto-step-leaves-move-isometry
Right multiples are built from a lift of an automorphism of the level module to the lattice. A lift
that respects the pairing modulo the level, and is invertible modulo it, comes from the two coprime
halves $M_b$ and $\bar M_b$ of the level ideal. Hensel lifting and the Chinese remainder theorem extend
it to every modulus.

% before: is-onto-step-leaves-move-character
An automorphism of determinant one acts on short series through a scalar that is multiplicative in
the automorphism. Over a finite ring whose residue fields have at least four elements, every such
character of the general linear group is trivial on the matrices of determinant one.

% before: is-onto-step-leaves-move-det-character
The level module is free over the finite ring $\mathcal O_E/M_b$, whose residue fields have at least
four elements when the level is odd and prime to $3$. So the previous lemma applies to the
automorphisms of the level module.

% before: is-onto-step-leaves-move-const
The level action of an automorphism carries a short series at a vector to a multiple of the short
series at a vector of which the first is a right multiple. The scalar depends only on the
automorphism and the weight.

% before: is-onto-step-leaves-move-last
Short series are cusp forms when the last positive entries of their weight are positive. For the
weights $\kappa'_l$ these entries are differences of the Harish-Chandra parameters, which the parameter
choices make large.

% before: is-onto-step-leaves-move
The data estimate compares traces of orbit families with short theta series. This needs short series
at every allowed vector, right multiples by every automorphism, and the effect of the level action of
a determinant-one automorphism on short series.

% before: is-onto-step-leaves-det-congruence
The level condition compares a lift of an automorphism of the level module with an adelic level lift
only on $\bar M_b\Lambda$ and modulo $M$. In coordinates adapted to the lattice, such a congruence
compares the two determinants modulo the closure of the level ideal.

% before: is-onto-step-leaves-det
The level action moves short series without a scalar factor only for automorphisms of determinant
one. The elements of $P_0$ are unitriangular at the places of the level, so their automorphisms of
the level module have determinant one.

% before: is-onto-step-leaves-data
The data estimate concerns the lattice $\Theta_0$ of traces of orbit families. A trace over the
group $U$ with the character $\psi$ lands in the $\psi$-eigenspace and multiplies eigenvectors by
$|U|$, and every short theta series is the common image of an orbit family under the level actions.

% before: is-onto-step-leaves
A record object is tied at every place to the cusp forms of the run used there. Through this tie the
index-parameter conditions give the first four geometric estimates.

% before: is-onto-step-index-bound-conclusion
The index bound draws on one part of the index-parameter conditions, their conclusion. It bounds,
at each completion of the coefficient field of the core, the integral cusp forms against the theta
module of a run inside its ring of forms.

% before: is-onto-step-index-bound-localize
The conclusion of the index-parameter conditions is stated after completion, while indices are lengths
of localizations. Over a discrete valuation ring the completion can only lengthen a quotient.

% before: is-onto-step-index-bound-passage
With the comparison of lengths, the conclusion becomes a bound on indices of the integral cusp forms
against those of them that lie in the theta module. It also shows that these span the cusp forms.

% before: is-onto-step-index-bound-base-change
The record object lives over the coefficient field $L$, the cusp forms over the smaller field of the
core. An index of lattices multiplies by the ramification index when the lattices are carried to the
larger field.

% before: is-onto-step-index-bound-extend
A bound on indices for all subspaces is equivalent to a single bound for one auxiliary lattice, and
that bound is carried to the larger field by the previous lemma. Enlarging the second lattice only
lowers the index.

% before: is-onto-step-index-bound-sum
The remainders over $L$ are the remainders over the field of the core, weighted by ramification
indices. Their weighted sum is at most the sum over the field of the core.

% before: is-onto-step-index-bound-exponent
The conclusion of the index-parameter conditions applies to integers above $T$ times the exponent
function of the run. The exponent $E(v,j)$ of the record object is built from the normalized exponent
of the depth data, and the two are compared here.

% before: is-onto-step-index-bound-ties
Being recorded ties the lattices and exponents of the record object to the run used at each place.
The index bound uses the form tie, the depth tie and the values of a few numbers.

% before: is-onto-step-index-bound-remainder
The previous lemmas combine into remainders over $L$ for one run. They bound the index of the lattice
spanned by the integral cusp forms against the lattice spanned by those in the theta module.

% before: is-onto-step-index-bound
The index bound measures the lattice of the record object against its theta lattice, place by place.
Its remainders come from the conclusion of the index-parameter conditions, one family for the places
over $p_{de}$ and one for the others.

% before: is-onto-step-depth
The depth estimates compare the local numbers of the record object with the local terms $F_0$ of the
numerical bounds. They rest on three local statements at a prime, for the split places, the
Hodge-Tate lattices and the unramified companions.

% before: is-onto-step-compare
The rank at the twisted places needs a lattice of the run inside the base theta lattice. Away from
the own prime of the place the two are compared through the class sums of the base tuple.

% before: is-onto-step-theta-j-ties
Two more lattices of the record object are tied to the lift. They are the base theta lattice and the
trace lattice of the run used at the place, pulled back along the isomorphism of the tie.

% before: is-onto-step-theta-j-low-types-vectors
An element of a multiplicity space is continuous in each variable only for a Sobolev norm, so the
argument works with smooth vectors and their derivatives. The models of discrete series are unitary
representations of $U(J_\epsilon)$ moved to the unitary group of a Gram matrix by a change of basis, and
the smooth vectors are taken for the moved group.

% before: is-onto-step-theta-j-low-types-smooth
The derivatives of smooth vectors are skew-adjoint, and the Gårding vectors are strongly smooth.
These two facts make the smooth vectors dense, the derivative linear in the Lie algebra, and a
Sobolev-continuous functional a finite sum of continuous functionals of derivatives.

% before: is-onto-step-theta-j-low-types-translates
In an irreducible unitary representation the translates of one nonzero vector span a dense subspace.
For a Sobolev-continuous functional this density has to be carried to the smooth vectors, and smoothing
by Gårding operators does that.

% before: is-onto-step-theta-j-low-types
A vector of the rational structure is read through its automorphic form, evaluated on lowest $K$-type
vectors. Elements of a multiplicity space of a model of discrete series lose nothing under this
evaluation.

% before: is-onto-step-theta-j-determined
Combined with the injectivity of the rational structure, the evaluation on lowest $K$-types determines
a vector. This is what lets an identity of theta series identify vectors.

% before: is-onto-step-theta-j-fg
The base theta lattice is spanned by theta series at all elements of the finite adelic group. By the span
over the classes, finitely many representatives suffice, and at each of them the integral coefficient
vectors form a finitely generated module.

% before: is-onto-step-theta-j-type
The record object asks its lattice $\Theta_J$ to lie in the $\psi$-eigenspace of its group. Through the
tie this is a statement about the base theta lattice and the subgroup $P_0$.

% before: is-onto-step-theta-j-idempotent
The projection $e_j$ of the record object is the good operator moved along the tie, and the good
operator is the image of the sum of the idempotents of the characters of the family. These characters
are absolutely irreducible and take values in a number field. Their idempotents are central and
orthogonal over any field of characteristic $0$, by Schur's lemma and the orthogonality relations.

% before: is-onto-step-theta-j
The lattice $\Theta_J$ of the record object is the base theta lattice transported along the tie. The
span over the classes and the comparison away from the own prime make it finitely generated and of
the right type, and give the containment of $\Theta_0$ in it after localization at the place.

% before: is-onto-step-run-model
The share bound counts cusp forms on the two-level data of
Definition~\ref{is-onto-share-two-level}. The run used at a place gives such data, with the group
acting through the level structure of the pair of runs.

% before: is-onto-step-run-model-exists-components-lattice
A self-dual module for the Hermitian form comes from the lattice of a component when the trace form on the
lattice is perfect. The different of the field converts the one duality into the other, and an ideal
whose product with its conjugate, the different and $\xi_0$ is the unit ideal turns the lattice into a
self-dual module.

% before: is-onto-step-run-model-exists-components-flag
The laws of a core give its components the absolute value of the host element and, through the family
uniformization, a perfect trace form on their lattices. The ideal of the host has the norm that the lemma on
lattices needs, so the components extend to flag data once their groups contain the principal congruence
subgroup of the level.

% before: is-onto-step-run-model-exists-components-congruence-exponential
Two periods of the family uniformization with the same point give two exponentials onto the same abelian
variety. They differ by a linear change of coordinates, which carries one period lattice to the other
through an isometry of the lattice that respects the level structures.

% before: is-onto-step-run-model-exists-components-congruence-rigidity
An isometry of a period lattice that is congruent to $1$ modulo the level and matches two periods by a
linear map gives the same point. The isometry that the exponentials give is congruent to it modulo the
level, and their quotient has finite order, so Minkowski's argument makes them equal.

% before: is-onto-step-run-model-exists-components-congruence-orbit
An element of the principal congruence subgroup moves the periods at a point of the flag domain by an
isometry of the lattice of the component. By the rigidity of the family uniformization the two framed points
coincide up to a scalar, so the moved point lies in the orbit of the component's group.

% before: is-onto-step-run-model-exists-components-congruence-roots
Since the element moves every point into its orbit, a scalar multiple of it lies in the group. The group
has finite index in the stabilizer of the lattice and acts freely on the flag domain, which forces the
scalar to be a root of unity.

% before: is-onto-step-run-model-exists-components-congruence-lattice
The dictionary relates rational points and adelic points of the unitary group. Comparing them needs the
rational vectors of an adelic translate of the standard lattice, a lattice on which a rational element
inherits the congruence of an element of the level.

% before: is-onto-step-run-model-exists-components-congruence-root
A rational element of the dictionary's level is conjugate at the archimedean places to the scalar multiple
of the congruence element. Both are congruent to $1$ modulo the level, on two lattices, and this forces the
root of unity between them to be $1$.

% before: is-onto-step-run-model-exists-components-congruence
The extension to flag data and the action of a component's group on its self-dual module both need the
principal congruence subgroup of the level inside that group. The family uniformization and the dictionary
of the core provide it.

% before: is-onto-step-run-model-exists-components-back-periods
A matrix that carries the periods of one point to upper triangular combinations of the periods of another
determines the second point from the first. The frames then differ by an isometry and an upper triangular
matrix.

% before: is-onto-step-run-model-exists-components-back
An element of a component's group moves a point of the flag domain to another point of the same component.
The family uniformization compares the two framed points, and the comparison is an element of the principal
congruence subgroup up to an upper triangular matrix.

% before: is-onto-step-run-model-exists-components
A run model is built from the components of the unitary data of a run. Its axioms ask the components
to share their signature data, to have trivial characters on their arithmetic groups, and to carry
self-dual modules on which these groups act trivially modulo the level ideal.

% before: is-onto-step-run-model-exists-transport
The group $G$ acts on the components through the level changes of the pair of runs. Each level change
carries a component to another by a transport matrix, unique up to the arithmetic groups, and the
cusp forms are read compatibly along it.

% before: is-onto-step-run-model-exists
Every lift has a run model at the points of a family. It is built from the unitary data of the run
used at one place.

% before: is-onto-step-model-reading
The cusp forms of a run model are compared with the rational structure of the lift. The comparison
is a semilinear bijection that respects the actions of the group, so the dimensions of eigenspaces
can be read on either side.

% before: is-onto-step-model-level-places
The level hypothesis is a congruence at every prime over $g_0$ or $g_1$, while the group of the
realization sees only the $f$ places of $P_g$ over each of them. The other primes over $g_0$ and $g_1$
are the complex conjugates of these places.

% before: is-onto-step-model-level-half
At a place of $P_g$ the congruence for a run model reduces to one on the level module of the core,
through the identification of (4) of Definition~\ref{is-onto-step-run-model}. The elements of
$\Gamma_c$ act trivially modulo the level ideal, so only the level action enters.

% before: is-onto-step-model-level-conj
At the complex conjugate of a place the congruence follows from the one at the place itself. The
module $L'$ is self-dual, and the adjoint of an isometry moves the ideal to its conjugate.

% before: is-onto-step-model-level-coordinates
The level action lives on a lattice of the core, while the level lifts of elements of $G$ are adelic
matrices that stabilize the closure of the lattice of $W_i$. Coordinates in a basis of that closure,
reduced modulo a prime $P$, turn a congruence between the two into an equality of vectors over the
residue field at $P$.

% before: is-onto-step-model-level-condition
The level condition on a pair of runs ties the maps $\ell_i$ to the group $G$ of the realization. It
says that each $\ell_i(g)$ is congruent, on part of the lattice of the core, to a level lift of $g$, an
adelic element of the unitary group of $W_i$ whose reductions at the places of $P_g$ are the components
of $g$.

% before: is-onto-step-model-level-action
The congruence at a place of $P_g$ needs the level action of an element of the half-unipotent subgroup
to be unipotent modulo that place. This holds for every element that is upper unitriangular there.

% before: is-onto-step-model-level
The dimension statements for two-level data need the level hypothesis for the subgroups to which
they are applied. For a run model it holds for the half-unipotent subgroups once $g_0$ and $g_1$
exceed $3r$, three times the rank.

% before: is-onto-step-model-arch
In the indefinite case the ranks of the cusp forms come from the archimedean condition. For run
models it holds at the parameter $T_O$ of the member at large scales.

% before: is-onto-step-models
The share bound uses, for every lift, two-level data whose cusp forms read the rational structure of
the lift and satisfy the hypotheses of the dimension statements. Lemmas~\ref{is-onto-step-run-model-exists},
\ref{is-onto-step-model-reading}, \ref{is-onto-step-model-level} and~\ref{is-onto-step-model-arch} give
such data.

% before: is-onto-step-bad-weight-whittaker
The killed part is bounded by the count of Lemma~\ref{cv-bad-part}, which needs the type of the face in
Whittaker form. The level tie among the point conditions gives it: the subgroup of the face is the
upper unitriangular subgroup, and its character is a product of additive characters of the
superdiagonal.

% before: is-onto-step-bad-weight-regular-determinant
A space of functions embeds into copies of the regular representation of a finite group when a single
function has translates that are independent in a strong sense: their group determinant is a
nonzerodivisor. The number of copies is then bounded by the dimension of a space of products.

% before: is-onto-step-bad-weight-regular-cocycle
The group acts on the cusp forms through pullbacks twisted by scalars. These scalars can be chosen to
form a cocycle, so that the twisted pullbacks define an action on polynomials in one variable over the
holomorphic invariant functions.

% before: is-onto-step-bad-weight-regular-one-lift
The function with independent translates is built from cusp forms that separate the points of one
orbit. Such a form is nonzero at a point without contraction and vanishes at the other translates of
that point.

% before: is-onto-step-bad-weight-regular-translates
Summing such forms over representatives of the orbits gives a single family whose matrix of translates
is, at a suitable point of each component, a multiple of a permutation matrix. Its determinant is then
nonzero.

% before: is-onto-step-bad-weight-regular
In the indefinite case the cusp forms of the model are compared with the regular representation of a
half-unipotent subgroup. The number of copies needed is bounded by the dimension of a space of cusp
forms of a shifted weight.

% before: is-onto-step-bad-weight-model
The count needs, for each of the two half-unipotent subgroups $H_k$, an embedding into copies of the
regular representation whose size is bounded by the dimension of an eigenspace. Both cases of a model
give it, with the constant $2$.

% before: is-onto-step-bad-weight-transfer
The model reads the rational structure of the lift through an embedding $\sigma$ of the coefficient
field, while the count is made through $\sigma_0$. An automorphism of $\mathbf C$ carries one base
change to the other.

% before: is-onto-step-bad-weight-share
The count bounds the killed part by the shares of the tuples of labels that are not separated. With
Green's description the weights are explicit, and the share is the measure of a set of tuples of
polynomials.

% before: is-onto-step-bad-weight-share-sharper
The first term of the bound comes from the tuples with a component that is not squarefree. The
squarefree labels carry all but $2/g_k$ of the total weight at each place, which gives a smaller
constant than the one stated.

% before: is-onto-step-bad-weight
The bad-place bound compares the part of the rational structure of a lift killed by the good operator
with the whole. Green's description of the cuspidal characters, read on the model of the lift, makes
that part small.

% before: is-onto-step-bad-numerics-bound
The inequality of Lemma~\ref{is-onto-step-bad-weight} bounds a sum of five terms by $\delta$. At large
scales each term is at most a power of $x$ times $e^{-x^{\varepsilon/6}}$, which is much smaller than
$e^{-x^{\varepsilon/7}}$.

% before: is-onto-step-bad-numerics
Lemma~\ref{is-onto-step-bad-weight} bounds the killed part by $\delta$ under numerical conditions on the member. Along
the standard families these conditions hold for $\delta=\exp(-x^{\varepsilon_z/7})$, and this
$\delta$ is at most $\frac14$.

% before: is-onto-step-exp-dominant
The weights of the exponent dictionary are shifts of the weights that the Harish-Chandra parameters
give on the spaces $V_J$. Those weights are explicit in the parameters, and the interlacing and the
gaps of Definition~\ref{th-hc-parameter} make them integral and dominant.

% before: is-onto-step-exp-dictionary
The exponent bound compares the local exponents of the block lattices of a point with exponents
computed from weights. The weights and a parameter record are read off the size data through the exponent
dictionary.

% before: is-onto-step-exp-cut-exponent
The exponents of the weights compare two integral forms of the same representation of $GL_N$, the dual
Weyl module and the Weyl module over $\mathbf Z$. Away from a prime $l$ only the power of $l$ matters.

% before: is-onto-step-exp-cut-attained
Over $\mathbf Q$ the dual Weyl module and the Weyl module coincide. So some nonzero integer carries the
first lattice into the second, and the least exponent of $l$ that works is attained.

% before: is-onto-step-exp-cut-beta
The projector used to cut a weight into two blocks has coefficients $j!/(k!(k+j+1)!)$. A single integer,
prime to $l$ up to a power $l^{\lfloor\log_lw\rfloor}$, clears all these denominators at once.

% before: is-onto-step-exp-cut-grid
The cut is made by an operator built from the action of $\mathfrak{gl}_N$ on endomorphisms. It is a
product of truncated series in a raising and a lowering operator, one for each root joining the two
blocks.

% before: is-onto-step-exp-cut-adjoint
The grid projector is applied to an endomorphism of the dual Weyl module, so the action of
$\mathfrak{gl}_N$ on endomorphisms needs the same properties as the action on the module. Its weights
are differences of weights of the module.

% before: is-onto-step-exp-cut-invariance-rank-one
The invariance of the grid projection is proved one simple root at a time. For a single root the
truncated series is an $\mathfrak{sl}_2$ computation, and the bound on the weights makes the truncation
exact.

% before: is-onto-step-exp-cut-invariance-rank-two
Two roots that share an index interact through a third root. The projectors of the two roots, with
consecutive indices, together kill the effect of that third root.

% before: is-onto-step-exp-cut-invariance
An endomorphism that commutes with the Levi subalgebra $\mathfrak{gl}_p\times\mathfrak{gl}_q$ becomes
one that commutes with all of $\mathfrak{gl}_{p+q}$ after the grid projector is applied. The weights
of the action on endomorphisms are bounded by $2K$, twice the length of the truncated series.

% before: is-onto-step-exp-cut-scalar-root
The value of the grid projector at a highest weight vector is computed root by root. For one root it
is again an $\mathfrak{sl}_2$ computation, in which the truncated series sums to the inverse of a
shifted weight difference.

% before: is-onto-step-exp-cut-scalar
On the highest weight vector the grid projector acts, up to the terms that the endomorphism kills, by
the inverse of an explicit product of differences of shifted weights. This product is a nonzero
integer.

% before: is-onto-step-exp-cut-integrality
The denominators of the grid projector are bounded by the factorial divisibility of
Lemma~\ref{is-onto-step-exp-cut-beta}. Each root joining the two blocks costs one factor $\Lambda$.

% before: is-onto-step-exp-cut-module
The row operators on the dual Weyl module satisfy the hypotheses of the three statements on the grid
projector. Schur's lemma for this irreducible module identifies the projected endomorphism with a
scalar.

% before: is-onto-step-exp-cut-levi-duality
The denominators of an endomorphism of the dual Weyl module are read off through the Fock pairing.
Over the integers the dual Weyl lattice and the Weyl lattice are dual to each other for the pairing
normalized by the highest weight vector.

% before: is-onto-step-exp-cut-levi-substitution
Restricting a polynomial to matrices of a special shape keeps it in a dual Weyl module of smaller
size. The substitution only has to respect the column operations that define these modules.

% before: is-onto-step-exp-cut-levi-restriction
Restriction to block diagonal matrices relates the dual Weyl module of the whole weight to those of
the two blocks. It respects the lattices, the highest weight vectors and the row operators of each
block, and it kills the vectors whose weight is too small on the first block.

% before: is-onto-step-exp-cut-levi
The endomorphism to be projected comes from the two blocks. It is built from the restriction of
polynomials to block diagonal matrices, and its denominators are those of the two blocks.

% before: is-onto-step-exp-cut-blocks
Together these give the inclusion for the whole weight from the inclusions for the two blocks. The
projected endomorphism is the identity up to an integer, and its denominators cost $\Lambda^{pq}$.

% before: is-onto-step-exp-cut
The exponent of a Weyl lattice at a prime measures how far the dual lattice is from it. Cutting the
weight into two blocks changes the exponent by at most a term controlled by the spread of the
shifted weights.

% before: is-onto-step-exp-sum-shift
The weights of the dictionary differ from the natural weights by constants. Multiplication by the
determinant identifies the modules of a weight with those of its shift, so the exponent does not see the
constant.

% before: is-onto-step-exp-sum-localize
The lattices attached to a family of lattices in $K^n$ are defined by conditions at every finite place.
At one place they are the standard lattices moved by the matrix of a local basis.

% before: is-onto-step-exp-sum-local
At a place over $l$, the exponents of the factors bound the power of $l$ that carries the dual Weyl
lattice of a lattice into its Weyl lattice. The local description of the two lattices reduces this to
the standard lattices, factor by factor.

% before: is-onto-step-exp-sum-lattices
The presentation of the lift forms by algebraic modular forms carries two lattices at each
representative. A family tie identifies them with the dual Weyl lattice and the Weyl lattice of the
lattice of the representative.

% before: is-onto-step-exp-sum
The block lattices of the point lie, at every finite place, in the lattices of the weights of the
dictionary. This bounds their local exponents by sums of exponents of Weyl lattices.

% before: is-onto-step-exp-model
At the twisted places whose residue characteristic is in the range of the weights, the exponent
bound uses another count. It comes from Hida settings built with the Galois representations.

% before: is-onto-step-exp-bound-lift-exponent
The exponent function $J_0$ of the numerical bounds of a point measures, place by place, how far the
cut lattice is from the Weyl cut lattice. It is normalized by the ramification index, so that a power of
the residue characteristic counts once.

% before: is-onto-step-exp-bound-lift
An inclusion by a power of the residue characteristic bounds the lift exponent by that power. This
turns the inclusions of the block lattices, and those that the exponent models give, into bounds on
$J_0$.

% before: is-onto-step-exp-bound-weights
The parameter records and weights of the exponent dictionary are explicit numerical data. The
estimates below use only the inequalities listed here, which relate the sizes of the member to the
values $U_\tau(a)$ of the weights.

% before: is-onto-step-exp-bound-recursion
The recursive exponent chooses, prime by prime, between two bounds. At small, excluded and special
primes it is the sum of the Weyl exponents; at the other primes it is the coarse weight bound, which
counts the residues of the weights near the block of each thick element.

% before: is-onto-step-exp-bound-numbers
The comparisons below rest on a few inequalities between the numbers of a parameter record. They
follow from the hypotheses and largeness by short computations.

% before: is-onto-step-exp-bound-total-rank-two
The cutting of rows below starts from matrices of size at most one, and there the Weyl lattice and its
dual agree. They are the cases without pairs of rows of an explicit inclusion in size at most two,
where the exponent is the largest $l$-adic valuation of a binomial coefficient.

% before: is-onto-step-exp-bound-interval
When all the values $\lambda_a+N-1-a$ lie in an interval shorter than $l$, the Weyl lattice and its
dual agree after inverting the integers prime to $l$. Cutting off one row at a time with the exact cut
bound gives this, with a bound for longer intervals.

% before: is-onto-step-exp-bound-beyond
Beyond the height $h$ every Weyl exponent of the weights vanishes. For a thick element the values of
the weight lie between $0$ and $h$, and for the others they lie in an interval of length at most $2b'$.

% before: is-onto-step-exp-bound-compare
The recursive exponent is at least the Weyl exponent sum outside the range of the weights, and equals
the coarse weight bound inside it. So a bound on $J_0$ by either of these is a bound by the recursive
exponent.

% before: is-onto-step-exp-bound-tail
The sum bound of the record object has the form $30TmbL^2+2mb$. The terms of the total bound beyond the
main term are absorbed into $2mb$ once the parameter record has rows for a big pair.

% before: is-onto-step-exp-bound-total-numbers
The sums below split the primes at $l_J$ and at $h$, and they bound sums over primes by elementary
estimates in $b$ and $L$. These need a few more inequalities between the numbers of a parameter
record, again short computations from the hypotheses and largeness.

% before: is-onto-step-exp-bound-total-primes
Sums of $\log p$, of $\log p/p$ and of $1/p$ over the primes up to a bound enter the total bound.
Chebyshev's bound and Mertens' first theorem, in elementary forms, give the constants needed.

% before: is-onto-step-exp-bound-total-excluded
At an excluded prime the recursive exponent is the whole Weyl exponent sum, so the excluded primes
beyond $l_J$ must be few. Each of them divides an integer of a short interval below $l_J^2$, and such
an integer has at most one prime factor beyond $l_J$.

% before: is-onto-step-exp-bound-total-no-window
Beyond $l_J$, at the primes that are not excluded, the window counts vanish. The depth of a thick
element is then $q_\tau h$, which keeps its coarse bound small.

% before: is-onto-step-exp-bound-total-exponents
The total bound uses four kinds of information about the Weyl exponents of the weights. Outside $T_W$
the width controls them, and on $T_W$ the height $h$ does. For a thick element and $l\geq2\eta+t$, a cut
after the first $k-1$ rows leaves $\log(3h)$ only $k(m+t)$ times.

% before: is-onto-step-exp-bound-total-real
The coarse bound of a thick element involves the integer parts $h_0$ and $i_\tau(l)$. Replacing them by
the real numbers they come from, and $k-1$ by $k$, gives an upper bound whose terms are simple
functions of $l$.

% before: is-onto-step-exp-bound-total-weight-bound
At the primes beyond $l_J$ that are neither excluded nor special, the recursive exponent is the coarse
weight bound. For a thick element its term is at most the real coarse bound.

% before: is-onto-step-exp-bound-total-cost
A single prime contributes at most $fN_r^2L$ to the weighted sum of the Weyl exponent sums. This
handles the excluded and the special primes, which are few.

% before: is-onto-step-exp-bound-total-thin
For an element outside $T_W$ only the primes up to its width contribute, and each contributes at most
$\binom{N_r}2$ times the logarithm of the width. The prime number bound turns this into a bound linear
in $b'$.

% before: is-onto-step-exp-bound-total-small
For a thick element and a prime $l\le l_J$, the Weyl exponent times $\log l$ is at most $2mbL/l$.
Summing $1/l$ over these primes costs a factor $\log l_J\le L$.

% before: is-onto-step-exp-bound-total-large
Beyond $l_J$ the term of a thick element is capped by its real coarse bound, and only the primes up to
$h$ contribute. Over these primes each part of that bound sums to a multiple of $mbL^2$ or of
$k(S+1)(L+2)$.

% before: is-onto-step-exp-bound-total
Summed against $\log l$ over a set of primes, the recursive exponents have an explicit total bound. Its
main term is $30mbL^2$ for every thick element.

% before: is-onto-step-exp-bound-sum
The special set of the dictionary contains the prime divisors of $Q_S$, which are small. Removing them
does not change the recursive exponents, so the total bound applies, and the tail bound puts it in the
form of the record object.

% before: is-onto-step-exp-bound-window
In the range of the weights, no value of a weight outside its block is congruent modulo $l$ to a point
near the block. Such a value differs from such a point by a nonzero amount below $l$, except for a
thick element, where a difference divisible by $l$ would make $l$ excluded.

% before: is-onto-step-exp-bound-frobenius
The exponent models link the Hida settings to the weights twisted by Frobenius data. In the range of
the weights the window counts vanish, so the coarse weight bound does not see these data.

% before: is-onto-step-exp-bound-hida
The exponent models are families of Hida settings: modules of forms with values in tensor products of
dual Weyl modules, with a Hecke algebra and a localization. When a family of them covers the block
data at a twisted place, their inclusions of forms give the inclusion of the block lattices, and their
constituents carry weights to be compared with those of the dictionary.

% before: is-onto-step-exp-bound-link-depth-series
The linkage of a weight is stated through lattices over power series rings, on which the operations
below act. The substitution $\iota$ replaces the variable by $\pi_0$, a product of the series
$(1+X)^d-1$ over the roots of unity $d$ in $\mathbf Z_l$, and a unit series describes how
$\mathrm{act}_u$ moves $\pi_0$. On reductions two substitutions remain: $\gamma_v$, which moves $X$ by a
multiple of $X^2$, and $X\mapsto X^l$.

% before: is-onto-step-exp-bound-link-depth-data
A linked weight comes with Witt lattice data: at every label a lattice with a matrix $P_i$ for Frobenius
and a matrix $D_i$ for the action of $u$, and rank-one lattices at the $t$ positions of the block. The
height bounds the divisibility of $P_i$ by powers of $X-l$, and the local height describes the same
matrices after the local exponentials. The frame ties the reductions $\Phi_i$ and $\Gamma_i$ to the
rank-one lattices at the positions of the block.

% before: is-onto-step-exp-bound-link-depth-reflected
The labels of a linkage come in pairs exchanged by a reflection about $h$, which carries the window of
a label to the window of its partner. Only one label of each pair has to satisfy the linked conditions,
and the lattice data carry the sequences of the labels or of their partners. The window counts are
taken here with an arbitrary width $\kappa$.

% before: is-onto-step-exp-bound-link-depth-content
Two products of linear factors that agree in characteristic $l$ have the same roots there. This turns
the polynomial identity of a linkage into an equality of residues modulo $l$.

% before: is-onto-step-exp-bound-link-depth-shift
The residues of a shifted block lie among those of a set of integers that contains the block itself.
When the shifts are small, a shifted entry can only be matched by an element of the set congruent to a
point near the block, and the window count counts these elements. This bounds the total shift.

% before: is-onto-step-exp-bound-link-depth-arith
The depth bound comes from comparing two sequences with the same sum. Below the block their difference
is the depth, after the block each entry of the second sequence is at most $h$, and on the last $t$
rows the shifts enter.

% before: is-onto-step-exp-bound-link-depth-chain
For two linked pairs of sequences, the residues of the shifted blocks lie among those of the primed
sequences, and so among those of the unprimed ones. The shift bound applies to both pairs, and the
arithmetic lemma combines the two bounds; the shifts of the second pair enter with the factor $l$.

% before: is-onto-step-exp-bound-link-depth-datum
The reductions of the lattice data are matrices over a field $F$ and over $F[[X]]$ with a block
structure. The bound uses only the eigenvalues of $\Theta_i$, its diagonal entries at the positions of
the block, and the orders of the diagonal entries and of one minor of $\Phi_i$.

% before: is-onto-step-exp-bound-link-depth-assembled
From the weight properties the hypotheses of the chain lemma follow. The orders of the diagonal entries
of $\Phi_i$ bound the shifts, the eigenvalues of $\Theta_i$ give the residues, and the minor bounds the
sum over the last rows.

% before: is-onto-step-exp-bound-link-depth-exchange
With the shifts bounded by $\lfloor h/(l-1)\rfloor$, the chain lemma applies at $\tau$ and at
$\sigma(\tau)$. The depth at $\tau$ is then bounded by the window counts at $\tau$ and at
$\sigma(\tau)$.

% before: is-onto-step-exp-bound-link-depth-matrices
The Frobenius matrices have inverses up to a power of $X$, and these inverses inherit the block
structure. On a block of size one, the two diagonal entries multiply to that power of $X$.

% before: is-onto-step-exp-bound-link-depth-reduction
Reduction to $F$ kills $l$, so $X-l$ becomes $X$, and a power $1+lr$ of $u$ determines $r$ modulo $l$ as
a multiple of $c$. By these two facts the conditions on the lattices become conditions on their
reductions.

% before: is-onto-step-exp-bound-link-depth-rank-one
At a position of the block the lattice data have rank one. There the height determines the Frobenius
entry up to a unit and the reduced constant term of $D_i$.

% before: is-onto-step-exp-bound-link-depth-frobenius-entry
The frame relates a diagonal entry of $\Phi_i$ to the rank-one Frobenius entry, up to units and powers
of $X$ in the Laurent series. Comparing orders of vanishing gives the order of the diagonal entry.

% before: is-onto-step-exp-bound-link-depth-base-change
The matrices $\Gamma_i$ and $\Phi_i$ are the reduced matrices of the lattice in a changed basis $g_i$.
The change of basis keeps the characteristic polynomial of the coefficient of $X$ in $\Gamma_i$, the
existence of an inverse up to $X^h$, and the divisibility of minors.

% before: is-onto-step-exp-bound-link-depth-theta-entry
The second identity of the frame computes a diagonal entry of $\Gamma_i$. Its coefficient of $X$ comes
from the first two coefficients of the factors.

% before: is-onto-step-exp-bound-link-depth-unit-series-binomial
The unit series is built from the series $(1+X)^c$ for $c\in\mathbf Z_l$ and the substitutions they
define. Their laws reduce to identities between binomial coefficients and to a first-order differential
equation that $(1+X)^c$ satisfies.

% before: is-onto-step-exp-bound-link-depth-unit-series-roots
The series $\pi_0$ is a product over the roots of unity of order dividing $l-1$ in $\mathbf Z_l$. Three
facts about these roots are used: they lift the nonzero residues, one of them lifts a generator, and
their product is $-1$.

% before: is-onto-step-exp-bound-link-depth-unit-series-pi
The substitution $\iota$ of $\pi_0$ carries $X-l$ to $q_0$. The series $\pi_0$ is fixed by the
substitutions of the roots of unity, and for odd $l$ it is $X^{l-1}$ times a unit.

% before: is-onto-step-exp-bound-link-depth-unit-series-division
A series in $X$ can be divided by a unit times $X^n$, with remainder of degree less than $n$. Iterating
the division writes every series in terms of the powers of $wX^n$ below $X^n$, and a series fixed by a
substitution that fixes $wX^n$ uses only the first of them.

% before: is-onto-step-exp-bound-link-depth-unit-series-invariants
For $wX^n=\pi_0$ the division shows that $\iota$ is injective. The series fixed by the roots of unity
are the series in $\pi_0$.

% before: is-onto-step-exp-bound-link-depth-unit-series-q
The series $q_0$ and $q$ differ by a unit. So divisibility by a power of $q$ after $\iota$ descends to
divisibility by the same power of $X-l$.

% before: is-onto-step-exp-bound-link-depth-unit-series
The constant term of the unit series enters the coefficient of $X$ in $\Gamma_i$. Its reduction is that
of $c$.

% before: is-onto-step-exp-bound-link-depth-alpha
At a position of the block, the order of the diagonal entry of $\Phi_i$ is the shifted weight. It is at
least $0$ by the frame, and at most $h$ because $\Phi_i$ has an inverse up to $X^h$.

% before: is-onto-step-exp-bound-link-depth-properties
The weight inputs give a weight datum with the weight properties. The eigenvalues of $\Theta_i$ come
from the lattice $N$, its diagonal entries from the rank-one lattices and the unit series, and the
minor from the height of $N$.

% before: is-onto-step-exp-bound-link-depth-inputs
The depth bound follows from the weight inputs. The sum of the last $t$ entries of the decreasing
sequence $U'_\tau$ is at most its sum over any $t$ entries, which gives the condition on the minor.

% before: is-onto-step-exp-bound-link-depth-norm-division
The height over $O$ is read off at the linear factors $X-c$ with $c$ in the ideal of definition. Over an
adically complete domain such a factor divides with a constant remainder, so it is prime, and powers of
distinct such factors that divide a series divide it together.

% before: is-onto-step-exp-bound-link-depth-norm-descent
Divisibility after the substitution $\iota$ and after a local exponential is brought back to $O[[X]]$ in
three ways. A matrix inverse up to a power of $q$ descends to one up to the same power of $X-l$; the
factors of $q$ at the roots of unity together give divisibility by $X-l$; and a local exponential at
$x$ detects the power of $X-x$.

% before: is-onto-step-exp-bound-link-depth-norm-local
A local exponential at $\zeta^b-1$ sees only one of the linear factors of $q_0$; the others become units.
It also reads the constant term of $\iota(y)$ from the class of $y$ modulo $X-l$.

% before: is-onto-step-exp-bound-link-depth-norm-matrices
Four facts on matrices and units finish the comparison. They give the minors of a diagonal matrix of
powers, the constant terms of a matrix identity twisted by $X\mapsto\upsilon X$, the characteristic
polynomial of a matrix conjugate to a diagonal one, and the first-order expansion of $u^{-m}$.

% before: is-onto-step-exp-bound-link-depth-norm
The inputs of the period change give the local height of the lattices, read through the local
exponentials. The height over $O$, which the weight inputs require, follows from it.

% before: is-onto-step-exp-bound-link-depth-period
Through the heights, the inputs of the Frobenius period change give the weight inputs. So the depth
bound holds under them.

% before: is-onto-step-exp-bound-link-depth-orient
In a linkage only one label of each reflected pair satisfies the linked conditions, and the lattice data
carry the sequences of the labels or of their partners. The reflection about $h$ transfers the linked
conditions to the partners, so the period change bound applies at every label.

% before: is-onto-step-exp-bound-link-depth
The linkage of a constituent weight to the weights of the dictionary bounds its depth on the first
$k-1+j_\tau$ rows. This is the bound the weight inputs require.

% before: is-onto-step-exp-bound-link
A Hida setting linked to the weights has the weight inputs, with the same constituent data. Only the
depth bound has to be checked.

% before: is-onto-step-exp-bound-short-branching
The inclusion of forms rests on a bound for the power of $l$ that carries a weight space of a dual Weyl
module into the Weyl module, when a block of $p$ equal entries $w$ is added to the weight. The power
grows by the valuation of an explicit product $D_n$ over contents of shapes. It is used in a form graded
by a depth, with the block placed last.

% before: is-onto-step-exp-bound-short-graded-block
The graded form is proved for representations of $\mathfrak{gl}$ of $I_1\sqcup I_2$ with a subspace on
which the operators within $I_1$ act by the scalar $w$ and those raising $I_2$ into $I_1$ vanish. The
lattices are spanned by divided powers of the lowering operators applied to a lattice in this subspace.
Tensors with values in the subspace carry the Casimir operators of partial sums, whose joint eigenspaces
are indexed by ways of adding boxes to $\lambda_2$ one at a time.

% before: is-onto-step-exp-bound-short-graded-skeleton-jucys
The Casimir operators of partial sums are compared through the Jucys--Murphy operators of the tensor
positions. These combine the action on the values with the transpositions of the $I_2$-entries, and they
commute with permutations of the $I_1$-entries.

% before: is-onto-step-exp-bound-short-graded-skeleton-jucys-relations
The relations below are those of the degenerate affine Hecke algebra. The step from $C_j$ to $C_{j+1}$ adds
twice the Jucys--Murphy operator of position $j$, and adding a box of content $c$ raises the Casimir value
by $2c+q$.

% before: is-onto-step-exp-bound-short-graded-skeleton-path-eigen
Along a box path, each Jucys--Murphy operator therefore acts by the content of the box added at its step.
The polynomial $A(X)$ then acts by the product of $X$ minus these contents.

% before: is-onto-step-exp-bound-short-graded-skeleton-commute
When the operators $E_{bb'}$ satisfy the relations of $\mathfrak{gl}_q$, the Jucys--Murphy operators commute, and
their symmetric functions commute with the permutations of the positions. The coefficients of $A(X)$ then
descend to commuting operators on the symmetrized tensors.

% before: is-onto-step-exp-bound-short-graded-skeleton-maps
A tensor in $T_n$ is turned into a vector of the block module by applying $n$ lowering operators to its values.
In the other direction, raising operators read off a vector as a function on multisets, with values in $V^+$.

% before: is-onto-step-exp-bound-short-graded-skeleton-mu
The map $\mu_n$ intertwines the action of the block module with explicit operators on tensors. The sum of the
diagonal operators on $I_1$ acts on its image by $pw-n$, which separates the images for different $n$. A raising
operator moves the image down by one step, and on symmetric tensors it acts through the last Jucys--Murphy
operator.

% before: is-onto-step-exp-bound-short-graded-skeleton-iota
Iterating the raising step shows that reading off $\mu_n(z)$ by raising operators gives the symmetrization of
$A(w)z$. This is how the polynomial $A$ in the Jucys--Murphy operators enters the denominators.

% before: is-onto-step-exp-bound-short-graded-skeleton-quotient
The denominator is cleared by dividing a polynomial $D$ by the characteristic polynomial of commuting operators
in a generic ring. The remainder vanishes on every joint eigenspace, and the quotient stays in the algebra
generated by the operators.

% before: is-onto-step-exp-bound-short-graded-skeleton-weights
The Jucys--Murphy operators preserve the lattices and the weights carried by the values of a tensor. The
weight of a value is shifted by the multiset of its position, and the shift is dominant for $\xi$.

% before: is-onto-step-exp-bound-short-graded-skeleton-assembly
These pieces give the bound up to the factor $D(w)$. A dual vector of depth $n$ is read off as a function on
multisets with values in $\nabla'$; the operator from the division lemma moves it into $\Delta_L$ after
multiplying by $l^e$, and the injectivity of reading off returns it to $\Delta_n(\Delta_L)$.

% before: is-onto-step-exp-bound-short-graded-skeleton-form
A contravariant form makes reading off adjoint to the lowering operators. It identifies the dual of the
lattice in depth $n$, and with generation and nondegeneracy it makes reading off injective on the image of
$\mu_n$.

% before: is-onto-step-exp-bound-short-graded-skeleton-shapes
On a box path space, $A(X)$ acts by the product of $X$ minus the contents of the added boxes. These products
divide the polynomial $D$ of the depth bound when no added box has content $w$, and they vanish at $w$ otherwise.

% before: is-onto-step-exp-bound-short-graded-skeleton
In a block module whose tensors decompose into box path spaces, and whose symmetrized tensors vanish on
shapes with a column of more than $p$ added boxes, the denominators of a dual vector of depth $n$ are
bounded. The bound is the exponent of the lattice in $V^+$ plus the valuation of $D_n$.

% before: is-onto-step-exp-bound-short-graded-rational
The dual Weyl module of $w^p\lambda_2$ over $\mathbf Q$ is a block module, with the first $p$ rows as
$I_1$ and the last $q$ rows as $I_2$. Its subspace $W_0$ consists of the vectors of weight $w$ on the first
$p$ rows that are killed by the operators raising the last $q$ rows into the first $p$. The lattice in
$W_0$ is generated from $h_\nu$ by the divided powers on the last $q$ rows.

% before: is-onto-step-exp-bound-short-graded-rational-module-closed
Generation of the dual Weyl module rests on its irreducibility over $\mathbf Q$. A nonzero subspace
stable under the divided powers contains a vector killed by the raising operators, which is a multiple of
$h_\nu$. Its orthogonal complement for the Fock pairing is stable as well, and positivity makes it $0$.

% before: is-onto-step-exp-bound-short-graded-rational-module
The structural hypotheses of the bound hold for this module. Contravariance and positivity come from the
Fock pairing, and generation from the vector $h_\nu$, which lies in $W_0$.

% before: is-onto-step-exp-bound-short-graded-rational-lattice
The lattice $\Delta(\Delta_L)$ lies in the Weyl module over $\mathbf Z_{(l)}$, as divided powers preserve the
Weyl module. By the duality of the Weyl and dual Weyl modules under the Fock pairing, the dual Weyl module
over $\mathbf Z$ pairs integrally with it.

% before: is-onto-step-exp-bound-short-graded-rational-pieri-irreducible
Both the decomposition into box path spaces and the recovery of a vector from its last $q$ rows use that $W_0$
is irreducible for the operators on these rows. A nonzero stable subspace of $W_0$ generates the whole module,
and the depth decomposition then places $h_\nu$ in it.

% before: is-onto-step-exp-bound-short-graded-rational-pieri-highest
The Pieri rule is proved for representations of $\mathfrak{gl}_q$ with a positive definite contravariant form,
in which highest weight vectors control the stable subspaces. Their weights are dominant, and the Casimir
operator acts on them by $\mathrm{cas}$.

% before: is-onto-step-exp-bound-short-graded-rational-pieri-tower
Adding one tensor position adds a copy of the standard representation, whose matrix units commute with the
earlier operators. A highest weight vector of the larger action projects to one of the smaller action whose weight
has one box fewer, and the Casimir values then record a box path.

% before: is-onto-step-exp-bound-short-graded-rational-pieri-general
The tensors $T_n$ form such a tower, with the Casimir operators $C_j$ of the partial sums. When the module is
generated by highest weight vectors of one weight, the tower lemma decomposes $T_n$ into box path spaces.

% before: is-onto-step-exp-bound-short-graded-rational-pieri-real
For the rational block module these hypotheses come from the Fock pairing, the diagonal action on monomials and
the irreducibility of $W_0$. The highest weight vector of weight $\lambda_2$ is $h_\nu$, which lies in every nonzero
stable subspace.

% before: is-onto-step-exp-bound-short-graded-rational-pieri
The decomposition into box path spaces is Pieri's rule for the module $W_0$ of the last $q$ rows. Each
tensor factor adds one box to the highest weight, and the partial Casimir operators record the shapes
along the way.

% before: is-onto-step-exp-bound-short-graded-rational-columns-paths
The vanishing is proved by reordering the steps of a box path. Two consecutive boxes in decreasing
columns can be added in the other order. Two consecutive boxes in the same column cannot, and their
contents differ by $1$.

% before: is-onto-step-exp-bound-short-graded-rational-columns-weights
Exchanging two steps corresponds to the operator $s_j$, which moves a vector of one weight to the
reordered weight. When the reordered sequence is not a content sequence, $s_j$ acts by $-1$, and a vector
antisymmetric in more than $p$ consecutive positions symmetrizes to $0$, as two of these positions carry the
same element of $I_1$.

% before: is-onto-step-exp-bound-short-graded-rational-columns-general
Moving boxes to later steps when their columns decrease ends with the columns nondecreasing. A column with
more than $p$ added boxes then appears in more than $p$ consecutive steps.

% before: is-onto-step-exp-bound-short-graded-rational-columns
Symmetrization ties each tensor position to one of the $p$ elements of $I_1$. A box path space whose final
shape has a column of more than $p$ added boxes contributes nothing after symmetrization.

% before: is-onto-step-exp-bound-short-graded-rational-exponent-restriction
Setting the first $p$ rows and columns to those of the identity matrix restricts polynomials in a
$(p+q)\times(p+q)$ matrix to polynomials in a $q\times q$ matrix. It intertwines the divided powers on the last
$q$ rows, carries the dual Weyl module of $\nu$ to that of $\lambda_2$, and is injective on $W_0$.

% before: is-onto-step-exp-bound-short-graded-rational-exponent-dual
Integrality of the values at the identity matrix of all the images of $g$ under the divided powers controls the
coefficients of $g$ at the lower triangular monomials. These coefficients determine whether $g$ lies in the
dual Weyl module over the integers.

% before: is-onto-step-exp-bound-short-graded-rational-exponent
The exponent of the lattice in $W_0$ comes from the hypothesis on the pairs $(\lambda_2,\mu_2)$. A vector of
$W_0$ is determined by its restriction to the last $q$ rows and columns, which lies in the dual Weyl
module of $\lambda_2$.

% before: is-onto-step-exp-bound-short-graded
The block exponent bound in graded form is a statement about Weyl modules over $\mathbf Z$ alone. It
does not involve the Hida setting.

% before: is-onto-step-exp-bound-short-last-dual
The last graded form follows from the graded form by replacing each weight with entries at most $M$ by its
complement $M-\kappa$, read backwards. This turns a block placed after the weight into a block placed
before it, and keeps the depths.

% before: is-onto-step-exp-bound-short-last-weights
The complement needs the row degrees to be at most $M$. A nonzero vector of the dual Weyl module of a
weight bounded by $M$ has row degrees bounded by $M$, because the module is spanned by products of at
most $\nu_0$ minors.

% before: is-onto-step-exp-bound-short-last-cofactor
The complement is realized on polynomials by substituting for each variable an entry of the adjugate.
Up to powers of the determinant this substitution is an involution, and it carries row degree $\mu$ to
$\mu^\vee$.

% before: is-onto-step-exp-bound-short-last-nabla
The substitution carries the dual Weyl module of $\nu$ into that of $\nu^\vee$, and the Weyl module into
the Weyl module, after a power of the determinant. For minors this is Jacobi's theorem on the minors of
the adjugate.

% before: is-onto-step-exp-bound-short-last-duality
So the exponent of a pair of weights equals the exponent of the complementary pair. The powers of the
determinant cancel because the polynomial ring is a domain.

% before: is-onto-step-exp-bound-short-last
The top weight theorem uses the bound with the added block placed after the given weight. In that form
the product $D_n$ is taken for $M-w$ and for the complementary weight $\lambda_2^\vee$.

% before: is-onto-step-exp-bound-short-top-data
The inclusion of forms holds for Hida settings whose weights have a given shape, with a bound
$C'_\tau$ built from the depths and the exponents. The top inputs list what the proof uses of the
setting and of its constituents.

% before: is-onto-step-exp-bound-short-top-window-dominance
The exponent at an index of $T_2$ is bounded by adding rows to the weight, each time as a block of equal
entries placed last. The last graded form applies to a weight space once the depth of its block is known.
That depth is read off from the row degrees of the first rows, which are dominated by the weight.

% before: is-onto-step-exp-bound-short-top-window-products
The products $D_n$ are built from the contents of skew shapes with at most $p$ boxes in each column. So
their exponents are at most $p$, and their factors are positive integers up to $w+q-m-1$. Legendre's
formula then bounds their valuations.

% before: is-onto-step-exp-bound-short-top-window-peel
Rows whose entries need not agree are added one at a time, each as a block of one row. When the entries of
the weight spread by at most $h$, each such row raises the exponent by at most $(h-1)/(l-1)$.

% before: is-onto-step-exp-bound-short-top-window-block
The $t$ equal entries are added at once as one block. The exponent then grows by the largest valuation of
the products $D_n$ over the depths that occur.

% before: is-onto-step-exp-bound-short-top-window-steps
The weight at an index of $T_2$ is built in three stages: the $j$ rows after the first $k$, then the block
of $t$ equal rows, and then the last $q$ rows. Only the block stage depends on the depth on the first
$k+j$ rows.

% before: is-onto-step-exp-bound-short-top-window-label-sum-base
The valuation of $D_n$ is a sum of valuations of labels, weighted by the number of boxes with each label.
As $h<l^2$, a label $y$ contributes $v_l(y-\beta)\le1$, which is $0$ unless $y$ is congruent to $\beta$
modulo $l$. Labels far below $x^-$ carry few boxes, which gives the logarithmic term.

% before: is-onto-step-exp-bound-short-top-window-label-sum
The bound is needed for every integer $\beta$, with a logarithmic term that vanishes when $h<l$. For
$h<l$ no label is divisible by $l$. Otherwise a translation of the labels moves $\beta$ to $0$ and
changes neither side.

% before: is-onto-step-exp-bound-short-top-window-labels
For the block stage, the number of boxes with a given label is bounded through the positions of the first
$k+j$ rows. The rows after the first $k$ carry at most $j$ boxes of each label below $x^-$. A box in one of
the first $k$ rows with a label far below $x^-$ needs many boxes in its row, and the shape has only $i$
boxes in all.

% before: is-onto-step-exp-bound-short-top-window
At an index of $T_2$ the theorem needs the exponent of each weight of small depth. The last graded form of
the block exponent bound gives it, with the valuations of the products $D_n$ bounded by $C'_\tau$.

% before: is-onto-step-exp-bound-short-top-base
The exponents are defined over $\mathbf Z$, and the forms take values over $O$. Both modules over $O$ are
spanned by bideterminants with integer coefficients, so an inclusion over $\mathbf Z$ passes to $O$ once
the integers prime to $l$ are units.

% before: is-onto-step-exp-bound-short-top-forms
Forms are moved into a smaller submodule one simple subquotient at a time. A simple subquotient below a
power of $\varpi$ already occurs modulo $\varpi$, and a subquotient whose forms vanish after localization
can be crossed by an element of $S$.

% before: is-onto-step-exp-bound-short-top-core
The tensor factors are treated one at a time. At an index of $T_2$ the exponent $c'_\tau$ suffices,
because a simple subquotient met along the way has a weight too deep to carry forms that survive
localization.

% before: is-onto-step-exp-bound-short-top
The top weight theorem bounds the exponent of the inclusion of forms by the bound $C'_\tau$ at the
indices of $T_2$ and by the Weyl exponents elsewhere. It holds for any data with the top inputs.

% before: is-onto-step-exp-bound-short-top-bound
The bound $C'_\tau$ is at most the bound $C^{\mathrm s}_\tau$ of the same shape as the short bound. An
inclusion with a smaller exponent gives one with a larger exponent.

% before: is-onto-step-exp-bound-short-shape
For the weights of the dictionary, the thick indices have the shape that the top weight theorem
requires. The conditions come from (W1) to (W5) and from the bound on the weights beyond the first
$k-1$ rows.

% before: is-onto-step-exp-bound-short-inputs
The weight inputs of a Hida setting, together with the numerical conditions on the parameters, give
its top inputs. The Weyl exponents of the dictionary are the least exponents of the inclusions of the
Weyl lattices.

% before: is-onto-step-exp-bound-short
From the weight inputs, the forms of a Hida setting with values in the dual Weyl modules are carried
into those with values in the Weyl modules by a power of $l$. The exponent is the short weight bound.

% before: is-onto-step-exp-bound-coarse
The coarse weight bound exceeds the short one by six times $q_\tau+2j_\tau$ at each thick element. An
inclusion with a smaller exponent gives one with a larger exponent.

% before: is-onto-step-exp-bound-models
At a twisted place in the range of the weights, an exponent model bounds the block lattices by the
coarse weight bound. The Frobenius data of the model do not change this bound.

% before: is-onto-step-exp-bound-place
At one place, the inclusions of the block lattices and the exponent models together bound the lift
exponent by the recursive exponent. Which of the two is used depends on whether the residue
characteristic is in the range of the weights.

% before: is-onto-step-exp-bound
The exponent bound of the record object follows from the bound at one place and the bound on the sum,
applied to the weights of the dictionary at every lift. Its exponents $J'(j,l)$ are the recursive
exponents of the weights.

% before: is-onto-estimates
The surjectivity bound rests on estimates at the twisted places, each proved from a different part
of the published theory. They are stated for all record objects of the points at large scales, with
constants that may depend on the family.

% before: is-places-error-terms
The three conditions of the surjectivity bound are proved with explicit error terms. At a twisted
place the error term is a number attached to its residue characteristic, chosen according to the
kind of prime it is, plus a share of the slack in the index bound of the record object.

% before: is-places-prime-sums
The error terms at the good primes up to $Y$ enter the third condition through sums weighted by
$\log l$. Three explicit estimates for sums over primes bound them.

% before: is-places-prime-sums-rs
The proof of Lemma~\ref{is-places-prime-sums} works with elementary bounds: binomial coefficients,
Legendre's formula and a finite check. The sharper explicit estimates of Rosser and Schoenfeld for
$\theta(t)$ give the same three bounds in a few lines.

% before: is-places-error-sum
Grouping the twisted places by residue characteristic turns sums over places into sums over primes.
The good primes then contribute an explicit main term, and the bad primes a term that the page
estimates control.

% before: is-places-absorb
The bounds for the error sums are a list of terms of different shapes. One numerical inequality
absorbs all of them into a multiple of $x^\varepsilon fmb\log b$, and this multiple produces the
constant $C_E$.

% before: is-places-index-subspace
The local index bound of Lemma~\ref{is-places-local-index} assumes a bound for the lattices
$M_s\cap W$ and $\Theta_1\cap W$ for every subspace $W$, and its proof needs a bound for the single
lattice $M_s\cap a^{-1}\Theta$. The two kinds of bound are equivalent. After localizing at $w$, the
elementary divisors of the smaller lattice in the larger one single out the subspace for which the
first bound gives the second.

% before: is-places-index-projector
The proof of Lemma~\ref{is-places-local-index} passes from lattices in $Z$ to their
$\psi$-eigenvectors under $U$. A sum over $U$ weighted by the values of $\psi$ maps $Z$ onto these
eigenvectors and multiplies each of them by $|U|$, so it is a projection up to the factor $|U|$. The
last condition on a Mackey pair makes the sum vanish on the translates of eigenvectors by elements
outside $U$.

% before: is-places-index-eigen
Step (c) of the proof of Lemma~\ref{is-places-local-index} bounds the index of the
$\psi$-eigenvectors in two $\rho(P)$-stable lattices by the index of the lattices themselves. The
translates of $V_\psi$ by representatives of the cosets of $U$ are independent. Up to the factor
$|U|$ of the weighted sum of Lemma~\ref{is-places-index-projector}, the index on $Z$ therefore
contains the index on $V_\psi$ once for each coset.

% before: is-places-local-index
The first condition of the surjectivity bound bounds the index of each image lattice at a twisted
place and at its conjugate. That index is computed locally, from the lattices of the record object
and the $\psi$-eigenvectors of the subgroup $U$.

% before: is-places-combination
With the local index bounds and the numerical lemmas, the estimates of the record object combine
place by place into the three conditions. The hypotheses are the estimates of
Lemma~\ref{is-onto-estimates} together with the ties of the record object to the realization.

% before: is-places-level-index
For the record objects of the points, the group is a product of mirabolic groups over finite fields
and $U$ is the product of their unipotent subgroups. Their orders give the index bound and the
valuations of $|U|$ that the combination needs.

% before: is-ties-mirabolic-mackey
For the record object of a point, $P$ is a product of mirabolic groups over finite fields and $\psi$
is, on the unipotent subgroup, a product of nontrivial additive characters of the sums of the entries
just above the diagonal. The local index bound needs such a pair to satisfy the Mackey condition. An
element of $P$ outside $U$ conjugates a suitable unipotent element to another one with a different
value of $\psi$.

% before: is-ties-class-sums
The record object of a point is tied to the realization through the rational structures of the lifts
and the lift forms that transform under the face by its character. A class sum writes the theta lift
of such a form in the rational coordinates of the lift, and the twisted lattices of the record object
are built from the theta series of the tuples.

% before: is-ties-lattice-systems
The theta lattice of a tuple is spanned by theta series at all points of the finite adelic group,
with coefficient vectors integral for the lattices of these points, while the integral level forms
are integral only at representatives of the double cosets of the level. A lattice system attaches a
lattice of coefficients to every point, compatibly with the rational points and the level, and agrees
with the integral vectors at the representatives. A family tie identifies the rational coefficients
with a space of polynomials and so provides one.

% before: is-ties-lattice-facts
The theta lift of a lift form unfolds into a sum over the representatives of theta series, and the
form averaged over one stabilizer lifts to a single theta series. Along a lattice system, a theta
series at any point with an integral coefficient vector is, up to a root of unity, the theta lift of
an integral lift form. Tracing over the face with its character then gives an integral level form.

% before: is-ties-class-sum-facts
Class sums commute with the good idempotent and take values on which the face acts by its
character. Their values on integral level forms and the theta lattice of the tuple span the same
space, and on the good part they agree with the theta matrices of the rows by the base links.

% before: is-ties-level-tie
Being recorded ties every space, map and lattice of the record object to the realization. The part
of the tie used below consists of isomorphisms of coordinates that carry the group action, the
idempotents, the lattices and the maps $\theta'$ of the record object to those of the realization.

% before: is-twisted-places-chain
The estimates combine into the three conditions of the surjectivity bound with some error terms.
The combination is an inequality chain at each twisted place over the lattices of the record object.

% before: is-eigenvalue-lift-data
The lower norm bound rests on an explicit formula for the eigenvalues of a theta lift composed with its
adjoint. The formula is needed at every embedding of the coefficient field, so its data are a Weil
datum, a level vector and coefficients for the lift, with characters conjugate to those of the
realization.

% before: is-eigenvalue-zeta-inputs
The eigenvalues are computed from the Rallis inner product formula, which writes the norm of a theta
lift as a doubling integral. That integral splits into a finite part, an Euler product given by the
unramified computation, and an archimedean part, bounded through the formal degree of the discrete
series.

% before: is-eigenvalue-local-factors
The other factors of the doubling integral are local computations at the infinite places and at the
places of the level. The volume of the definite quotient is the index of the level times the mass of
the standard lattice, which is bounded below by counting the automorphisms of that lattice.

% before: is-eigenvalue-lift-computations
Three identities connect these inputs to the adjoint of a theta lift. The adjoint is a sum of theta
integrals over a basis of the coefficients, the archimedean integral factors through the matrix
coefficients of the lowest $K$-types, and an operator whose inner products are a fixed multiple of
those of the space acts by that multiple.

% before: is-eigenvalue-size-bounds
Each factor of the eigenvalue is then estimated in terms of the ranks and the sizes of the lifts. The
estimates for the roots of unity, the partial $L$-values and the normalizing factors depend only on the
definite-side datum.

% before: is-eigenvalue-numeric
The estimates combine into one inequality, which loses at most $8fn^2L'$ against the hook-length term.
It is stated for abstract numbers, since it is applied at every embedding of the coefficient field.

% before: is-eigenvalue-numeric-direct
The proof of Lemma~\ref{is-eigenvalue-numeric} first bounds the mass and the factors at the
embeddings together, after moving the roots of unity into the mass. Bounding each factor separately
gives the same inequality with a shorter count.

% before: is-eigenvalue-adelic
With these inputs the eigenvalue of a theta lift on each Hecke eigenspace is an explicit positive
number. Its logarithm is the hook-length sum of the weights at the embedding, up to a loss controlled
by the sizes of the lifts.

% before: is-eigenvalue-type-parts
The eigenvalue formula holds on each Hecke eigenspace separately, so a bound on the whole space needs
the eigenspaces to span it. The same comparison of local components identifies the conjugate characters
at other embeddings with those of lift data.

% before: is-eigenvalue-projection
The lifts of the realization are theta lifts followed by the projection given by a central idempotent
of the group. A lower bound for a map on spanning eigenspaces of its square passes to its composite
with an orthogonal projection, on the orthogonal complement of the kernel.

% before: is-eigenvalue-conjugate-datum
At an embedding of the coefficient field other than the base one, the lift of the realization is
compared with a theta lift built from Galois conjugates of the data of the realization. A conjugate
datum consists of these conjugate characters, coefficients, forms and lift.

% before: is-baselinks-rational-tests
The rationality tests of the primitive rows are functionals that take rational values on theta lifts
of rational forms. Coordinates that such functionals see as rational are themselves rational, and
the image of the adjoint of such a lift is then spanned by images of rational forms.

% before: is-baselinks-group-algebra
The definite side of a face uses the image, on the lift forms, of an idempotent of the group algebra
of $G$, which acts on the rational forms and, through the semilinear maps of the realization, on
their completions. Theta lifts that commute with $G$ carry these idempotents along, so statements
about such images pass between the rational forms, their completions and the target of the lift.

% before: is-baselinks-lift-cut
The first base link asks that the theta lift of the coordinate vectors of the lift forms be a
bijection onto the image of the coordinate vectors of the rows. This follows by linear algebra once
the coordinate vectors span the image of the adjoint inside a subspace that the lift and its adjoint
exchange with the image of the rows, and the lift multiplies norms by a positive scalar on each part.

% before: is-baselinks-lifted-parts
The definite subspace of the lift forms is the orthogonal complement of the kernel of the theta lift.
Inside a subspace stable under the Hecke operators it is the sum of the Hecke eigenspaces whose
systems also occur in the image of the rows, provided the adjoint of the lift reaches each of those
eigenspaces.

% before: is-baselinks-definite-coordinates
The definite side of a face is built from a basis of the definite subspace inside the level part,
with the class-sum pairing as its Hermitian matrix. When that subspace is cut out by Hecke
eigenvalues of operators with rational matrices, it has a basis over $K_0$, and the metric is then
the class-sum pairing at every infinite place.

% before: is-eigenvalue-core-group
The rational lifts are compared through the functions that the elements of the Hecke modules
$H_W(i,\sigma)$ give at a single vector of the models. The finite group $G$ acts on these functions by
right translation by elements of a compact level group, so it acts unitarily, and the inner product
can be read off from the functions at any tensor of unit vectors.

% before: is-eigenvalue-conjugate-forms
Two conjugate data at one embedding are compared through invariant pairings of their theta functions.
A pairing invariant under an irreducible family of isometries is a multiple of the inner product, and
on tensors such pairings are multiples of products of inner products. Two maps whose pairings all have
this form, and which agree up to a unit at one point, then agree up to one unit everywhere.

% before: is-eigenvalue-conjugate-invariant-form
The pairing used for theta functions is their $L^2$ pairing on a fundamental domain for the rational
points. A Weil datum preserves it, so a scalar by which a Weil datum acts on a theta function has
absolute value $1$.

% before: is-eigenvalue-conjugate-coefficients
The coefficient data of two conjugate data are identified with the same rational coefficient space by
semilinear maps with the same pairings. The map between them that matches these identifications is
unitary, and weak approximation at the archimedean places makes it intertwine the adelic actions.

% before: is-eigenvalue-conjugate-k-types
The comparison of two conjugate data takes place on the lowest $K$-types of the models of $W_i$. The
pairings of theta functions indexed by these spaces are multiples of the inner product when the
lowest $K$-type is a single irreducible representation of the maximal compact subgroup and consists of
smooth vectors, which is the condition named here.

% before: is-eigenvalue-conjugate-theta-closed
The vectors of a conjugate datum are obtained from theta functions by Heisenberg translations and by
derivatives along the archimedean directions. These operations preserve theta functions, and the
difference quotients of a theta function along a real direction converge to its derivative uniformly,
at a rate linear in the step.

% before: is-eigenvalue-conjugate-theta-schur
Two Weil data with the same characters differ on theta functions by operators that commute with the
Heisenberg translations and preserve norms. Such an operator is a unitary scalar, because the
products of two theta functions are periodic functions whose Fourier coefficients it preserves.

% before: is-eigenvalue-conjugate-vacuum-exists
The comparisons below start from a vacuum vector in the Fock model. The classical theta series of the
Gaussian provides one for any pair of self-dual lattices, and normalizing its norm gives a Fock vacuum
vector.

% before: is-eigenvalue-conjugate-split-form
At a split place the Weil datum acts on linear forms invariant under one half of $\mathbb W$ through a
positive scalar times the characters. To compare two Weil data there, one needs such a form that does not
vanish on all theta functions. It comes from the Fock vacuum, which does not lie in the span of the
differences that the invariance kills.

% before: is-eigenvalue-conjugate-split-ideals
The decomposition of an adelic unitary element below moves its denominators to places that split in
the host. Such denominators are cleared by elements of the ring of integers with prescribed prime
divisors. These exist because the class group is finite: a power of an ideal is principal, and its
generator has the prime divisors of the ideal.

% before: is-eigenvalue-conjugate-split-norm-one
A quasi-reflection with a norm-one eigenvalue $u$ is a rational unitary element whose denominators may
involve primes fixed by the conjugation. Replacing $u$ by a suitable element $u'$ of norm one moves these
denominators to split primes and to primes dividing $N$, while $u'/u$ stays congruent to $1$ at the
primes dividing $N$. The element $u'$ is written down from a presentation $u=\alpha/\bar\alpha$ given by
Hilbert's Theorem 90.

% before: is-eigenvalue-conjugate-split-correction
Every rational unitary element is a product of quasi-reflections. Correcting the eigenvalue of each
factor gives a rational element congruent to the given one at the primes dividing $N$, whose
denominators lie at split primes and at primes dividing $N$.

% before: is-eigenvalue-conjugate-split-approximation
An adelic unitary element can be brought close to $1$ at the primes dividing $N$ by a rational one. This
is weak approximation for the unitary group at those places, with closeness measured by the lattice
$\hat\Lambda$.

% before: is-eigenvalue-conjugate-split-finite
Over the finite adeles the decomposition is built in two steps. A rational element brings the given one
close to $1$ at the primes dividing $N$ and leaves its remaining denominators at split primes; the
idempotent of the places under those primes then separates a factor supported at split places from a
factor congruent to $1$ modulo $N$.

% before: is-eigenvalue-conjugate-split-decomposition
A continuous character of the adelic unitary group is controlled by its values on a few kinds of
elements. Every element factors into a rational one, an archimedean one, one supported at split places
and one close to $1$ at the finite places.

% before: is-eigenvalue-conjugate-split-character
The ratio of two Weil data with the same characters is a continuous unitary character of the adelic
unitary groups. It is trivial on the rational points and at the split places, and these two conditions
force a unitary character to be trivial.

% before: is-eigenvalue-conjugate-vacuum-classical
In real coordinates a Fock vacuum vector becomes a classical theta function for a principal
polarization. The space of such functions with a given characteristic is a line, which gives the
uniqueness of the vacuum.

% before: is-eigenvalue-conjugate-vacuum-coordinates
The archimedean part of $\mathbb W$ carries a complex structure and a symplectic form coming from the
Cartan involutions and the standard character. In these coordinates the annihilation operators become
the Cauchy-Riemann equations twisted by a Gaussian factor.

% before: is-eigenvalue-conjugate-vacuum-unique
A vacuum vector is determined by its values at the archimedean points, and those values satisfy the
equations of a classical theta function. The rational points whose finite parts lie in the lattice of
the pair form a self-dual lattice for the symplectic form, so the classical uniqueness applies.

% before: is-eigenvalue-conjugate-vacuum-span
The unitary groups at the archimedean places act on Fock vacua through scalars, and this has to be
transferred to every theta function killed by the annihilation operators. Such a function is fixed by a
lattice of finite index in the lattice of the pair, and a Fourier argument on the finite quotient writes
it through translates of the vacuum.

% before: is-eigenvalue-conjugate-theta-facts
The comparison also uses properties of the Fock model and of Weil data that do not involve conjugate
data. Vacuum vectors and Weil data are determined up to a unit by their defining conditions, the
creation operators move with the bases under the Weil representation, and the lowest $K$-types at the
embeddings that agree with $\sigma_0$ on $K$ are irreducible.

% before: is-eigenvalue-conjugate-level-vectors
The level vector of a conjugate datum is obtained from a Fock vacuum vector by a construction at the
places of $P_g$ that is linear in the vacuum. Since the vacuum is unique up to a unit, so is the level
vector.

% before: is-eigenvalue-conjugate-pure-tensors
The theta integral of a conjugate datum is built from the vectors $\phi\otimes_B\mathrm{io}(z\otimes w)$.
For two choices of the Weil datum, the bases and the map $\mathrm{io}$ with the same characters, these
vectors agree up to a unit, because both families are equivariant for the same Weil representation
and both maps $\mathrm{io}$ take the value $\Delta$ at unit vectors.

% before: is-eigenvalue-conjugate-good
The weight tie describes the core forms at an embedding of $L$ in two ways, according to the square
root of $-d$ to which the embedding sends that of the core. At a good embedding every label carries
the embedding of the host attached to it to its own type, which makes the models holomorphic there.
An embedding is good when it agrees with $\sigma_0$ on $K$, and exactly one of an embedding and its
complex conjugate is good.

% before: is-eigenvalue-conjugate-models
At an embedding that is not good, the functions of the core forms are obtained from those at the
conjugate embedding, which is good. The models of
discrete series at the two embeddings are related by conjugate-linear isometries, and these carry
multiplicity maps to multiplicity maps with the conjugate level character.

% before: is-eigenvalue-conjugate-lowest-types
A theta integral and an element of a Hecke module are both multilinear in vectors of the lowest
$K$-types, and both transform under the compact groups at the archimedean places. Because each lowest
$K$-type is irreducible, two such maps that agree at one tensor of nonzero vectors agree everywhere.
This turns an agreement at the extremal vector into an agreement of the whole families of functions.

% before: is-eigenvalue-conjugate-class-sums
The rational lifts and their conjugates are computed from sums of core forms, so every lift form must
be written as such a sum. The pins of the one-map data read a lift form at representatives of the
classes, and the one-map tie expresses the section built from these values through core forms, with
coefficients in the coefficient field of the realization.

% before: is-eigenvalue-conjugate-dictionary
The core describes its forms by sections on flag domains, one for each component of a uniformization,
and the Hecke modules of the realization consist of adelic functions. A dictionary passes from the
first to the second: it attaches to each adelic point a component and a frame, and to a family of
sections its adelization. Its good properties are the ones under which the adelization of a cusp form
is a holomorphic cusp form on the adelic unitary group.

% before: is-eigenvalue-conjugate-dictionary-facts
The proofs below use a few computations with the archimedean points and with the transvection
subgroups. The archimedean points split off the adelic group, the frames of a holomorphic model are
orthonormal, and a transvection subgroup is a copy of the adeles whose rational points are cocompact.
Along an isotropic vector, the flow of a flag domain is a one-parameter group of unimodular matrices.

% before: is-eigenvalue-conjugate-spectral-facts
To produce a multiplicity map from a single holomorphic form, the form is placed in the space of
square-integrable functions on the adelic quotient and the archimedean group is made to act. This uses
the regular representation, the rigidity of holomorphic functions on the archimedean group with a given
torus law, and standard facts on smooth vectors and on the centre of the enveloping algebra.

% before: is-eigenvalue-conjugate-dictionary-good
At a good embedding the weight tie imposes laws on the dictionary of the core. They give each condition
of a good dictionary, with the level of the trivial level datum of $W_i$ as the level group, on which
the multiplier and the level character are trivial.

% before: is-eigenvalue-conjugate-adelization
The adelization of bounded sections through a good dictionary inherits the invariance of the sections
under the arithmetic groups, and the base frame turns the holomorphy of the sections into holomorphy
along the archimedean points. For cusp forms, an average over a transvection subgroup reduces to the
constant term of a section along an isotropic vector of its lattice.

% before: is-eigenvalue-conjugate-matrix-coefficient
A holomorphic form generates under the archimedean group a representation whose matrix coefficient at
the form is that of an extremal vector. Both coefficients are holomorphic functions on the archimedean
group with the same torus law, and such a function is determined by its value at $1$.

% before: is-eigenvalue-conjugate-multiplicity
The matrix coefficient gives a map from the model of discrete series to the square-integrable functions
that carries the extremal vector to the form. Its values are smooth vectors, whose continuous
representatives form the required multiplicity map.

% before: is-eigenvalue-conjugate-characters
A conjugate datum needs conjugates of the characters $\mu^0_i$ and $\eta$ and a Weil datum for them.
The conjugate of $\eta$ is obtained by twisting $\eta$ with a character of the class group, which the
product formula and the agreement of $\sigma$ with $\sigma_0$ on $K$ make well defined. The Weil datum is
a twist of the one of the realization, allowed because the twisting characters are trivial on the
ideles of $F_h$.

% before: is-eigenvalue-conjugate-coefficient-forms
The coefficients and the forms of a conjugate datum are those of the realization with $\sigma$ applied
to their rational structure. The coefficient space is realized with the Gram matrices at the embeddings
through which $\sigma$ factors, and the form space is obtained from the class presentation of the forms.

% before: is-eigenvalue-conjugate-level-data
The theta integral of a conjugate datum is built from a level vector, bases at the archimedean places and
a map from the lowest $K$-types and the coefficients to polynomials. These come from the vacuum vectors
and the Fock model, and the integral converges because $V_i$ is positive definite.

% before: is-eigenvalue-conjugate-phase
The comparison with the one-map data produces a function of absolute value one on the adelic group of
$W_i$, invariant under the rational points, the archimedean points and the level. Multiplication by such
a function preserves multiplicity maps and their Petersson pairings.

% before: is-eigenvalue-conjugate-extension
The one-map data at a lift have their own coefficient field, which the pins identify with $L$. An
embedding of $L$ therefore gives an embedding of this field, and the two agree on the coefficient field
of the core, so that both see the same dictionary of the core.

% before: is-eigenvalue-conjugate-phase-function
The constants that the comparison attaches to the core live on the components of its dictionary. A
function on the adelic group of $W_i$ that depends only on the component inherits the invariances of a
good dictionary, which are the invariances a phase needs.

% before: is-eigenvalue-conjugate-class-values
Class sums are read through the coefficient map of the one-map data, and a conjugate datum through the
coefficient map of a family tie. When the two maps agree after base change, a value of a lift form at a
representative determines a rational vector over $L$, which the conjugate datum carries to its own
value.

% before: is-eigenvalue-conjugate-class-action
The central element of the comparison acts on sections of the core over the field of the one-map data,
while the core forms are defined over the field of the core. The action moves the level structures of
framed points, so it commutes with extension of scalars.

% before: is-eigenvalue-conjugate-unfolding
The theta integral of a conjugate datum is an integral over a fundamental domain of the definite group
$U(V_i)$. When the integrand is invariant under the rational points and the level, the integral is a
finite sum over double cosets, weighted by the index of the level and the orders of the stabilizers.

% before: is-eigenvalue-conjugate-frame-vacuum
The level of $V_i$ acts on the conjugate theta functions through the two pieces from which a level
vector is built, its vacuum and its level function. An element that stabilizes both lattices multiplies
the lattice vacuum by values of the characters of the Weil datum at determinants, and the level function
by values of the residue characters. A place model carries these two scalars to the level vector.

% before: is-eigenvalue-conjugate-frame-level-vector
A level vector is a place model of its vacuum applied to a level function. A pair of elements that
stabilizes the two lattices moves the vacuum and the level function by scalars, so it multiplies the
level vector by their quotient.

% before: is-eigenvalue-conjugate-frame-level-character
For the level of $V_i$ the scalar of the vacuum is a value of the Hecke character $\mu^0_i$ at a
determinant, while the action to be proved involves the level character of the realization. The two
agree on the level: both are products of the residue characters at the places of $P_g$, at the
determinants of the lower right blocks of the reductions.

% before: is-eigenvalue-conjugate-frame-fibres
The moved vacuum of the conjugate Weil representation is compared with the core through the complex
uniformization of the fibres of its universal abelian scheme. A frame gives coordinates on
$\mathbf C^d$ and periods of lattice vectors, and a symplectic basis gives a period matrix and a Riemann
form. An analytic fibre presents the fibre at a framed point as the quotient of $\mathbf C^d$ by the
periods, with a trivialization of the line bundle of the polarization, and a theta reading is a section
written in that trivialization.

% before: is-eigenvalue-conjugate-frame-situation
The constants of the frames are found by comparing the Weil representation of a pair $(S,T)$ with theta
functions of a flag datum whose matrix and element come from $T$. The archimedean components of $T$ act on
families of matrices that are orthonormal for the signature of $T$, and the theta function of a
Schwartz-Bruhat function is evaluated at the columns of the real points of $\mathbb W$.

% before: is-eigenvalue-conjugate-frame-glue
The archimedean components of $U(V)(\mathbf A)$ are read through the members of the archimedean family,
one for each embedding in $\mathcal T$. As $\mathcal T$ contains one embedding of each conjugate pair,
these readings determine the archimedean component, and every family of unitary matrices is the reading of
an element with finite component $1$.

% before: is-eigenvalue-conjugate-frame-signature
A family orthonormal for the signature is a point of the flag domain and gives a Cartan involution of
$T$. Moving it by the archimedean components of $U(T)(\mathbf A)$ keeps it orthonormal, and an element
fixes it exactly when its archimedean component is $1$; any two such families differ by such a move.

% before: is-eigenvalue-conjugate-frame-det-roots
The scalars by which the Weil representation acts are values of characters at determinant roots,
ideles whose quotient by their conjugate is the determinant. An element with archimedean component $1$
that stabilizes a lattice has a determinant root that is a unit at every finite place, as $E_h/F_h$ is
unramified at the finite places.

% before: is-eigenvalue-conjugate-frame-finite-vacuum
On a Fock vacuum vector, the elements with archimedean component $1$ that stabilize the two lattices act
through the characters at determinant roots. This is the finite part of the splitting of the Weil
representation, as in Lemma~\ref{is-eigenvalue-conjugate-frame-vacuum}(1).

% before: is-eigenvalue-conjugate-frame-heisenberg
The finite part of a vacuum is described by a Schwartz-Bruhat function on rational matrices, on which the
lattice acts through Heisenberg operators written in the coordinates of a symplectic basis. A vacuum of the
tensor lattice is a function of norm one fixed by these operators twisted by signs, and the level-one theta
functions are the corresponding functions on $\mathbf C^d$.

% before: is-eigenvalue-conjugate-frame-context
The theta functions of the next lemmas are read at frames of a flag datum attached to $T$ with the conjugate
embeddings. Its matrix is the transpose of the Gram matrix of $T$, so a basis of $T$ at each embedding that
is orthonormal for the signs of the flag datum gives a frame, and its Cartan involution acts on the real
points.

% before: is-eigenvalue-conjugate-frame-conjugate-model
The archimedean elements act on theta functions through their real points, and it is convenient to read
these points in a second flag datum with the conjugate embeddings and the opposite signature. Complex
conjugation carries the coordinates, periods and theta series of the first flag datum to those of the
second, so the two describe the same theta functions.

% before: is-eigenvalue-conjugate-frame-zak-coordinates
The theta functions at the real points are read through the coordinates of a frame. These coordinates identify the
real points with $(\mathbf C^d)^n$, carry the Riemann form to the imaginary part of $H_x$, and are moved by the unitary
elements together with the frame.

% before: is-eigenvalue-conjugate-frame-zak-series
At a point of the flag domain the theta series is a Siegel theta series. It converges, is holomorphic in the variable,
and determines the Schwartz-Bruhat function.

% before: is-eigenvalue-conjugate-frame-zak-weil
A unitary element moves the point of the flag domain. The theta series at the moved point of a suitable
Schwartz-Bruhat function is the theta series at the original point, and this Weil operator is determined by the
property.

% before: is-eigenvalue-conjugate-frame-zak-real
The Zak theta functions are built from the theta functions of Schwartz-Bruhat functions at the real points.
These depend linearly and injectively on the Schwartz-Bruhat function, and the Heisenberg operators and the
unitary elements act on them by translating and moving the point.

% before: is-eigenvalue-conjugate-frame-zak-points
The real points of $\mathbb W$, those with finite component $0$, are read through the columns of their archimedean
components at the embeddings of $\mathcal T$. This identifies them with the real vector space on which the theta
functions of Schwartz-Bruhat functions are defined, compatibly with the unitary groups and the real directions.

% before: is-eigenvalue-conjugate-frame-zak-forms
On the real points the pairing of $\mathbb W$ becomes the Riemann form of the second flag datum, and on the rational
and finite points it becomes the trace of its alternating form. Each identity is a computation with the standard
character.

% before: is-eigenvalue-conjugate-frame-zak-space
A theta function is controlled by its values at the real points, and the translations by rational elements
move these points among themselves. This is what allows a theta function to be defined by prescribing its
values there.

% before: is-eigenvalue-conjugate-frame-zak-exists-lattice
The adelic closure of a tensor lattice is a compact group of finite points. By strong approximation every finite
point is a rational point modulo it, and its multiples shrink into any open subgroup.

% before: is-eigenvalue-conjugate-frame-zak-exists-extension
A function on an open subgroup that contains the real points extends to all of $\mathbb W$ through the rational
points, which cover the quotient. The extension is a theta function when the function transforms under the rational
points of the subgroup as a theta function does.

% before: is-eigenvalue-conjugate-frame-zak-exists
Prescribed values at the real points extend to a theta function. The extension uses an open subgroup that
contains the real points and the rational points that cover the rest.

% before: is-eigenvalue-conjugate-frame-zak-fock
At an orthonormal frame the annihilation and creation operators can be read in the coordinates of the frame.
On a Gaussian times a holomorphic function the annihilation operators vanish, and the polynomial vectors
become derivatives of the holomorphic function at $0$.

% before: is-eigenvalue-conjugate-frame-zak
A theta function on $\mathbb W$ is determined by its values at the real points, and every Schwartz-Bruhat
function has one whose values there are its theta function at a frame. These theta functions depend linearly
and injectively on the Schwartz-Bruhat function, and at an orthonormal frame they are annihilated by the
annihilation operators of its Cartan involution.

% before: is-eigenvalue-conjugate-frame-root-values
The scalars of the Weil representation are values of $\chi_S$ at determinant roots. Two determinant roots
of an element differ by an idele of $F_h$, on which $\chi_S$ is a power of the quadratic character with
even exponent, so the value does not depend on the root.

% before: is-eigenvalue-conjugate-frame-compact-character
A multiplicative function on the archimedean elements is determined by its values on those that commute
with a Cartan involution. The polar decomposition writes every archimedean element as such an element times
the square of an element that the involution inverts.

% before: is-eigenvalue-conjugate-frame-fock-character-lattice
The Weil datum with the powers of $\tilde\eta$ as characters is built from two partial actions, one through the
finite places and one through the archimedean places. The first is a lattice model for given self-dual lattices. At
a place that does not split it fixes the vectors invariant under the part of the lattice of the pair at that place.

% before: is-eigenvalue-conjugate-frame-fock-character-oscillator
The archimedean action comes from an oscillator action of a unitary group on a Heisenberg representation of a
complex vector space. A vector annihilated by the annihilation operators of a Cartan involution is an eigenvector of
the elements that commute with the involution, with a determinant as eigenvalue.

% before: is-eigenvalue-conjugate-frame-fock-character-fock
The Fock model is the oscillator action for the real points of $\mathbb W$, which act on a space of quasi-periodic
functions by the Heisenberg operators. The Cartan involutions of $S$ and $T$ give one of the whole space, and the
determinant of the oscillator lemma splits into one factor for each archimedean place.

% before: is-eigenvalue-conjugate-frame-fock-character-product
The two models act on theta functions through different places, so they commute. Their product is continuous, and
it fixes the value at $0$ under the rational points of the two unitary groups when the characters are the powers of
$\tilde\eta$.

% before: is-eigenvalue-conjugate-frame-fock-character-model
The Weil datum built from the lattice model and the Fock model has the powers of $\tilde\eta$ as its
characters. At the elements supported at one finite place it acts through the lattice model, and at the archimedean
elements through the Fock model, which multiplies a vector annihilated by the annihilation operators by a product of
determinants.

% before: is-eigenvalue-conjugate-frame-fock-character
On a vector annihilated by the annihilation operators, the archimedean elements that commute with the
Cartan involutions act by a scalar. It is a product of values of the characters of the Weil datum at
determinant roots and of determinants on the holomorphic parts.

% before: is-eigenvalue-conjugate-frame-compact-vacuum
When the Cartan involution of $S$ is $1$, the scalar becomes the value of $\chi_S$ at a determinant root
times the quotient of the determinants on the two eigenspaces of the Cartan involution of $T$. A
determinant root whose archimedean components are square roots of the determinants turns the values of the
Hecke character into these determinants.

% before: is-eigenvalue-conjugate-frame-theta-space
Two theta functions are compared below through their values at the points with finite component $0$, where
both are given by theta series. These values determine a theta function: strong approximation writes every
point as a rational point plus such a point plus a multiple of a lattice point, and the three conditions
on theta functions remove the first and the last.

% before: is-eigenvalue-conjugate-frame-compact-move
Moving an orthonormal family by an element that commutes with its Cartan involution changes the frame by a
block diagonal matrix. The theta function at the moved frame is the same multiple of the theta function at
the original frame.

% before: is-eigenvalue-conjugate-frame-overlap-doubling
The inner product of two theta functions is computed below by doubling. The product of a theta series at one point
with the conjugate of a theta series at another is a Hermitian theta series on $n\times r$ matrices over $E$, with
an explicit period matrix and an explicit kernel in front. The value at $0$ of its Schwartz-Bruhat function and the
kernel at $0$ give the inner product.

% before: is-eigenvalue-conjugate-frame-overlap-determinants
The kernel and the period matrix of the doubling involve the inverses of two matrices built from the two points. On
the flag domain the first columns of a point span a positive subspace and their orthogonal complement is negative, and
this makes both matrices invertible and the doubled period matrix lie in the Hermitian upper half space.

% before: is-eigenvalue-conjugate-frame-overlap-series
A Hermitian theta series with such a period matrix is a Siegel theta series in rational coordinates. So it converges
absolutely, and it determines its Schwartz-Bruhat function.

% before: is-eigenvalue-conjugate-frame-overlap-doubling-exists
For an even number of rows the doubling exists. As the doubled theta series determines its Schwartz-Bruhat function,
the doubling is unique.

% before: is-eigenvalue-conjugate-frame-overlap-kernel-shift
Translating both arguments by the same real point multiplies the product of the two Gaussians and the kernel of the
doubling by a phase. That phase is the one in the mean below, so the translation leaves the rest of the integrand a
Fourier series in the translating point.

% before: is-eigenvalue-conjugate-frame-overlap-mean
After the doubling and the translation law, the product of two theta functions at the real points is a Fourier series
on a real torus. Its mean over the torus is its constant term, the value at $0$ of the doubling times the kernel.

% before: is-eigenvalue-conjugate-frame-real-character-overlap
The theta functions carry an inner product, an integral over a fundamental domain of the rational points. The inner
product of two Zak theta functions is the value at $0$ of the doubling times the kernel at the two points. So the inner
product of the theta functions of one Schwartz-Bruhat function at two orthonormal frames is not zero, and it does not
change when both frames are moved by the same element.

% before: is-eigenvalue-conjugate-frame-real-character-scalar-fock
The operators of the Weil action move the Heisenberg operators by the action on $\mathbb W$. Differentiating along the
real directions, they move the annihilation operators for one Cartan involution to those for the moved one, so they
carry Fock vacua at an orthonormal family to Fock vacua at the moved family.

% before: is-eigenvalue-conjugate-frame-real-character-scalar-schur
The Heisenberg operators act irreducibly on the Schwartz-Bruhat functions. So a linear map commuting with all of them is
a scalar, as in Schur's lemma.

% before: is-eigenvalue-conjugate-frame-real-character-scalar-norm
At an orthonormal frame the kernel of the overlap formula does not depend on the frame, and the value at $0$ of the
doubling is a fixed multiple of the inner product of the Schwartz-Bruhat functions. So the norm of a Zak theta function at
an orthonormal frame is a fixed multiple of the norm of its Schwartz-Bruhat function.

% before: is-eigenvalue-conjugate-frame-real-character-scalar
An element with finite component $1$ moves the theta function of a Schwartz-Bruhat function at one orthonormal frame
to a multiple of the theta function at the moved frame. The multiple has absolute value $1$, because the Weil action
preserves the inner product.

% before: is-eigenvalue-conjugate-frame-real-character
The archimedean elements carry the theta function at one orthonormal frame to the theta function at the
moved frame. The scalar is a character of the archimedean elements, the same for every frame and every
Schwartz-Bruhat function.

% before: is-eigenvalue-conjugate-frame-real-move
The elements with finite component $1$ move the theta function of a Schwartz-Bruhat function to the theta
function at the moved family, up to the value of $\chi_S$ at a determinant root. Together with the
previous lemma this controls $\omega(1,g)$ on a vacuum for every $g$.

% before: is-eigenvalue-conjugate-frame-moved-vacuum
To compare the vacuum at different points one moves it by the Weil representation. The moved vector is
again a Fock vacuum vector, for the moved Cartan involutions and the moved lattice.

% before: is-eigenvalue-conjugate-frame-move
If two elements of $U(T)(\mathbf A)$ carry the lattice of $T$ to the same lattice, the moved vacua are
given by one Schwartz-Bruhat function. The two moves differ by an element with archimedean component $1$
stabilizing the lattice and an element with finite component $1$, and the two previous lemmas give the
factor $\mu(g_1)\mu(g_0)^{-1}$.

% before: is-eigenvalue-conjugate-frame-vacua-mean
The inner product of two Schwartz-Bruhat functions is the mean of a third, a normalized finite sum over the
classes modulo a multiple of the integral matrices. The mean is linear and invariant under translations,
because translating a function only permutes these classes.

% before: is-eigenvalue-conjugate-frame-vacua-heisenberg
The Heisenberg operators act on Schwartz-Bruhat functions through the coordinates in the symplectic basis.
In these coordinates the Riemann form becomes the standard symplectic pairing, so the operators compose up
to a root of unity and preserve the inner product. On theta series they act by translating the variable by a
period, with an exponential factor.

% before: is-eigenvalue-conjugate-frame-vacua-tensor
A vacuum is fixed by the operators of a whole lattice, and these can only fix a common vector when the lattice
is its own dual for the Riemann form and the signs form a semicharacter on it. The tensor lattice of two
self-dual modules has both properties, because its members are read through contractions with a dual family.

% before: is-eigenvalue-conjugate-frame-vacua-lattice
For a self-dual lattice of rational coordinate pairs with a semicharacter of finite order, the fixed vectors
can be written down. A base point is chosen so that the pairs without translation part act trivially there,
and the vector is then determined along the translates of the base point.

% before: is-eigenvalue-conjugate-frame-vacua
The tensor lattice of $\mathcal O_E^n$ and a self-dual module has a vacuum, and it is unique up to a
scalar. Its theta series is the candidate for the moved vacuum at an orthonormal frame.

% before: is-eigenvalue-conjugate-frame-theta-product-vacuum
In the coordinates of a symplectic basis of the self-dual module, the periods of a point give a matrix in the Siegel
upper half space. The theta series of the vacuum is then a product of classical theta functions with half-integral
characteristics, one for each column.

% before: is-eigenvalue-conjugate-frame-theta-product-unique-fourier
A holomorphic function periodic under $\mathbf Z^g$ has Fourier coefficients, and a second, quasi-periodic law under
the columns of a matrix $\tau$ relates the coefficients at $m$ and at $m+ne_j$. So the coefficients at the $n^g$
residues determine the function.

% before: is-eigenvalue-conjugate-frame-theta-product-unique-normal-form
A symplectic basis of a lattice for the imaginary part of a positive Hermitian form puts the lattice in the form
$\mathbf Z^g+\tau\mathbf Z^g$ with $\tau$ in the Siegel upper half space. In these coordinates the factor of automorphy
becomes the classical one, up to the exponential of a quadratic function.

% before: is-eigenvalue-conjugate-frame-theta-product-unique-count
Multiplying by the exponential of the quadratic function turns a theta function of level $n$ into a function with the
classical laws. So the normal form and the Fourier coefficients bound the number of independent theta functions of
level $n$ for a unimodular form by $n^g$.

% before: is-eigenvalue-conjugate-frame-theta-product-unique-proportional
A lattice that is self-dual for an integral rational form gives, through an injective period map, a discrete lattice
on which the imaginary part of the Hermitian form is unimodular. The bound for level one then says that two theta
functions for it are proportional.

% before: is-eigenvalue-conjugate-frame-theta-product-unique
The functional equations of a level-one theta function are those of the sections of a principal polarization. So
they determine the function up to a scalar.

% before: is-eigenvalue-conjugate-frame-theta-product
A product of level-one theta functions, one for each column, satisfies the same functional equations as
the theta series of the vacuum. The two agree up to a constant.

% before: is-eigenvalue-conjugate-frame-vacuum-line-idempotents
An integer $w_0$ with $w_0+\bar w_0=1$ and $w_0\bar w_0$ even forces every place above $2$ to split, and at each of
them it is congruent to $1$ at one place and to $0$ at the other. So $w_0$ is a sum of split idempotents modulo $2$.

% before: is-eigenvalue-conjugate-frame-vacuum-line-signs
The lattice vacuum conditions ask for invariance under twice the lattice and under its halves at the split
places. For the lattice of a tensor product they amount to one condition for each element of the tensor
lattice, with the sign of that element.

% before: is-eigenvalue-conjugate-frame-vacuum-line-series
A holomorphic function whose product with the Gaussian is invariant, up to the Heisenberg factor, under a multiple of
the lattice periods is quasi-periodic for a sublattice. Such functions are finite combinations of theta series, and
so they are theta series of Schwartz-Bruhat functions.

% before: is-eigenvalue-conjugate-frame-vacuum-line-fock
At the real points a theta function is a Gaussian times a function of the coordinates of the frame. When the
annihilation operators of an orthonormal frame kill it, this function is holomorphic, and the theta function
comes from a Schwartz-Bruhat function.

% before: is-eigenvalue-conjugate-frame-vacuum-line
A Fock vacuum vector at an orthonormal frame is the theta function of a Schwartz-Bruhat function, and its
invariance under the lattice of the pair becomes the vacuum condition on that function. By the uniqueness of
the vacuum, the vector is a multiple of the theta function of the vacuum.

% before: is-eigenvalue-conjugate-frame-torus-mean-unfold
An integral over a fundamental domain of a lattice in a locally compact group can be unfolded along a tube around a
section of a projection. When the integrand is invariant under the compact fibre of the tube, the integral becomes a
mean over a fundamental domain of the projected lattice.

% before: is-eigenvalue-conjugate-frame-torus-mean
The normalization of a Fock vacuum vector is an integral over a fundamental domain of the rational points.
For theta functions it equals a mean over a real torus, which can be compared with the mean of a theta
reading.

% before: is-eigenvalue-conjugate-frame-tensor-lattice
The Fock vacuum vectors at the points of the core are invariant under the lattice of a pair of lattices,
while the vacuum condition on Schwartz-Bruhat functions is stated for a tensor lattice. The two subgroups
of $\mathbb W$ agree when the second lattice is the first one multiplied by the ideal $\mathfrak A$.

% before: is-eigenvalue-conjugate-frame-unit
At a single point the moved vacuum is compared with the analytic fibre. A vacuum is a theta function of a
vacuum of the tensor lattice, the product of level-one theta functions is the theta series of that vacuum
up to a constant, and the normalizations of both sides make this constant a unit.

% before: is-eigenvalue-conjugate-frame-core
The remaining inputs concern the dictionary of the core at level $i+1$ and its uniformization at the
restriction $s$ of $\sigma$. The situation collects the data of the frame constants, without the Weil
datum.

% before: is-eigenvalue-conjugate-frame-framing
At every point of the dictionary the core gives a rational element that moves the point to the base
frames, together with a lattice, a symplectic basis, letters and an analytic fibre. These data place the
point in the situation of the general lemmas on the Weil representation.

% before: is-eigenvalue-conjugate-frame-marking
The action of $G$ on the core changes the level structures of framed points, so comparing frame constants
at $h$ and at $hk$ needs to know how a level structure is read in an analytic fibre. A marking at an adelic
point is a map from the lattice of the core to $E_h^r$ that agrees with that point modulo the level, and a
marked analytic fibre reads the level structure through it.

% before: is-eigenvalue-conjugate-frame-uniformization
The framed points of the core are read through a uniformization of its moduli scheme at a good embedding.
The facts used about it hold for any uniformization that satisfies the compatibility conditions of the
moduli problem, and they involve the family uniformization behind it, whose lattices give each component a
marking of the lattice of the datum.

% before: is-eigenvalue-conjugate-frame-dictionary-point
The dictionary of the core sends a point of $U(W_i)(\mathbf A)$ to a component and a frame. The lattice law
of the dictionary and the base frames of the model show that a rational element moves the point so that the
component's module and the frame come from it.

% before: is-eigenvalue-conjugate-frame-component-lattice-positive
The sign of $\kappa_0$ is read from positivity. At every point of the flag domain the periods carry a
positive form whose imaginary part is the form of the flag data; it is built place by place from a
triangular reduction of the Gram matrix of the frame.

% before: is-eigenvalue-conjugate-frame-component-lattice
The uniformization attaches to each component a lattice with a Riemann form. Comparing it with the form of
the flag datum gives the normalization of $\kappa_0$ and the self-duality of the module.

% before: is-eigenvalue-conjugate-frame-letter-dictionary
The polynomials that the one-map data attach to coefficient vectors are written in letter variables of the
datum, while the theta series are differentiated in the coordinates of a flag datum. A letter dictionary
identifies the two sets of variables through the indexing and the bijection onto the letters, and matches
the highest weight polynomials of the datum with those of the flag datum.

% before: is-eigenvalue-conjugate-frame-letters
The theta readings are formed with a symplectic basis and with coordinates labelled by the letters of the
datum. A self-dual module of the right signature has such a basis and indexing, matched with any given
labelling by the letters.

% before: is-eigenvalue-conjugate-frame-analytic-fibre
At a frame whose minors have absolute value one, the uniformization provides the period map, the
trivializing functions of the line bundle and a theta reading of mean one. These are the data of an
analytic fibre.

% before: is-eigenvalue-conjugate-frame-points
The facts below are what the uniformization of the core and the conditions on the realization supply
for the frame constants. They give a framing at every point, the base frames as an orthonormal family, the
moved lattices, and the values of the multiplier at determinant roots.

% before: is-eigenvalue-conjugate-frame-ortho-constant-norm-theta-mean
The mean of the squared absolute value of a theta function with characteristics, weighted by the Gaussian of the
imaginary part, is a classical Gaussian integral. It depends only on the determinant of the imaginary part of the period
matrix.

% before: is-eigenvalue-conjugate-frame-ortho-constant-norm-mean
A multiple of a theta function with characteristics that has mean one therefore has a multiplier of known absolute
value. The weighted square of its absolute value is invariant under the lattice, so its mean does not depend on the
fundamental domain.

% before: is-eigenvalue-conjugate-frame-ortho-constant-norm-det
At a frame whose minors have absolute value $1$, the triangular change to an orthonormal frame has diagonal entries of
absolute value $1$. So the Hermitian form in period coordinates has the same determinant as at the orthonormal frame,
the product of the numbers $c_\tau^r$.

% before: is-eigenvalue-conjugate-frame-ortho-constant-norm
At a frame whose minors have absolute value $1$, a level-one theta function of mean one is normalized. So the factor
relating the product of its values to the theta series of the vacuum has an absolute value that does not depend on the
frame.

% before: is-eigenvalue-conjugate-frame-ortho-constant-unit-factor
A family of pairs of holomorphic functions that are proportional at every point, with one factor for the whole family
at each point and factors of one absolute value, is proportional with one constant. Near a point where a denominator does
not vanish the ratio has constant modulus, so it is constant, and the identity theorem carries the proportionality to the
whole connected domain.

% before: is-eigenvalue-conjugate-frame-ortho-constant-unit
Two holomorphic functions with the same weight law that are proportional at the frames with unit minors, with factors of
one absolute value, are proportional with one constant. The constant can be chosen for all such pairs at once.

% before: is-eigenvalue-conjugate-frame-ortho-constant
At the orthonormal points of one component, with one symplectic basis and indexing, the products of theta
readings and the theta series of a vacuum satisfy the same functional equations. The constant relating them
is the same at every such point.

% before: is-eigenvalue-conjugate-frame-hermitian-riemann
The period matrices of a point of the flag domain satisfy Riemann's relations. A triangular change of frame
moves the point to an orthonormal one and multiplies the periods by an invertible matrix, and at an
orthonormal point the Riemann form on the periods is the imaginary part of a positive diagonal form.

% before: is-eigenvalue-conjugate-frame-hermitian
The Gaussians and the transition laws of analytic fibres are written with the Hermitian form $H_x$ built
from the period matrices of a symplectic basis. This form is the Hermitian form of the Riemann form on the
periods, so it is determined by its imaginary part there, and it is invariant under the unitary group of
the flag datum.

% before: is-eigenvalue-conjugate-frame-level-one
The generating section of the line bundle at a framed point, read in an analytic fibre, satisfies the
functional equation of a level-one theta function. Its period law comes from the transition functions of
the fibre, and its sign is the one fixed by $w_0$ on a self-dual module.

% before: is-eigenvalue-conjugate-frame-renumbering
The symplectic basis and the indexing are choices, and a theta reading depends on them. Changing them
renumbers the coordinates of $\mathbf C^d$, and the analytic fibre and its theta reading can be transported
so that the product of the Gaussian and the theta readings is unchanged.

% before: is-eigenvalue-conjugate-frame-series
A component has a single theta series: the products of theta readings at its orthonormal points are the
theta function of one Schwartz-Bruhat function, whatever symplectic basis and indexing are used. This is
what makes the constant of a component independent of the point.

% before: is-eigenvalue-conjugate-frame-marking-laws
Markings exist at every adelic point, and a level lift of an element of $G$ carries a marking at $h$ to a
marking at $hk$, composed with the change of level structure of that element. The level lift also fixes the
lattice of $W_i$, and the multiplier at it is the inverse of the residue characters at $P_g$.

% before: is-eigenvalue-conjugate-frame-level-move
Changing the level structure of a framed point does not change its abelian variety or its polarization. So
an analytic fibre at the point gives one at the moved point, with the same theta readings and the composed
marking.

% before: is-eigenvalue-conjugate-frame-compare-scalar
An analytic fibre with its marking pins down the framed point up to its generator. Two framed points that carry
analytic fibres with the same periods and the same marking therefore differ only by a scalar on the generator, and
the theta readings and the values of sections change by that scalar and its powers.

% before: is-eigenvalue-conjugate-frame-compare
Two framed points whose markings correspond under a rational element are isomorphic, and the isomorphism
multiplies the line bundle of the polarization by a scalar. This scalar relates the values of sections at
the two points and the products of theta readings at corresponding real points.

% before: is-eigenvalue-conjugate-frame-extremal-orthogonal
A smooth vector killed by the upper triangular part of the complexified Lie algebra, in the frame of the model,
has matrix coefficients that are holomorphic sections over the frame domain. Such a section is determined by its
torus law and its value at $1$, so in an irreducible unitary representation a vector of the same torus weight
orthogonal to it vanishes.

% before: is-eigenvalue-conjugate-frame-extremal-fock
Polynomial vectors over a theta function are formed with creation operators. These commute with each other,
and the annihilation operators satisfy the canonical commutation relations with them, because the derivatives
of the Heisenberg action have the symplectic form as their commutator. Over a nonzero vacuum this makes the map
from polynomials to polynomial vectors injective.

% before: is-eigenvalue-conjugate-frame-extremal-compact
The diagonal torus of the maximal compact subgroup acts on polynomial vectors through a character and a
rescaling of the creation variables. The exponents of the character and the weight of the lowest $K$-type of
the model are both read from the parameter of the label.

% before: is-eigenvalue-conjugate-frame-extremal-line
Four facts about the holomorphic model and the delta polynomial are used in the next lemmas. The vectors of
extremal weight in a lowest $K$-type form a line, distinct polynomials give distinct polynomial vectors over the
level vector, the diagonal torus acts on the vector of the delta polynomial through the extremal weight, and the
lowest $K$-type is irreducible.

% before: is-eigenvalue-conjugate-frame-rebase-substitution
Changing the bases of $W$ by unitary matrices that commute with the signature matrix changes each creation
operator into a combination of creation operators of the same sign. On polynomial vectors this is a linear
substitution of the creation variables, with the entries of the matrices or their conjugates as coefficients.

% before: is-eigenvalue-conjugate-frame-rebase
The polynomial vectors of the conjugate datum are formed with some orthonormal bases of $W_i$ adapted to the
Cartan involution, while the frames of the model give a particular such basis. A change of bases at the
archimedean places moves the polynomials, so the map from lowest $K$-types and coefficients to polynomials
can be taken with the frames of the model.

% before: is-eigenvalue-conjugate-frame-extremal-unit
The map from lowest $K$-types and coefficients to polynomials gives the delta polynomial at some unit
vectors. An extremal vector of the holomorphic model is unique up to a phase, so it gives the delta
polynomial at a unit coefficient vector.

% before: is-eigenvalue-conjugate-frame-phase-proportional
Two linear maps out of an irreducible representation that intertwine it with the same operators, up to scalars,
are proportional once they agree at one vector. The scalars are forced to be $1$ on elements of nonzero trace, and
these generate the group.

% before: is-eigenvalue-conjugate-frame-phase-eta-extension
The Hecke character needed for the Fock vacua is built from local data: trivial on the units at the finite places,
a sign at the uniformizers of the places that do not split, and a phase at the archimedean places. Such data define a
character of the group generated by the ideles that are units at every finite place and the principal ideles, once
they agree on the principal ideles that are units. That agreement comes down to the roots of unity of $E$, which is
why the hypothesis on them appears.

% before: is-eigenvalue-conjugate-frame-phase-eta-product
The agreement on principal ideles uses the product of the local signs at an element of $F_h$. For the host field this
product is $1$ by the reciprocity law of class field theory, since no finite place ramifies.

% before: is-eigenvalue-conjugate-frame-phase-eta-unique
Class field theory also characterizes the quadratic character of $E/F$ among the characters of the ideles of $F$.
This characterization identifies the restriction of the constructed character with it.

% before: is-eigenvalue-conjugate-frame-phase-eta-roots
For the host field, the hypothesis on roots of unity is met by elements that are units away from the primes dividing
$N$. They are $\sqrt{-d}$ for the root $-1$ and $1+\zeta$ for the others, and their phases cancel in pairs over the
embeddings of $\mathcal T$.

% before: is-eigenvalue-conjugate-frame-phase-eta-split
These elements may fail to be units at the primes dividing $N$. The places of $F_h$ over those primes split in $E_h$,
so the parity condition does not concern them.

% before: is-eigenvalue-conjugate-frame-phase-eta-reading
The constructed character is known on the units, on the uniformizers and on the archimedean parts of principal
ideles. Continuity extends this to every idele with archimedean component $1$, and the density of $E_h$ in its
archimedean completion to every idele with finite component $1$.

% before: is-eigenvalue-conjugate-frame-phase-eta
The action of a compact element on a Fock vacuum is computed by comparing the Weil datum with one whose characters are
powers of a fixed Hecke character. That character must be unramified at every finite place, restrict to the quadratic
character and have exponent one half at every embedding, and every definite-side datum has one.

% before: is-eigenvalue-conjugate-frame-phase-vacuum
When the source space is positive definite, the scalar by which a pair of archimedean elements moves a Fock vacuum is a
product of determinants. The exponents come from the archimedean exponents of the two characters, the rank of the
source and the trace of the Cartan involution of the target.

% before: is-eigenvalue-conjugate-frame-phase-connected
The rational points of $U(V_i)$ generate the group and act irreducibly because they are dense in the archimedean
points, which form a connected group. Connectedness comes from the Cayley transform, which joins every unitary matrix
$k$ with $1+k$ invertible to the identity, and a scalar of norm one moves every matrix into that set.

% before: is-eigenvalue-conjugate-frame-phase
Two maps from the coefficients of $V_i$ to theta functions are compared: the one given by the conjugate
datum and the one given by the polynomials that the one-map data attach to rational coefficient vectors.
Both are equivariant under the rational points of $U(V_i)$ and both give the theta function of the delta
polynomial at the highest weight vector, so they differ by a constant of absolute value one.

% before: is-eigenvalue-conjugate-frame-central
The class sums of the core are moved by central elements of $G$ that act on the level structures by an
integer. Such an element exists for every integer prime to the level, as a scalar at the places of $P_g$.

% before: is-eigenvalue-conjugate-frame-cover-sections
Sections of the moduli scheme are rules on framed points over reduced finitely generated algebras, and their values
at complex framed points are read through such points by base change. Since these algebras are Jacobson rings, a
section is determined by its complex values, and a section of weight zero does not see the frame of a point.

% before: is-eigenvalue-conjugate-frame-cover-situation
A section of weight zero is constant once its value at every complex point is known, and the dictionary of a
uniformization gives these values at its own points. To reach every complex point one compares the lattices that the
uniformization attaches to its components with the lattices of adelic elements of $U(W)$, through markings.

% before: is-eigenvalue-conjugate-frame-cover-isometry
Two markings of components by adelic elements that differ by a matrix $\delta$ on the finite vectors give an isometry
of the two component lattices. It is $\delta$ in the coordinates of the flag data, and it respects the level
structures modulo $M$.

% before: is-eigenvalue-conjugate-frame-cover-compare
A point of the moduli scheme carries Tate systems for two lattices: the one of its level structures and the one read
off from a period. Comparing them at each level gives maps between the lattices that are similitudes modulo that
level, with factors that agree from level to level.

% before: is-eigenvalue-conjugate-frame-cover-norm
The similitude factors are removed by multiplying the maps by elements of $\mathcal O_E$ whose norms invert them.
Such elements exist modulo every integer, and a limit argument makes them compatible across the levels and congruent
to $1$ modulo $\mathfrak M$.

% before: is-eigenvalue-conjugate-frame-cover-tower
The tower of maps from the lattice of the datum to the lattice of a component is built at a period of that component.
The maps preserve the forms modulo each level and agree with the maps of the family uniformization modulo $M'$ on the
part where the level structure is defined.

% before: is-eigenvalue-conjugate-frame-cover-marked
Every component lattice is the lattice of some adelic element, at which its marking is a marking. The element is
the limit of the maps that the level structures of a point of the component give at all levels.

% before: is-eigenvalue-conjugate-frame-cover-over
Every complex framed point has the point of a framed point at an orthonormal frame of some component. Moving that
frame and a marked element of the component to a dictionary point changes the periods by a rational unitary matrix,
and the isometry of the component lattices shows that the two points agree.

% before: is-eigenvalue-conjugate-frame-cover-constant
The sections built from the theta class data carry the delta section of the one-map data as a factor. On
the complex points over $\sigma'$ this section is a constant of absolute value one.

% before: is-eigenvalue-conjugate-frame-inputs-roots
The class sums of the delta data are indexed by tuples of classes modulo the level ideal, and each class is lifted to a
root of a marking divided by the level. Conjugating the coefficients dilates the classes by an integer prime to the
level, and the roots, the markings and the sums over the classes follow the dilation.

% before: is-eigenvalue-conjugate-frame-inputs
The comparison with the one-map data uses three facts about the conjugate Weil representation and the
frames of the core at $\sigma$. The level of $V_i$ acts on the conjugate theta functions through the
conjugated level character. The frames give constants of absolute value one for which the action of $G$
becomes a translation, and an identity between the conjugated theta series at the class
representatives and the conjugated sections of the core.

% before: is-eigenvalue-conjugate-theta-sum
With the action of the level on the theta functions, the integrand of the theta integral is invariant
on the right under the level. The unfolding then writes the theta integral as a sum over
representatives of the double cosets, which is the form in which the one-map data present the lift.

% before: is-eigenvalue-conjugate-inputs
The comparison of a rational lift with its Galois conjugates rests on three properties of the core
forms and of conjugate data. The core forms give elements of the Hecke modules with prescribed
functions at the extremal vectors, the theta integral of a conjugate datum is determined up to a unit
by its characters, and away from the base embedding a conjugate datum exists whose theta integral is
the conjugated class sum of core forms.

% before: is-eigenvalue-base-inputs
The hypotheses of the lower norm bound give the conditions used below: the core conditions, the
rationality of the lift matrices of the rows, and the comparison of the rational lifts with their
Galois conjugates. No place of the level of the base tuple divides the integer of the host, so the
eigenvalue formula applies at the base tuple.

% before: is-eigenvalue-base-place
At the base place the lift of the realization is a theta lift followed by an orthogonal projection.
Its positive eigenvalues are therefore at least the smallest eigenvalue of the theta lift on the Hecke
eigenspaces.

% before: is-eigenvalue-other-places
At the other infinite places the metric of the rows is the class-sum pairing at the corresponding
embedding. The rational lifts identify the lift there with a conjugate theta lift, to which the
eigenvalue formula applies at that embedding.

% before: is-lift-eigenvalue-bound
The lower norm bound is an inequality for the eigenvalues of the lifts at the base tuple. With the
adelic constants of the realization these eigenvalues are explicit, and they are bounded below by the
hook-length term of Lemma~\ref{is-hook-lower-bound}.

% before: is-surjectivity-and-lower-norm
The two bounds on the lifts are where the published theory of unitary groups enters. The lower norm
bound compares a lift's effect on Petersson metrics with hook-length sums of its weights, and the
surjectivity bound controls the image lattices at twisted places through Galois representations and
Green's description of cuspidal characters.

% before: is-cauchy-binet
A step of a chain of lifts is a sum over pairs of lattice classes, each term an operator between
the spaces of one class. Cauchy-Binet data record this decomposition at every infinite place. The
norm of the step is then controlled by the number of classes and by the largest single term.

% before: is-cauchy-binet-bound
The decomposition turns the operator norm of a step into a Cauchy-Schwarz estimate over the
classes. The bound loses the number of classes at the two levels, which is why the counts of
lattice classes enter the sizes of the lifts.

% before: is-definite-side-condition
The period of the iterated lift must be written through explicit rational matrices whose norms are
known, one step for each lift. The definite-side condition records such a factorization for every
word, with a bound on the norm of the final period map and Cauchy-Binet bounds for each step in
terms of the logarithmic sizes of the lifts.

% before: is-slot-rows
The bundles at the members are built from a common recipe, and the conditions on them are verified
for every selection built from that recipe at once. The designed families satisfy, in addition, the
standing conditions on denominators, twists and constants of Chapter~\ref{chapter-assembly}.

% before: is-definite-rows
Chapter~\ref{chapter-assembly} uses the definite side through a single list of bounds at the
realizations of the selected bundles. This theorem assembles the list from the separate conditions
and from Theorem~\ref{is-surjectivity-and-lower-norm}, with one constant $C$ that serves for all of
them.

% before: is-polynomial-model
The theta kernels and the period at the top act on vector-valued forms whose values lie in
representations of $GL_N$. Norms of these operators are computed in the polynomial model of the
representations, with the Fock pairing, in which highest weight vectors are products of minors.

% before: is-standard-basis
The proofs below need two facts about the polynomial model in characteristic $0$. The dual Weyl
module has a basis of products of minors with integer coefficients, which reduces questions about
the Fock pairing to rational polynomials. It is also generated by its extremal vector under the
divided row operators, which lets the extremal vector stand in for the whole representation.

% before: is-fock-facts
The computations in the polynomial model use a few properties of the Fock pairing: it is
multiplicative on tensors, nondegenerate on the representation spaces, and the action of $GL_N$ is
adjoint to the transposed action. Positivity of the Hermitian pairings makes the Cauchy-Schwarz
inequality available.

% before: is-branching-data
Restriction from $GL_{n+1}$ to $GL_n$ is made explicit by a coefficient map and a branching vector
built from minors. The constants $c_{\lambda,\sigma}$ and $\kappa_{\lambda,\sigma}$ are the ones that
appear in the norm of the period at the top.

% before: is-capelli
The functional that compares a vector with the branching vector is controlled by its behaviour under
Capelli's column determinant. It commutes with the row operators, on a dual Weyl module it is a
scalar, and the scalars for different weights can be separated by a choice of the parameter $u$.

% before: is-top-component
A polynomial in the representation of $GL_{n+1}$ is sorted by its degrees in the last row. The
component that is largest in the order $\omega$ has the extremal shape of a branching pattern: a
monomial in the last row times a vector of the representation of $GL_n$ of the complementary weight,
on which the Capelli operators act by the scalars of Lemma~\ref{is-capelli}.

% before: is-capelli-separation
The Capelli operators tell the possible top components apart by their scalars. Two different
admissible indices have different scalar functions of $u$, and a finite list of parameters, read off
from the branching intervals, contains a zero of every such function except that of $b$. The product
of the values of that function over the list is the constant $c_{\lambda,\sigma}$.

% before: is-capelli-kill
Each top component of a polynomial is an eigenvector of the Capelli operators, so products of the
operators $C_n(u)-c$ peel the components off one at a time. This shows that the product of the
Capelli operators over the list of Lemma~\ref{is-capelli-separation} kills every polynomial of
the representation that is homogeneous of degree $d$ in the last row and whose component of index
$b$ vanishes.

% before: is-branching-vector
The norm of the branching map is computed from its value on the branching vector and from its
equivariance. The branching vector lies in the representation of weight $\lambda$, has the extremal
vector of $\sigma$ as its coefficient, and represents the coefficient functional up to a scalar.

% before: is-branching-vector-order
The proof of Lemma~\ref{is-branching-vector}(2) kills the difference $y-a\xi$ with the product
of Capelli operators and then divides by $c_{\lambda,\sigma}$. Ordering the polynomials by their top
value of $\omega$ first makes that difference smaller than $y$, so the constant is not needed there.

% before: is-abstract-coisometry
A map whose composite with its adjoint is a positive scalar $\kappa$ has squared norm at most $\kappa$
for any positive Hermitian forms. This is the only analytic input of the branching norm.

% before: is-branching-norm
The period at the top pairs the branching image of a form with a form on the smaller group. Its
norm is the scalar $\beta$, a product over the types of $T_u^{-n_\tau}c_\tau^2\kappa_\tau$, and this
gives the term $-a_t\log T_u$ of $\mathrm{MAIN}(u)$ and the last-side numbers.

% before: is-height-numbers
The definite-side condition asks for height numbers, a list of explicit inequalities between the
sizes of the lifts, the norm of the period map and the quantities of the record. The proof produces
them in two halves: the kernel side bounds the lifts one step at a time, and the last side bounds
the period map at the top. Each half has its own constants.

% before: is-height-numbers-from-sides
The two halves combine into height numbers once the size constant $c$ of the record is large
compared with their constants. This is where the threshold $C_0$ of the definite-side condition
comes from.

% before: is-chain-from-kernels
Given bounds for the lifts and for the period map in terms of inner products, the chain of
matrices required by the rows condition is obtained by choosing bases. The metrics of the chain are
the Gram matrices of the inner products, so step bounds become operator norm bounds, and these are
the Cauchy-Binet bounds with a single class.

% before: is-kernel-ties
The bounds are proved for the theta kernels of a realization, which carry their own fields, primes,
weights and characters. These must agree with the definite-side datum, and the scale must control
the primes and the weights. The evaluation conditions collect what a bundle has to satisfy for the
proof to apply.

% before: is-jet-operator
A theta kernel between two levels sends a polynomial in the variables of the lower level to the
polynomial whose coefficients are the Taylor coefficients of a theta function. Jet operators
formalize this, and $t(b,\nu,S,\Omega)$ is the bound that the Fock model gives for them.

% before: is-fock-substitution
Substituting a matrix into a polynomial of a given highest weight multiplies its Fock norm by at
most a product of powers of the singular values. This is the dual form of the growth condition on
theta functions.

% before: is-gaussian-moments
Two Gaussian integrals underlie the estimate for jets. The first reads the Taylor coefficients of an
entire function of Gaussian growth off its integral against conjugate monomials, and the second bounds
a product of singular values against a Gaussian, which is where the factors
$\exp(\varphi(\nu_{\tau,i}))$ come from.

% before: is-fock-cauchy
The jets of an entire function of Gaussian growth are Gaussian integrals of the function. In the Fock
model the estimate applies to any functional that is bounded on substitutions, with the bound
$t(b,\nu,S,\Omega)$.

% before: is-block-bound
The two estimates combine into a bound for the jet operator of a function of Gaussian growth, for
step weights. It is the source of the main term $\phi(\nu)-\nu\log\Omega$ of the size of a lift.

% before: is-vacuum-theta-vectors
The vacuum vectors whose norms enter the sizes of the lifts are theta functions on the spaces
$V_{a,b}$ of the host. They are built from the genus-one theta function of a CM lattice by tensor
powers along orthonormal bases and by averaging over cosets of commensurable lattices. Their
conjugates are formed from the conjugates of the lattice and of the types.

% before: is-theta-kernel
The kernel entries of the chain of lifts are sums, over cosets of lattices of homomorphisms, of jet
operators of coset theta functions. Their number and size are governed by the lattice classes, the
level primes and the coupling sets.

% before: is-growth-theta-functions
The size of a kernel entry at a conjugate embedding is controlled by the growth of a theta function
on the Fock space after conjugating its Taylor coefficients. These are the theta functions and theta
vectors up to units of a CM lattice of the host and a pair of Hermitian lattices. They are translated
by torsion points, and conjugated by arbitrary embeddings of $\overline{\mathbf Q}$.

% before: is-growth-period-lattice
A sup bound for a theta function needs a lower bound for the nonzero vectors of the subgroup
$\Lambda_{L,L'}$ and the covolume of that subgroup. Both come from the self-duality of the two
Hermitian lattices and of $\mathfrak b^{-1}$, which makes the imaginary part of the Hermitian form
integral and unimodular on $\Lambda_{L,L'}$. The same computation shows that vectors coming from
matrices over $E$ are torsion points.

% before: is-growth-sup-mean-square
A function whose squared modulus times the Gaussian weight is periodic is bounded by its mean square
over a fundamental domain. The constant depends on the dimension and on a lower bound for the lattice
vectors, through the volume of a ball whose translates do not overlap. The argument is the mean value
inequality for holomorphic functions on that ball.

% before: is-growth-translate-bound
For a theta function for $\Lambda_{L,L'}$, Lemma~\ref{is-growth-sup-mean-square} applies after a
rescaling, with the root discriminant as the lower bound for the lattice vectors. A translate only
multiplies the weighted modulus by a Gaussian factor, which gives the bound for all translates at
once.

% before: is-growth-theta-genus-one
The Galois behaviour of the theta vectors reduces, coordinate by coordinate, to that of the genus-one
theta functions of the CM lattice. An automorphism of $\mathbf C$ that fixes $K$ carries the elliptic
curve of the lattice to that of the conjugate lattice, with its points of finite order and its complex
multiplication. Shimura's reciprocity law then relates the Taylor coefficients of the translates of the
normalized sigma function by corresponding points.

% before: is-growth-theta-algebraic
Along an orthonormal basis a tensor power is a product of genus-one functions of the coordinates, and
its translate by a torsion point is the product of the genus-one translates. So the algebraicity of the
genus-one Taylor coefficients passes to tensor powers, to their translates and to their averages. This
makes the Taylor coefficients of a theta vector up to units algebraic, so that its conjugates are
defined.

% before: is-growth-theta-galois
For an embedding that fixes $K$, the correspondence of torsion points of the elliptic curve extends
coordinatewise along an orthonormal basis. It carries the subgroups $\Lambda_{L,L'}$ to their conjugates,
and conjugation of Taylor coefficients turns a translate by a torsion point into the translate by the
corresponding point, with the phases controlled by the Weil pairing.

% before: is-growth-theta-odd-norm
The norm of a conjugate theta vector is a power of the genus-one norm $n(\Lambda,t)$, and the
logarithm of this power enters the growth bound. That needs $n(\Lambda,t)\ne0$. On a lattice with
$\Delta=1$ the genus-one norm is the norm of $\sigma^*_\Lambda$, a fixed positive number, divided by the
value of the translate at $0$.

% before: is-growth-padded-basis
Alpöge conjugates the translates of a theta function for $\Lambda_{L,L'}$ on a space with more columns, where
the theta function becomes the restriction of a combination of translates of a tensor power of
$\sigma^*_\Lambda$. Such tensor powers need an orthonormal basis of the larger space of matrices. With four
times as many columns, one comes from an orthogonal basis and the four-square theorem.

% before: is-growth-padded-average
A nonzero entire function that is periodic for a multiple of a self-dual subgroup averages to a theta
function for the subgroup, and the average is a combination of translates of the function by torsion
points. The averaging uses that $\alpha_{\sqrt d/2}$ refines the pairing $e^{-\pi i\operatorname{Im}H}$ on the
subgroup. Translates of a theta function by torsion points also do not change when the translation vector
moves by twice a suitable multiple of a vector of the subgroup.

% before: is-growth-padded-restriction
Appending zero columns embeds $V_{a,b}$ isometrically in $V_{a',b}$, and a multiple of $\Lambda_{L,L'}$ lands
in the lattice of an orthonormal basis of the larger space. A translate of the tensor power of
$\sigma^*_\Lambda$ by a half-period, restricted to $V_{a,b}$, is then periodic for a multiple of
$\Lambda_{L,L'}$, and averaging it gives a theta function for $\Lambda_{L,L'}$. As these form a line, every
theta function for $\Lambda_{L,L'}$ is such a restriction.

% before: is-growth-padded-galois
On the larger space the coefficients of the combination are arbitrary complex numbers, so the conjugation
is by an automorphism of $\mathbf C$. The correspondence of translation vectors has to pass from the
matrices that generate $\Lambda_{L,L'}$ to the orthonormal basis, at the cost of a deeper level, and
conjugation has to commute with restriction to $V_{a,b}$.

% before: is-growth-theta-conjugate
Averaging over corresponding representatives turns the conjugate of a theta vector up to units into a
theta vector up to units for the conjugate lattice and the conjugate types. Its norm is then a power of
the genus-one norm, after the types are moved back by an automorphism of the host field. Embeddings that
move $K$ reduce to those that fix it by complex conjugation.

% before: is-growth-theta-conjugate-direct
The proof of Lemma~\ref{is-growth-theta-conjugate} follows Alpöge's route, which reaches the translates
of a theta vector through a larger space and the positivity through the norm of $\sigma^*_\Lambda$. For
a theta vector up to units the orthonormal basis that defines it already serves, and the positivity is
visible from the value at $0$.

% before: is-growth-coset-split
A coset theta function is built from the vacuum vector of a pair of levels by averaging over a pair of
self-dual lattices that split the period lattice along the places of the level. The operators below
carry out this averaging for any Hermitian space, so that the identities they satisfy can be proved
once. They send a function on one of the two lattices to a finite combination of translates of a
theta function.

% before: is-growth-coset-identities
The semicharacter law lets translates of a theta function pass through one another up to roots of
unity. With it, the image of an arbitrary function expands along the translates of the image of an
indicator, and the image of an indicator for one split refines into images for a finer split.

% before: is-growth-coset-lift
For a pair of levels the split lattices come from the lattice of matrices between the two levels,
enlarged at the places of the level and shrunk at their conjugates, or the reverse. The coset theta
function of a class is a translate, by a split representative of the class, of the image of the
indicator of one of these lattices.

% before: is-growth-coset-lattices
The lattices of matrices that enter the construction are compared one place at a time. At the places of
the level, membership in the scaled lattices becomes a congruence modulo a multiple of $H_S(N,N')$, and
the Chinese remainder theorem separates the places of the level from their conjugates.

% before: is-growth-coset-hermitian
The normalization of the translation operators is a semicharacter on a lattice of matrices because of
the arithmetic of the Hermitian form there. Its values have integral imaginary part, and their real
part is an integer multiple of $1/\sqrt d$ with the same parity as the imaginary part. This comes from
the ring of integers of $K$, whose elements have half-integral coordinates of equal parity, together
with the integrality of the trace pairing of the matrices.

% before: is-growth-coset-triples
The identities for the averaging operators need split triples and a vector fixed by the period lattice.
Both follow from the self-duality of the lattices of matrices between self-dual lattices, applied to the
second level scaled at the places of the level and at their conjugates.

% before: is-growth-coset-indicator
With split triples at every modulus, a coset theta function can be compared with the Heisenberg lift of
the indicator of its coset at any modulus that its own modulus divides. The comparison runs term by term
over representatives of the finer split lattice, and it shows that the coset theta functions depend
linearly on the indicators of the cosets.

% before: is-growth-theta-facts
The conjugates of the theta vectors up to units under the embeddings of $\overline{\mathbf Q}$ are
governed by the complex multiplication of the lattice. A conjugate of a translate by a torsion point
is a translate of the conjugate, for a conjugate lattice that is again a CM lattice of the host, and
the norms of the conjugates do not depend on the pair. The kernel data supply the vacuum vectors of
the pairs of levels and the refinement of a coset theta function along a smaller modulus.

% before: is-growth-conjugate-translate
With the Galois action, the bound for translates becomes a growth bound for every conjugate of a
torsion translate of a theta vector up to units. The conjugation may permute the types, which
preserves the Hermitian form and the norm, and embeddings that move $K$ are reduced to those that
fix it by complex conjugation.

% before: is-growth-coset-translates
A coset theta function is induced from the vacuum vector through representatives of a finite
quotient of lattices, and the commutation of the translation operators writes it as a sum of torsion
translates of the vacuum vector. The coefficients are averages of roots of unity, because the
Hermitian form takes rational values on torsion points, and the number of terms is bounded by an
index of lattices.

% before: is-growth-coset-refinement
A kernel entry sums coset theta functions modulo $c_B$, while its coefficient function is invariant
under the coarser modulus $T_ug_0g_1$. Grouping the cosets along the coarser modulus leaves a number
of coset theta functions that does not depend on the finer one.

% before: is-growth-coset-places
The refinement of coset theta functions needs a set of places with no complex-conjugate pair, at which
the ideal $\mathfrak c\mathfrak b^{-1}$ is a unit. For the places of the level both follow from the shape
of the level primes: each factors in $\mathcal O_E$ as an element times its conjugate, the two factors
are coprime, and the prime does not divide the conductor of the host. The ideal $\mathfrak b$ divides the
different of $E$ over $K$, which involves only places over primes dividing that conductor.

% before: is-kernel-entry-bound
A single kernel entry is one jet operator of a combination of theta functions. Its norm at a
conjugate embedding is bounded by the block bound, with the periods and the vacuum norm of the
conjugated CM lattice.

% before: is-transference
A self-dual lattice is counted through a basis of short vectors, and short vectors in a lattice
follow from a lower bound for the vectors of its dual. The transference theorem turns the one bound
into the other, with a loss polynomial in the rank.

% before: is-selfdual-classes
The class counts below rest on a bound for the number of classes of self-dual Hermitian lattices.
It is polynomial in the root discriminant of the host, with an exponent of order $fn^2$, so its
logarithm is absorbed into the error term $U$.

% before: is-class-counts
The step bounds add the counts of lattice classes and of cosets to the bound for a single entry. Each
count is at most the exponential of a constant times $U=fm^2N_m^2\log(g_0g_1T_u)$, which is an error
term.

% before: is-vacuum-averaging
The averaging operators move a theta function from one lattice to another. For the lattices that
occur here they preserve norms and compose up to an eighth root of unity, so a theta vector up to
units stays one when it is moved.

% before: is-vacuum-tensor-norm
An orthonormal basis gives coordinates on $V_{a,b}$ in which its lattice is a product of copies of
$\Lambda$. In these coordinates the norm of a tensor power is a power of the norm in genus one.

% before: is-vacuum-genus-one
The norm in genus one is computed from a theta constant of the lattice. It does not depend on the
choice of the unkilled two-torsion point, and complex conjugation of the lattice does not change it.

% before: is-vacuum-unkilled-splitting
The proof of Lemma~\ref{is-vacuum-genus-one}(1) finds the unkilled point from an ideal of odd norm
and separates the four classes of two-torsion points by hand. The decomposition of $2$ in $K$ gives
both facts at once, because it identifies the two-torsion points with a product of two fields with
two elements.

% before: is-vacuum-eta-weber
When the discriminant of a lattice is $1$, its covolume is given by the eta function. The theta
constant of the norm in genus one, divided by $\eta$ and raised to the twelfth power, is up to sign
the inverse cube of an algebraic number built from the two-torsion values of $\wp_\Lambda$, and
Weber's functions show that it is an algebraic integer.

% before: is-vacuum-eta-mean
The two formulas of Lemma~\ref{is-vacuum-eta-weber} bound the mean of the norms in genus one over the
conjugates of the CM lattice. The theta constants contribute at least as much as the eta values,
because a nonzero algebraic integer has norm at least $1$, and the mean of the eta values is given by
the Chowla-Selberg formula.

% before: is-vacuum-units-intertwiners
The vacuum vector of the host is given on four coordinates at a time, by identity (4) of
Definition~\ref{pd-cm-datum}, while a theta vector up to units is built on all of $V_{a,b}$ from an
orthonormal basis. Comparing the two moves averaging operators along a linear isometry and across a
decomposition into orthogonal blocks. On lattices stable under $\omega_K$ the normalization
$\alpha_{\sqrt d/2}$ is a sign that is multiplicative up to $e^{\pi i\operatorname{Im}H}$, so the sums that
define these operators do not depend on the representatives.

% before: is-vacuum-units-host
Every vacuum vector of a pair is obtained by averaging the tensor power $\theta_0^{\otimes ab}$. For
even ranks the matrices split into $2\times2$ blocks, each a copy of $E^4$, and the isometry
$\gamma_0$ of the kernel data turns the standard basis of each block into an orthonormal basis of the
whole space. Identity (4) on each block then writes $\theta_0^{\otimes ab}$ as an averaged tensor power
of the genus-one theta function.

% before: is-vacuum-norm-bound
The vacuum norms enter the sizes of the lifts through their mean over the conjugate embeddings. The
Chowla-Selberg formula bounds that mean by the logarithmic root discriminant and the logarithmic
derivative of $L(s,\chi_d)$ at $1$.

% before: is-vacuum-determinant-dual
The proof of Lemma~\ref{is-vacuum-norm-bound} shows $\mathfrak d_L\bar{\mathfrak d}_L=\mathcal O_E$
for a self-dual lattice $L$ one prime at a time, from local bases. Duality gives the same identity
globally, since passing to the dual lattice inverts the determinant ideal and conjugates it.

% before: is-kernel-step-bound
Each lift multiplies the inner products by at most the exponential of its logarithmic size. The
main term $\phi(\nu)-\nu\log\Omega$ comes from the Fock norm of a theta kernel, and its average over
the conjugates is computed by the Chowla-Selberg formula. The counts of lattice classes only add
error terms.

% before: is-top-period-bound
The last step of the chain is the period at the top, which restricts a form on $U(X_m)$ to the
subgroup of corank one. Its norm is the norm of a branching map, which is explicit for interlacing
weights. At a place of $T$ the factor $T_u^{-n_\tau}$ of this norm gives the term $-a_t\log T_u$ of
$\mathrm{MAIN}(u)$, up to the error terms.

% before: is-definite-side-condition-at-bundle
With the two halves and the chain lemma, the definite-side condition holds at a single bundle as
soon as its size constant exceeds a threshold depending on $C_Y$ and $C_T$. The remaining clauses
of the rows condition come from the division lattices: they must be lattices of positive rank in
the column.

% before: is-definite-side-condition-holds
The bundles built from the linked parameters satisfy the evaluation conditions, and along designed
families the size constant grows beyond any threshold fixed by the other constants. So the
definite-side condition holds for them, which is hypothesis (9) of Theorem~\ref{is-definite-rows}.

% before: is-local-indices
The division inequality compares covolumes of lattices that agree at almost all places. Local
indices, local bases and Gram determinants are the quantities in which these comparisons are made,
place by place.

% before: is-lattice-facts
Covolumes are computed through a basis and the local indices of the lattice relative to it. These
rules, with the product formula, are used at every step of the proof of the division inequality.

% before: is-gram-covolume
The covolume of the base lattice of the rows is bounded through the Gram determinants of its local
bases. The quotient metric is not defined over the subfield $K_0$, so the bound is an inequality,
and it comes from a Cauchy-Schwarz inequality for determinants at pairs of conjugate places.

% before: is-local-index-identity
At a twisted place the glued lattice agrees with the image of the lattice of a twisted tuple, and
its index relative to the image at the base tuple has to be computed. The determinant identities at
the place $v_0$ relate this index to Gram determinants and to the indices of the image lattices of
the lifts.

% before: is-division-package
The division inequality is proved for linear-algebra data that collect a realization at one
member: the rows, which carry the covolume and surjectivity bounds, and the columns, which carry the
lower norm bound and the column maps $B_u$. The passage compares the dual glued lattices of the two.

% before: is-glued-covolume
The slope of the dual of the glued lattice of the rows is computed by gluing the base lattice with
the twisted images. The covolume bound controls the base, and the surjectivity bound controls the
indices at the twisted places.

% before: is-column-height
The height of a column map is bounded by the norms of the steps of its chain of lifts and of the
period at the top. Inverting the bijection $\iota_{\mu_0}$ costs the smallest of the scalars by
which it multiplies metrics, and the lower norm bound controls these scalars.

% before: is-shift-into-dual
The multiplicity bound says that the periods of the division vectors are divisible at each place.
Dividing the division lattice by these valuations keeps its image in the dual glued lattice, and
raises its slope by the logarithmic size of the shift.

% before: is-record-numbers
The division inequality ends in a numerical comparison between quantities of the size data and the
main term of a word. The record numbers collect these quantities, with the inequalities between
them that the size data guarantee.

% before: is-record-height
The height of a column map is bounded by the period exponent and the mean sizes of the lifts, less
the lower norm bound. This lemma compares that bound with the height $h^{\mathrm{rec}}$ of the record,
the kernel side and the last side separately.

% before: is-record-division
Once the slope of the dual glued lattice is bounded, the rest of the division inequality is
arithmetic. The main terms of the covolume bound and of the height cancel up to the term
$\mathrm{MAIN}(u)$, and the remaining terms fit in the error budget.

% before: is-division-inputs
Besides the package, the division inequality for a word uses its chain of lifts with height
numbers, the data that bound the slope of the rows, and the valuations of its periods. These are
the inputs at one word.

% before: is-division-inequality
The division inequality bounds the slope of the auxiliary lattice $CC$ by the main term less the
height of the word. It combines the slope bound of the rows, the passage to the columns, the height
bound of the column map and the arithmetic of the record.

% before: is-division-link
The division link is the form in which the upper half of the comparison inequality of
Chapter~\ref{chapter-inert-prime-balance} is used. Its central clause bounds
$\mu(CC)+h_u$ by the main term $\mathrm{MAIN}(u)$, which collects the cost $\phi(a_t)-a_t\log\Omega$
of each of the $m$ lifts and the size of $T_u$. The other clauses compare words with one another and
record the identity of the periods with Bessel periods.

% before: is-division-kernel-bounds
The division link is reduced to inputs that can be checked at each member separately. Apart from the
division packages, these inputs are the multiplicity bound of
Chapter~\ref{chapter-period-divisibility}, comparisons of the explicit heights, and the
identification of the Bessel pairings.

% before: is-division-link-holds
Once the inputs hold at every member, the division link follows member by member. The passage
from the inputs to the link is uniform along a class of families, so the thresholds of the class
serve for both.

% before: is-orthogonality-link
The slope inequality of Chapter~\ref{chapter-hermitian-slopes} needs the columns of the glued
integral structure to form an almost tight frame. The orthogonality link expresses this through the
normalized Gram matrices of the columns: off the diagonal they are no larger than for random
subspaces of the expected dimension, and on the diagonal they are close to the identity.

% before: is-wu-decay
The columns are Poincaré series built from archimedean Bessel functions $W_u$, one for each word. For
the series to converge and to be square integrable, $W_u$ must decay fast enough against a power of
the norm on the stabilizer quotient. The gap hypothesis on the Harish-Chandra parameters is the
condition under which this decay is proved.

% before: is-gram-terms
The Gram matrices of the columns are evaluated by unfolding the Poincaré series into a sum over
double cosets. The terms with vanishing invariant that are not the identity form the leading part,
which has to be treated separately from the rest.

% before: is-unfolding-term-data
The four clauses use only a few features of the unfolded sum: how its terms split by their
invariant, and how the norms of partial sums are controlled by vectors attached to the terms and a
matrix indexed by their types. Term data keep these features and nothing else, so the norm bounds can
be proved once for an abstract sum of operators.

% before: is-unfolding-identities
A partial sum over a set of terms is bounded by the squared norms of its vectors grouped by type.
Row orthogonality turns these squared norms into sums over pairs of terms in the same row, which the
corank then sorts. The bound for the degenerate terms and the bound on the diagonal both use them.

% before: is-unfolding-pair-bound
After the identity part is removed, the rest of the sum consists of the generic and the non-generic
terms. Bounds on the four kinds of pairs therefore bound its norm, with a constant that depends only
on the bound for the matrix indexed by the types.

% before: is-unfolding-evaluation
The term data of the Gram matrices come from an automorphic setting that realizes the columns as
Poincaré series. Its terms are double cosets of rational isometries between rows, and their operators
are orbital integrals of the archimedean Bessel functions weighted by a character of the level
subgroup.

% before: is-unfolding-computations
The structural hypotheses of the norm bounds have to be checked for the unfolded series. Each is a
computation with the orbital integrals, the level characters or the character sums over the general
linear groups of the residue fields.

% before: is-unfolding-transport
The clauses concern the Gram matrices of the datum, while the bounds concern the term data of the
setting. The isomorphism of Bessel-period data and the identification of columns with Poincaré series
are isometries, so they carry the norms in the clauses to norms of partial sums of term data.

% before: is-unfolding-schur
For distinct words the degenerate terms must contribute nothing. Their orbital integrals are covariant
under a group that contains a compact unitary group at one archimedean place, on which the
representations of the two words have different types.

% before: is-unfolding-degenerate-pairs
For a single word the degenerate terms do not vanish, and their size has to be bounded. The
positivity of the row types bounds the products of the norms of their vectors, pair by pair.

% before: is-unfolding-degenerate
The pair bound controls the leading part on the diagonal through the norm bound for partial sums.
The resulting bound is exponentially small in $kf\log k$, which is the form the combination of the
clauses needs.

% before: is-unfolding-diagonal-numerics
The diagonal bound produced by the pair bound is a square of a sum of square roots of the constants of
the row bounds. It has to be compared with $\Gamma$, which holds once $C_\Gamma$ exceeds the
constants of the main-term identity and of the non-generic row bound.

% before: is-unfolding-diagonal
On the diagonal the identity part of the unfolded sum is the identity operator. What remains is
bounded by the row bounds of the main-term identity, the non-generic terms and the loops, and the
numerical comparison turns this into the factor $\Gamma$.

% before: is-gram-clauses
Near-orthogonality reduces to four clauses about the leading parts and the remainders. Each clause
comes from a different input of Theorem~\ref{is-orthogonality-link-holds}.

% before: is-gram-combination
The four clauses and the point counts give the orthogonality bounds by the triangle inequality for
Hilbert-Schmidt norms. The sizes of the family then make the three constants exponentially small in
$x^{\varepsilon/8}$.

% before: is-orthogonality-link-holds
The Gram matrices are evaluated by unfolding the Poincaré series, which turns them into sums over
points with local Bessel integrals. Bounds for these sums, the main terms of the diagonal, which come
from the Siegel-Weil formula, and point counts give near-orthogonality.

% before: is-irreducibility-link
The frame inequality loses a factor $(1-\delta)^{-1}$, where $\delta$ measures how far the row space
falls short of its expected dimension. The irreducibility link asks for a deficit smaller than any
negative power of $x$.

% before: is-row-labelling
At the two level primes $g_0,g_1$ every row carries a label, a polynomial that records the cuspidal
support of its local component. Rows with bad labels are the ones that may be lost, so the deficit is
bounded by counting labels.

% before: is-irreducibility-link-holds
Green's description of the cuspidal characters of $GL_r(\mathbf F_g)$ \cite{Gre55} gives the
proportion of each label among the rows. Bad labels are rare once $g$ is larger than
$\exp(8x^{\varepsilon_z})$, and the row count turns this proportion into a bound for the deficit.
